% Fonctions : généralités, limites et continuité — Mathématiques 1ère A — Cours signé OnBuch+
% Programme officiel MINESEC. Compiler avec : tectonic NA02-fonctions-limites-continuite-a.tex
\newcommand{\DOCMATIERE}{Mathématiques}
\newcommand{\DOCNIVEAU}{1ère A}
\newcommand{\DOCMODULE}{Relations et opérations fondamentales dans l'ensemble des nombres réels}
\newcommand{\DOCLECON}{NA02}
\newcommand{\DOCTITRE}{Fonctions : généralités, limites et continuité}
% OnBuch+ — préambule « cours signés » ENRICHI (compilé avec tectonic / XeLaTeX)
% Système visuel « Pop » d'OnBuch+ : fond crème (jamais blanc), bordures encre
% épaisses, ombres « dures » décalées, titres Archivo Black en capitales.
% Version MATHÉMATIQUES : graphes (pgfplots/tikz), boîtes pédagogiques
% (définition, propriété, méthode, attention, le savais-tu, exercice résolu,
% exercice, TP, bilan), page de garde et signature OnBuch+ ancrée en bas.
%
% Macros attendues dans main.tex AVANT \input{preamble} :
%   \DOCMATIERE  \DOCNIVEAU  \DOCMODULE  \DOCLECON (numéro)  \DOCTITRE
\documentclass[11pt]{article}

\usepackage{fontspec}
\usepackage[french]{babel}
\usepackage[a4paper,margin=2cm,top=3.4cm,bottom=2.8cm,headheight=1.4cm,footskip=1.3cm]{geometry}
\usepackage{amsmath,amssymb,amsfonts}
\usepackage{mathtools} % \xmapsto, \coloneqq... souvent utilisés par les modèles
\usepackage{cancel} % simplifications barrées
% \abs{z} (module, valeur absolue) et \norm{u} (norme), utilisés par les modèles
\providecommand{\abs}{}\let\abs\relax\DeclarePairedDelimiter{\abs}{\lvert}{\rvert}
\providecommand{\norm}{}\let\norm\relax\DeclarePairedDelimiter{\norm}{\lVert}{\rVert}
\usepackage[table]{xcolor} % [table] : fournit aussi \rowcolors (alternance de lignes)
\usepackage{pagecolor}
\usepackage{tikz}
\usetikzlibrary{arrows.meta,positioning,calc,shapes.geometric,shapes.misc,shapes.arrows,decorations.pathmorphing,decorations.pathreplacing,decorations.markings,patterns,shadows,mindmap,trees,fit,backgrounds,matrix,intersections,angles,quotes}
\usepackage{pgfplots}
\usepackage{pgf-pie} % diagrammes circulaires \pie (statistiques)
\pgfplotsset{compat=1.18}
\usepgfplotslibrary{fillbetween,groupplots} % aires entre courbes (calcul intégral)
\usepackage{enumitem}
\usepackage{fancyhdr}
% [most] charge toutes les bibliothèques tcolorbox (listings, minted, raster,
% documentation...) et gonfle inutilement la mémoire TeX sur un document long
% (~300 boîtes) : on ne charge que ce qui est réellement utilisé.
\usepackage[skins,breakable]{tcolorbox}
\usepackage{graphicx}
\usepackage{colortbl}
\usepackage{booktabs}
\usepackage{tabularx}
\usepackage{diagbox} % case d'en-tête diagonale des tableaux à double entrée
% Colonnes extensibles centrée / gauche / droite (C, L, R), souvent utilisées par les modèles
\newcolumntype{C}{>{\centering\arraybackslash}X}
\newcolumntype{L}{>{\raggedright\arraybackslash}X}
\newcolumntype{R}{>{\raggedleft\arraybackslash}X}
\usepackage{array}
\usepackage{multicol}
\usepackage{siunitx}
% parse-numbers=false : accepte aussi \SI{2,5 \times 10^{-3}}{...}
\sisetup{locale=FR,output-decimal-marker={,},group-minimum-digits=4,parse-numbers=false}
\DeclareSIUnit\ohm{\ensuremath{\symup{\Omega}}} % U+2126 absent de la police
\usepackage{qrcode}
% \degree : commande fréquemment utilisée par les modèles pour un angle ou une
% température en dehors de \SI{}{} (non fournie par les paquets chargés).
\providecommand{\degree}{^{\circ}}
% \qed (fin de démonstration), utilisé par les modèles sans amsthm
\providecommand{\qed}{\hfill$\square$}
% \barre : double barre des valeurs interdites dans un tableau de variations (tkz-tab)
\providecommand{\barre}{\ensuremath{\|}}
% environnement proof (amsthm non chargé) utilisé par les modèles
\ifdefined\proof\else\newenvironment{proof}[1][Démonstration]{\par\noindent\textit{#1.}\ }{\qed\par}\fi
\usepackage{lastpage}
\usepackage{hyperref}
\hypersetup{hidelinks}

% ============================================================
% Palette « Pop » OnBuch+ — mêmes hex que la classe P (plus_ui.dart)
% ============================================================
\definecolor{poporange}{HTML}{F59321}
\definecolor{popdark}{HTML}{E07A0C}
\definecolor{popcream}{HTML}{FFF6E8}
\definecolor{popink}{HTML}{1C1714}
\definecolor{poppurple}{HTML}{7A5AE0}
\definecolor{popgreen}{HTML}{2E9E6A}
\definecolor{poppink}{HTML}{E0457B}
\definecolor{popblue}{HTML}{2F7DE1}
\definecolor{popgold}{HTML}{C98A12}
\definecolor{popmuted}{HTML}{8A7F76}
\definecolor{popline}{HTML}{EADFD2}
\definecolor{popgray}{HTML}{8A7F76}
\colorlet{popgrey}{popgray}
% Alias tolérants (le modèle nomme parfois les couleurs différemment)
\colorlet{popwhite}{white}
\colorlet{popblack}{popink}
\colorlet{popred}{poppink}
\colorlet{popyellow}{popgold}
% Teintes claires pour fonds de tableaux / zones
\colorlet{popblueL}{popblue!10!white}
\colorlet{poporangeL}{poporange!14!white}
\colorlet{popgreenL}{popgreen!12!white}
\colorlet{poppurpleL}{poppurple!12!white}
\colorlet{poppinkL}{poppink!12!white}
\colorlet{popgoldL}{popgold!14!white}

\pagecolor{popcream}

% ============================================================
% Typographie
% ============================================================
\defaultfontfeatures{Ligatures=TeX}
\setmainfont{PlusJakartaSans-Medium.ttf}[Path=../../../fonts/,BoldFont=PlusJakartaSans-Bold.ttf]
% Maths en sans-serif (Fira Math) avec chiffres et lettres droites de Plus
% Jakarta, pour que les formules se fondent dans le texte.
\usepackage[math-style=ISO,bold-style=ISO]{unicode-math}
\setmathfont{FiraMath-Regular.otf}[Path=../../../fonts/]
\setmathfont{PlusJakartaSans-Medium.ttf}[Path=../../../fonts/,range={up/num,up/{Latin,latin},\mathup}]
\usepackage{icomma} % virgule décimale sans espace en mode math
% Caractères mathématiques Unicode absents des polices de texte (Plus Jakarta,
% Archivo Black) : rendus par leur équivalent mathématique, en texte comme en
% formule — sinon ils s'affichent en carré vide (ex. « Arithmétique dans ℤ »).
\usepackage{newunicodechar}
\newunicodechar{ℝ}{\ensuremath{\mathbb{R}}}
\newunicodechar{ℤ}{\ensuremath{\mathbb{Z}}}
\newunicodechar{ℂ}{\ensuremath{\mathbb{C}}}
\newunicodechar{ℕ}{\ensuremath{\mathbb{N}}}
\newunicodechar{ℚ}{\ensuremath{\mathbb{Q}}}
\newunicodechar{𝔻}{\ensuremath{\mathbb{D}}}
\newunicodechar{α}{\ensuremath{\alpha}}
\newunicodechar{β}{\ensuremath{\beta}}
\newunicodechar{γ}{\ensuremath{\gamma}}
\newunicodechar{δ}{\ensuremath{\delta}}
\newunicodechar{ε}{\ensuremath{\varepsilon}}
\newunicodechar{θ}{\ensuremath{\theta}}
\newunicodechar{λ}{\ensuremath{\lambda}}
\newunicodechar{μ}{\ensuremath{\mu}}
\newunicodechar{σ}{\ensuremath{\sigma}}
\newunicodechar{τ}{\ensuremath{\tau}}
\newunicodechar{φ}{\ensuremath{\varphi}}
\newunicodechar{ψ}{\ensuremath{\psi}}
\newunicodechar{ω}{\ensuremath{\omega}}
\newunicodechar{Σ}{\ensuremath{\Sigma}}
% majuscules produites par \MakeUppercase dans les titres (α-aminés -> Α)
\newunicodechar{Α}{\ensuremath{\alpha}}
\newunicodechar{Β}{\ensuremath{\beta}}
\newunicodechar{Γ}{\ensuremath{\Gamma}}
\newunicodechar{Δ}{\ensuremath{\Delta}}
\newunicodechar{Ε}{\ensuremath{\varepsilon}}
\newunicodechar{Θ}{\ensuremath{\Theta}}
\newunicodechar{Λ}{\ensuremath{\Lambda}}
\newunicodechar{Μ}{\ensuremath{\mu}}
\newunicodechar{Τ}{\ensuremath{\tau}}
\newunicodechar{Φ}{\ensuremath{\Phi}}
\newunicodechar{Ψ}{\ensuremath{\Psi}}
\newunicodechar{Ω}{\ensuremath{\Omega}}
\newunicodechar{↦}{\ensuremath{\mapsto}}
\newunicodechar{∈}{\ensuremath{\in}}
\newunicodechar{∉}{\ensuremath{\notin}}
\newunicodechar{∧}{\ensuremath{\wedge}}
\newunicodechar{⇔}{\ensuremath{\Leftrightarrow}}
\newunicodechar{⟺}{\ensuremath{\Longleftrightarrow}}
\newunicodechar{⇒}{\ensuremath{\Rightarrow}}
\newunicodechar{⟹}{\ensuremath{\Longrightarrow}}
\newunicodechar{∘}{\ensuremath{\circ}}
\newunicodechar{∪}{\ensuremath{\cup}}
\newunicodechar{∩}{\ensuremath{\cap}}
\newunicodechar{≡}{\ensuremath{\equiv}}
\newunicodechar{⊂}{\ensuremath{\subset}}
\newunicodechar{⊥}{\ensuremath{\perp}}
\newunicodechar{∃}{\ensuremath{\exists}}
\newunicodechar{∀}{\ensuremath{\forall}}
\newunicodechar{∛}{\ensuremath{\sqrt[3]{\,}}}
\newunicodechar{∜}{\ensuremath{\sqrt[4]{\,}}}
\newunicodechar{⃗}{\ensuremath{{}^{\to}}}
\newunicodechar{ⁿ}{\ifmmode^{n}\else\textsuperscript{\ensuremath{n}}\fi}
\newunicodechar{ˣ}{\ifmmode^{x}\else\textsuperscript{\ensuremath{x}}\fi}
\newunicodechar{ᵇ}{\ifmmode^{b}\else\textsuperscript{\ensuremath{b}}\fi}
\newunicodechar{ʳ}{\ifmmode^{r}\else\textsuperscript{\ensuremath{r}}\fi}
\newunicodechar{ᵘ}{\ifmmode^{u}\else\textsuperscript{\ensuremath{u}}\fi}
\newunicodechar{ʸ}{\ifmmode^{y}\else\textsuperscript{\ensuremath{y}}\fi}
\newunicodechar{ᵃ}{\ifmmode^{a}\else\textsuperscript{\ensuremath{a}}\fi}
\newunicodechar{ᵉ}{\ifmmode^{e}\else\textsuperscript{\ensuremath{e}}\fi}
\newunicodechar{ᵏ}{\ifmmode^{k}\else\textsuperscript{\ensuremath{k}}\fi}
\newunicodechar{ᵖ}{\ifmmode^{p}\else\textsuperscript{\ensuremath{p}}\fi}
\newunicodechar{ᴺ}{\ifmmode^{N}\else\textsuperscript{\ensuremath{N}}\fi}
\newunicodechar{ᶜ}{\ifmmode^{c}\else\textsuperscript{\ensuremath{c}}\fi}
\newunicodechar{ʰ}{\ifmmode^{h}\else\textsuperscript{\ensuremath{h}}\fi}
\newunicodechar{ᵠ}{\ifmmode^{q}\else\textsuperscript{\ensuremath{q}}\fi}
\newunicodechar{ₙ}{\ifmmode_{n}\else\textsubscript{\ensuremath{n}}\fi}
\newunicodechar{ₐ}{\ifmmode_{a}\else\textsubscript{\ensuremath{a}}\fi}
\newunicodechar{ᵢ}{\ifmmode_{i}\else\textsubscript{\ensuremath{i}}\fi}
\newunicodechar{ᵣ}{\ifmmode_{r}\else\textsubscript{\ensuremath{r}}\fi}
\newunicodechar{ₖ}{\ifmmode_{k}\else\textsubscript{\ensuremath{k}}\fi}
\newunicodechar{ᵦ}{\ifmmode_{\beta}\else\textsubscript{\ensuremath{\beta}}\fi}
\newunicodechar{ₓ}{\ifmmode_{x}\else\textsubscript{\ensuremath{x}}\fi}
\newfontfamily\popdisplay{ArchivoBlack-Regular.ttf}[Path=../../../fonts/]
\newfontfamily\poplogo{SpaceGrotesk-Bold.ttf}[Path=../../../fonts/]
\newfontfamily\popsemi{PlusJakartaSans-SemiBold.ttf}[Path=../../../fonts/]
\newfontfamily\popxbold{PlusJakartaSans-ExtraBold.ttf}[Path=../../../fonts/]
\newfontfamily\popxboldls{PlusJakartaSans-ExtraBold.ttf}[Path=../../../fonts/,LetterSpace=3.0]

\setlist[enumerate]{leftmargin=1.7em,itemsep=5pt,topsep=4pt}
\setlist[itemize]{leftmargin=1.7em,itemsep=4pt,topsep=4pt,label={\color{poporange}\textbullet}}
\setlength{\parindent}{0pt}
\setlength{\parskip}{5pt}
\renewcommand{\arraystretch}{1.25}

% pgfplots : style Pop par défaut
\pgfplotsset{
  every axis/.append style={axis line style={popink,line width=1pt},tick style={popink},
    grid style={popline,line width=0.5pt},label style={font=\small},tick label style={font=\footnotesize}},
  popcurve/.style={poporange,line width=1.8pt,no markers,smooth},
}
% Même style disponible dans un \draw TikZ simple (hors pgfplots)
\tikzset{popcurve/.style={poporange,line width=1.8pt}}

% ============================================================
% Pill / squiggle
% ============================================================
\newcommand{\poppill}[3]{%
  \tikz[baseline=-0.55ex]\node[draw=popink,line width=1.2pt,rounded corners=2.6pt,%
    fill=#2,inner xsep=6.5pt,inner ysep=3pt]{%
    \popxboldls\fontsize{7}{7}\selectfont\color{#3}\MakeUppercase{#1}};%
}
\newcommand{\popsquiggle}[1]{%
  \tikz[baseline=-0.4ex]{
    \draw[#1,line width=2.6pt,line cap=round]
      plot [smooth] coordinates {(0,0.09) (0.34,0.01) (0.68,0.21) (1.02,0.01) (1.36,0.10)};
  }%
}
% Mot-clé surligné (marqueur orange sous le texte)
\newcommand{\cle}[1]{{\popxbold\color{popdark}#1}}

% ============================================================
% En-tête / pied
% ============================================================
\pagestyle{fancy}
\fancyhf{}
\renewcommand{\headrulewidth}{0pt}
\renewcommand{\footrulewidth}{0pt}
\lhead{\raisebox{-0.15ex}{\poplogo\fontsize{13}{13}\selectfont\color{popink}OnBuch\color{poporange}+}}
\rhead{\poppill{\DOCMATIERE\ \textbullet\ \DOCNIVEAU\ \textbullet\ Leçon \DOCLECON}{white}{popink}}
\lfoot{\popxbold\fontsize{7.6}{7.6}\selectfont\color{popmuted}\MakeUppercase{Cours signé OnBuch+}\ \textbullet\ {\popsemi\DOCTITRE}}
\rfoot{\tikz[baseline=-0.6ex]\node[rounded corners=8.5pt,fill=popink,minimum height=17pt,minimum width=17pt,inner xsep=5pt,inner sep=0]{\popxbold\fontsize{8}{8}\selectfont\color{popcream}\thepage\,/\,\pageref*{LastPage}};}

% ============================================================
% Titres
% ============================================================
\newcommand{\chaptitre}[2]{%
  \begin{tcolorbox}[enhanced,colback=poporange,colframe=popink,boxrule=2.6pt,arc=11pt,
      drop shadow=popink,left=15pt,right=15pt,top=12pt,bottom=11pt,before skip=3pt,after skip=18pt]
    \poppill{#2}{popink}{popcream}\par\vspace{9pt}
    {\raggedright\hyphenpenalty=10000\exhyphenpenalty=10000\popdisplay\fontsize{22}{24}\selectfont\color{popcream}\MakeUppercase{#1}\par}
  \end{tcolorbox}
}
% Section numérotée : pastille orange avec le numéro + titre Archivo
\newcounter{popsec}
\newcommand{\coursec}[1]{%
  \stepcounter{popsec}\setcounter{popsub}{0}%
  \par\vspace{16pt}\noindent
  \begin{minipage}{\linewidth}
  \tikz[baseline=-0.7ex]\node[draw=popink,line width=1.6pt,fill=poporange,rounded corners=4pt,minimum size=22pt,inner sep=0]{\popdisplay\fontsize{12}{12}\selectfont\color{popcream}\thepopsec};%
  \hspace{8pt}{\popdisplay\fontsize{15}{17}\selectfont\color{popink}\MakeUppercase{#1}}\par
  \vspace{4pt}\noindent{\color{poporange}\rule{36pt}{3.6pt}}
  \end{minipage}\par\nopagebreak\vspace{8pt}
}
\newcounter{popsub}
\newcommand{\courssub}[1]{%
  \stepcounter{popsub}%
  \par\vspace{9pt}\noindent{\popxbold\fontsize{11.5}{13}\selectfont\color{popink}{\color{poporange}\thepopsec.\thepopsub}\hspace{6pt}#1}\par\nopagebreak\vspace{4pt}
}

% ============================================================
% Boîtes pédagogiques — même recette : fond blanc, bordure épaisse colorée,
% ombre dure encre, badge plein. Titre optionnel [#1].
% ============================================================
\tcbset{popbase/.style={breakable,enhanced,colback=white,boxrule=2.3pt,arc=9pt,
  shadow={3pt}{-3pt}{0pt}{popink},left=14pt,right=14pt,top=11pt,bottom=11pt,before skip=11pt,after skip=15pt,
  fontupper=\popsemi\fontsize{10.6}{14.5}\selectfont\color{popink},
  fontlower=\popsemi\fontsize{10.6}{14.5}\selectfont\color{popink}}}
% Coche dessinée (le glyphe ✓ n'existe pas dans les polices du document)
\newcommand{\popcheck}[1]{\tikz[baseline=-0.5ex]\draw[#1,line width=1.6pt,line cap=round,line join=round](-0.13,0.0)--(-0.04,-0.09)--(0.13,0.1);}
\newcommand{\popboxhead}[3]{\poppill{#1}{#2}{white}\ifx\relax#3\relax\else\hspace{6pt}{\popxbold\fontsize{10.6}{12}\selectfont\color{popink}#3}\fi\par\vspace{7pt}}

\newtcolorbox{aretenir}[1][]{popbase,colframe=popgreen,before upper={\popboxhead{À retenir}{popgreen}{#1}}}
\newtcolorbox{exemplebox}[1][]{popbase,colframe=poporange,before upper={\popboxhead{Exemple}{poporange}{#1}}}
\newtcolorbox{definition}[1][]{popbase,colframe=popblue,before upper={\popboxhead{Définition}{popblue}{#1}}}
\newtcolorbox{propriete}[1][]{popbase,colframe=poppurple,before upper={\popboxhead{Propriété}{poppurple}{#1}}}
\newtcolorbox{methode}[1][]{popbase,colframe=popgold,before upper={\popboxhead{Méthode}{popgold}{#1}}}
\newtcolorbox{attention}[1][]{popbase,colframe=poppink,before upper={\popboxhead{Attention}{poppink}{#1}}}
\newtcolorbox{experience}[1][]{popbase,colframe=popblue,colback=popblueL,before upper={\popboxhead{Expérience}{popblue}{#1}}}
% « Le savais-tu ? » : carte violette pleine (contexte, culture, Cameroun)
\newtcolorbox{savaistu}[1][]{popbase,colback=poppurple,colframe=popink,
  fontupper=\popsemi\fontsize{10.4}{14.2}\selectfont\color{white},
  before upper={\poppill{Le savais-tu\,?}{white}{poppurple}\ifx\relax#1\relax\else\hspace{6pt}{\popxbold\color{white}#1}\fi\par\vspace{7pt}}}
% Exercice résolu : énoncé en haut, \tcblower puis la solution sur fond crème
\newcounter{popexo}\newcounter{popexr}
\newtcolorbox{exoresolu}[1][]{popbase,colframe=poporange,colbacklower=popcream,
  segmentation style={popink,line width=1.2pt,dashed},
  before upper={\stepcounter{popexr}\popboxhead{Exercice résolu \thepopexr}{poporange}{#1}},
  before lower={\poppill{Solution}{popink}{popcream}\par\vspace{6pt}}}
% Exercice (non résolu en place) — la correction est dans la partie Corrigés
\newtcolorbox{exercice}[1][]{popbase,colframe=popink,
  before upper={\stepcounter{popexo}\popboxhead{Exercice \thepopexo}{popink}{#1}}}
% Corrigé
\newtcolorbox{corrige}[1][]{popbase,colframe=popgreen,colback=popgreenL,
  before upper={\popboxhead{Corrigé}{popgreen}{#1}}}
% Objectifs / prérequis / bilan
\newtcolorbox{objectifs}{popbase,colframe=popink,colback=poporangeL,
  before upper={\popboxhead{Objectifs de la leçon}{popink}{}}}
\newtcolorbox{prerequis}{popbase,colframe=popink,colback=popblueL,
  before upper={\popboxhead{Prérequis}{popblue}{}}}
\newtcolorbox{bilan}{popbase,colframe=popink,colback=white,boxrule=2.8pt,
  before upper={\popboxhead{Fiche bilan}{poporange}{L'essentiel à connaître pour l'examen}}}
% Figure encadrée avec légende
\newtcolorbox{popfigure}[1][]{enhanced,breakable=false,colback=white,colframe=popline,boxrule=1.4pt,arc=8pt,
  left=10pt,right=10pt,top=10pt,bottom=8pt,before skip=10pt,after skip=12pt,halign=center,
  lower separated=false,halign lower=center,
  fontlower=\popsemi\fontsize{9}{11.5}\selectfont\color{popmuted},
  #1}
\newcommand{\legende}[1]{\par\vspace{6pt}{\popsemi\fontsize{9}{11.5}\selectfont\color{popmuted}#1}}

% ============================================================
% Signature OnBuch+
% ============================================================
\newcommand{\poplogoinline}[1]{{\poplogo\fontsize{#1}{#1}\selectfont\color{popink}OnBuch\color{poporange}+}}
\newcommand{\signatureonbuch}{%
  \begin{tcolorbox}[enhanced,colback=popink,colframe=popink,boxrule=2.2pt,arc=10pt,
      drop shadow=poporange,left=14pt,right=14pt,top=10pt,bottom=10pt,before skip=0pt,after skip=0pt]
    \tikz[baseline=-0.7ex]\node[circle,fill=popgreen,draw=popcream,line width=1.3pt,minimum size=22pt,inner sep=0]{\popcheck{white}};%
    \hspace{9pt}\begin{minipage}[c]{0.62\linewidth}
      {\popxbold\fontsize{12}{13}\selectfont\color{popcream}Signé\ \textbullet\ {\poplogo OnBuch\color{poporange}+}}\par\vspace{2pt}
      {\raggedright\popsemi\fontsize{8}{10.5}\selectfont\color{popline}\MakeUppercase{Équipe pédagogique OnBuch+}\\\MakeUppercase{Conforme au programme officiel MINESEC}\par}
    \end{minipage}\hfill
    {\poplogo\fontsize{22}{22}\selectfont\color{popcream}OnBuch\color{poporange}+}
  \end{tcolorbox}%
}

% ============================================================
% Page de garde
% #1 titre · #2 matière · #3 niveau · #4 module · #5 n° leçon · #6 durée ·
% #7 description / objectifs (texte ou itemize)
% ============================================================
\newcommand{\pagegarde}[7]{%
  \thispagestyle{empty}
  \newgeometry{a4paper,margin=1.7cm,top=1.6cm,bottom=1.4cm}%
  \begin{tikzpicture}[remember picture,overlay]
    \fill[poporange!18] ([xshift=-3.2cm,yshift=-3.6cm]current page.north east) circle (4.6cm);
    \fill[poppurple!14] ([xshift=2.4cm,yshift=7.6cm]current page.south west) circle (3.8cm);
  \end{tikzpicture}%
  {\poplogo\fontsize{30}{30}\selectfont\color{popink}OnBuch\color{poporange}+}\hfill
  \raisebox{0.6ex}{\poppill{Cours signé \textbullet\ Premium}{popink}{popcream}}\par
  \vspace{1pt}\popsquiggle{poppurple}\par
  \vspace{8pt}
  \begin{tcolorbox}[enhanced,colback=poporange,colframe=popink,boxrule=2.8pt,arc=13pt,
      drop shadow=popink,left=18pt,right=18pt,top=12pt,bottom=13pt,after skip=10pt]
    \poppill{#2 \textbullet\ #3}{popink}{popcream}\hspace{5pt}\poppill{#4}{white}{popink}\par\vspace{8pt}
    {\popdisplay\fontsize{14}{14}\selectfont\color{popink}LEÇON #5}\par\vspace{4pt}
    {\raggedright\hyphenpenalty=10000\exhyphenpenalty=10000\popdisplay\fontsize{25}{27}\selectfont\color{popcream}\MakeUppercase{#1}\par}
    \vspace{8pt}
    {\popxbold\fontsize{9.5}{9.5}\selectfont\color{popink}\MakeUppercase{Durée conseillée : #6}}
  \end{tcolorbox}
  \begin{tcolorbox}[enhanced,colback=white,colframe=popink,boxrule=2.2pt,arc=10pt,
      drop shadow=popink,left=16pt,right=16pt,top=10pt,bottom=10pt,after skip=8pt]
    \poppill{Dans cette leçon}{poppurple}{white}\par\vspace{5pt}
    \popsemi\fontsize{9.3}{12.6}\selectfont\color{popink}#7
  \end{tcolorbox}
  \noindent\begin{minipage}[t]{0.487\linewidth}
    \begin{tcolorbox}[enhanced,colback=popgreen,colframe=popink,boxrule=2.2pt,arc=10pt,
        drop shadow=popink,left=11pt,right=11pt,top=10pt,bottom=10pt,height=3.1cm,valign=center]
      \tcbox[colback=white,colframe=popink,boxrule=1.3pt,arc=5pt,boxsep=0pt,nobeforeafter,
        left=3pt,right=3pt,top=3pt,bottom=3pt,valign=center]{\qrcode[height=1.9cm]{https://play.google.com/store/apps/details?id=com.luvvix.onbuch&hl=fr}}%
      \hspace{8pt}\begin{minipage}[c]{0.5\linewidth}
        {\popdisplay\fontsize{11}{12.5}\selectfont\color{white}\MakeUppercase{Télécharge OnBuch}}\par\vspace{4pt}
        {\raggedright\popsemi\fontsize{8.4}{11}\selectfont\color{white}Gratuit \textbullet\ Google Play\\Tuteur IA, annales, concours, résultats\par}
      \end{minipage}
    \end{tcolorbox}
  \end{minipage}\hfill
  \begin{minipage}[t]{0.487\linewidth}
    \begin{tcolorbox}[enhanced,colback=poppurple,colframe=popink,boxrule=2.2pt,arc=10pt,
        drop shadow=popink,left=11pt,right=11pt,top=10pt,bottom=10pt,height=3.1cm,valign=center]
      \tikz[baseline=-0.7ex]\node[circle,fill=white,draw=popink,line width=1.3pt,minimum size=1.6cm,inner sep=0]{\popdisplay\fontsize{24}{24}\selectfont\color{poppurple}+};%
      \hspace{8pt}\begin{minipage}[c]{0.6\linewidth}
        {\popdisplay\fontsize{11}{12.5}\selectfont\color{white}\MakeUppercase{Passe à OnBuch+}}\par\vspace{4pt}
        {\raggedright\popsemi\fontsize{8.4}{11}\selectfont\color{white}Tous les cours signés, Tuteur IA illimité, fiches PDF — onglet OnBuch+ de l'app\par}
      \end{minipage}
    \end{tcolorbox}
  \end{minipage}
  \vspace{8pt}
  \popcontenu
  \vfill
  \signatureonbuch
  \restoregeometry
  \newpage
}

% Rangée « Ce cours contient » (page de garde)
\newcommand{\poptile}[3]{% #1 couleur #2 gros texte #3 légende
  \begin{tcolorbox}[enhanced,colback=#1,colframe=popink,boxrule=2pt,arc=9pt,shadow={2.5pt}{-2.5pt}{0pt}{popink},
      left=6pt,right=6pt,top=9pt,bottom=9pt,halign=center,valign=center,height=1.75cm,width=\linewidth,nobeforeafter]
    {\popdisplay\fontsize{12.5}{13}\selectfont\color{white}\MakeUppercase{#2}}\par\vspace{3pt}
    {\centering\popsemi\fontsize{7.8}{9.5}\selectfont\color{white}#3\par}
  \end{tcolorbox}}
\newcommand{\popcontenu}{%
  \par\noindent{\popxboldls\fontsize{7.5}{7.5}\selectfont\color{popmuted}CE COURS SIGNÉ CONTIENT}\par\vspace{7pt}
  \noindent\begin{minipage}[t]{0.235\linewidth}\poptile{popblue}{Cours}{complet, illustré, pas à pas}\end{minipage}\hfill
  \begin{minipage}[t]{0.235\linewidth}\poptile{poporange}{Méthodes}{et exercices résolus}\end{minipage}\hfill
  \begin{minipage}[t]{0.235\linewidth}\poptile{poppink}{Exercices}{type examen + corrigés}\end{minipage}\hfill
  \begin{minipage}[t]{0.235\linewidth}\poptile{popgreen}{Fiche bilan}{l'essentiel à retenir}\end{minipage}\par}

% Sommaire
\newtcolorbox{sommaire}{enhanced,colback=white,colframe=popink,boxrule=2.2pt,arc=10pt,
  drop shadow=popink,left=18pt,right=18pt,top=13pt,bottom=13pt,before skip=6pt,after skip=20pt,
  before upper={\poppill{Sommaire}{poppurple}{white}\par\vspace{9pt}\popsemi\fontsize{10.6}{14.5}\selectfont\color{popink}}}

% Signature de fin de cours — poussée tout en bas de la dernière page
\newcommand{\findecours}{%
  \par\vfill\nopagebreak
  \begin{center}\popsquiggle{poporange}\end{center}\vspace{4pt}
  \signatureonbuch
}
\AtBeginDocument{\def\llbracket{\mathopen{[\mkern-3.5mu[}}\def\rrbracket{\mathclose{]\mkern-3.5mu]}}} % intervalles d'entiers (glyphe ⟦ absent de la police)
% Glyphes absents de Fira Math : redéfinis à partir de glyphes disponibles
\AtBeginDocument{%
  \def\setminus{\mathbin{\backslash}}\let\smallsetminus\setminus
  \def\vdots{\mathord{\vbox{\baselineskip3.2pt\lineskiplimit0pt\kern4pt\hbox{.}\hbox{.}\hbox{.}}}}%
  \def\ddots{\mathinner{\mkern1mu\raise6.5pt\hbox{.}\mkern2mu\raise3.6pt\hbox{.}\mkern2mu\raise0.7pt\hbox{.}\mkern1mu}}%
  \def\triangle{\mathord{\scalebox{0.9}{$\symup{\Delta}$}}}%
  \def\nabla{\mathord{\raisebox{\depth}{\scalebox{1}[-1]{$\symup{\Delta}$}}}}%
  \def\checkmark{\popcheck{popdark}}%
  \def\star{\mathbin{\ast}}\def\bigstar{\mathord{\ast}}%
  \def\diamond{\mathbin{\scalebox{0.6}{$\Diamond$}}}%
  \def\bigcup{\mathop{\mathchoice{\raisebox{-0.3em}{\scalebox{1.7}{$\cup$}}}{\scalebox{1.25}{$\cup$}}{\cup}{\cup}}\displaylimits}%
  \def\bigcap{\mathop{\mathchoice{\raisebox{-0.3em}{\scalebox{1.7}{$\cap$}}}{\scalebox{1.25}{$\cap$}}{\cap}{\cap}}\displaylimits}%
  \def\lesseqqgtr{\mathrel{\lessgtr}}\def\bigoplus{\mathop{\oplus}\displaylimits}%
}
\providecommand{\ssi}{si et seulement si} % abréviation fréquente
\AtBeginDocument{\providecommand{\wideparen}{\overparen}} % arc de cercle

%% FIGFIX-BEGIN v5 — correctifs globaux des figures (docs/onbuch-generation/tools/figfix)
\usepackage{adjustbox}\usepackage{etoolbox}\tcbuselibrary{raster}
% toute figure TikZ/pgfplots plus large que la ligne est réduite à la largeur disponible
\BeforeBeginEnvironment{tikzpicture}{\begin{adjustbox}{max width=\linewidth}}
\AfterEndEnvironment{tikzpicture}{\end{adjustbox}}
% légendes pgfplots lisibles par-dessus les courbes
\pgfplotsset{every axis/.append style={legend style={fill=white,fill opacity=0.92,text opacity=1,draw=black!15,inner sep=3pt}}}
% étiquettes d'arêtes (syntaxe edge["..."]) avec fond blanc
\tikzset{every edge quotes/.append style={fill=white,inner sep=1.2pt}}
%% FIGFIX-END
\begin{document}

\pagegarde{Fonctions : généralités, limites et continuité}{Mathématiques}{1ère A}{Relations et opérations fondamentales dans l'ensemble des nombres réels}{NA02}{9 h}{Cette leçon introduit la parité des fonctions et ses conséquences graphiques, puis construit de manière intuitive les notions de limite (en un point, à droite et à gauche) et de continuité, par le calcul d'images, les tables de valeurs et la lecture graphique. Elle fournit les méthodes de justification attendues au Probatoire.
\vspace{4pt}
\begin{multicols}{2}\small
\begin{itemize}
\item À la fin de cette leçon, je sais définir une fonction paire et une fonction impaire sur un domaine symétrique par rapport à zéro
\item À la fin de cette leçon, je sais justifier algébriquement qu'une fonction est paire ou impaire sur un domaine borné symétrique par rapport à zéro
\item À la fin de cette leçon, je sais exploiter la courbe représentative d'une fonction pour conjecturer sa parité (symétrie par rapport à l'axe des ordonnées ou à l'origine)
\item À la fin de cette leçon, je sais conjecturer la limite d'une fonction en un réel par lecture graphique ou à l'aide d'une table de valeurs
\item À la fin de cette leçon, je sais déterminer la limite d'une fonction en un réel, à droite et à gauche d'un réel
\item À la fin de cette leçon, je sais exploiter la courbe de x ↦ a/x pour déterminer les limites à gauche et à droite d'une fonction homographique en −d/c
\end{itemize}\end{multicols}}

\begin{sommaire}
\begin{multicols}{2}
\begin{enumerate}
\item Parité d'une fonction
\item Approche intuitive de la limite en un réel
\item Limites à droite et à gauche
\item Continuité en un réel et sur un intervalle
\item Synthèse : lien entre limite et continuité, méthodes de rédaction
\item Activité d'intégration
\item Méthodes et exercices résolus
\item Exercices
\item Corrigés des exercices
\item Fiche bilan
\end{enumerate}
\end{multicols}
\end{sommaire}

\begin{objectifs}
\begin{itemize}
\item À la fin de cette leçon, je sais définir une fonction paire et une fonction impaire sur un domaine symétrique par rapport à zéro
\item À la fin de cette leçon, je sais justifier algébriquement qu'une fonction est paire ou impaire sur un domaine borné symétrique par rapport à zéro
\item À la fin de cette leçon, je sais exploiter la courbe représentative d'une fonction pour conjecturer sa parité (symétrie par rapport à l'axe des ordonnées ou à l'origine)
\item À la fin de cette leçon, je sais conjecturer la limite d'une fonction en un réel par lecture graphique ou à l'aide d'une table de valeurs
\item À la fin de cette leçon, je sais déterminer la limite d'une fonction en un réel, à droite et à gauche d'un réel
\item À la fin de cette leçon, je sais exploiter la courbe de x ↦ a/x pour déterminer les limites à gauche et à droite d'une fonction homographique en −d/c
\item À la fin de cette leçon, je sais justifier qu'une fonction est continue en un réel ou sur un intervalle ]a ; b[ par lecture graphique ou en l'identifiant à une fonction continue du programme
\item À la fin de cette leçon, je sais rédiger correctement les notations de limite en évitant les écritures proscrites comme l/0
\end{itemize}
\end{objectifs}

\begin{prerequis}
\begin{itemize}
\item Notion de fonction, ensemble de définition, image et antécédent (Seconde)
\item Calcul d'images f(x) et lecture de tableaux de valeurs
\item Représentation graphique d'une fonction dans un repère orthonormé
\item Symétrie axiale et symétrie centrale (collège)
\item Fonctions de référence : x ↦ ax + b, x ↦ x², x ↦ a/x
\end{itemize}
\end{prerequis}

\newpage

\coursec{Parité d'une fonction}

À la gare routière de Mvan à Yaoundé, le prix d'un trajet en taxi-moto dépend de la distance parcourue : c'est une fonction. Un commerçant remarque que sa recette du matin et celle du soir se ressemblent étrangement : sa fonction de recette serait-elle \emph{paire} ? Et lorsque le compteur d'électricité affiche des valeurs de plus en plus proches d'un seuil, que devient la facture : vers quelle valeur \emph{tend}-elle ? Enfin, une courbe de température relevée à l'hôpital de district ne doit présenter aucun saut brutal pour être exploitable : on dira qu'elle est \emph{continue}. Cette leçon te donne les outils pour étudier la \cle{parité}, les \cle{limites} et la \cle{continuité} des fonctions. Nous commençons par la parité, une propriété de symétrie qui, bien exploitée, divise par deux le travail d'étude d'une fonction.

\courssub{Ensemble de définition et domaine symétrique par rapport à zéro}

\begin{definition}[Ensemble de définition]
Soit $f$ une fonction. L'\cle{ensemble de définition} de $f$, noté $D_f$, est l'ensemble des nombres réels $x$ pour lesquels $f(x)$ existe (c'est-à-dire pour lesquels le calcul de l'image est possible).
\end{definition}

Par exemple, la fonction $x \mapsto \dfrac{1}{x-2}$ est définie sur $\mathbb{R}\setminus\{2\}$, car on ne peut pas diviser par zéro. La fonction $x \mapsto \sqrt{x}$ est définie sur $[0\,;\,+\infty[$, car une racine carrée n'existe que pour un nombre positif ou nul.

Pour parler de parité, il faut une condition sur le \emph{domaine} lui-même : il doit être \emph{symétrique par rapport à zéro}. Intuitivement, cela signifie que chaque fois qu'un nombre $x$ appartient au domaine, son opposé $-x$ y appartient aussi. Le domaine est alors « en miroir » de part et d'autre de $0$.

\begin{definition}[Domaine symétrique par rapport à zéro]
Un ensemble $D \subset \mathbb{R}$ est \cle{symétrique par rapport à zéro} si :
\[ \text{pour tout } x \in D, \quad -x \in D. \]
\end{definition}

\begin{exemplebox}[Domaines symétriques et non symétriques]
\begin{itemize}
\item $[-3\,;\,3]$ est symétrique par rapport à $0$ : si $x \in [-3\,;\,3]$, alors $-x \in [-3\,;\,3]$.
\item $]-2\,;\,2[$ est symétrique par rapport à $0$.
\item $[-5\,;\,-1] \cup [1\,;\,5]$ est symétrique par rapport à $0$, bien qu'il soit formé de deux morceaux : l'opposé d'un élément de $[1\,;\,5]$ est dans $[-5\,;\,-1]$, et réciproquement.
\item $[-2\,;\,3]$ n'est \textbf{pas} symétrique par rapport à $0$ : $3$ appartient au domaine, mais $-3$ n'y appartient pas.
\item $\mathbb{R}$ et $\mathbb{R}^*$ sont symétriques par rapport à $0$.
\end{itemize}
\end{exemplebox}

\begin{popfigure}
\begin{tikzpicture}[scale=1.1]
\draw[->,popline] (-3.6,0) -- (3.9,0) node[below left] {$x$};
\foreach \x in {-3,-2,-1,0,1,2,3}{\draw[popline] (\x,0.08) -- (\x,-0.08) node[below,font=\small] {$\x$};}
% Domaine [-2;3]
\draw[line width=2.2pt,poporange] (-2,0.35) -- (3,0.35);
\fill[poporange] (-2,0.35) circle (1.6pt);
\fill[poporange] (3,0.35) circle (1.6pt);
\node[poporange,above] at (0.5,0.4) {domaine $[-2\,;\,3]$};
% Point 3 et son opposé manquant
\fill[popink] (3,0) circle (1.6pt) node[above,yshift=2pt] {$3 \in D$};
\draw[popink,fill=white] (-3,0) circle (1.8pt) node[above,yshift=2pt] {$-3 \notin D$};
\draw[<->,popmuted,dashed] (-3,-0.7) -- (3,-0.7) node[midway,below,font=\small] {symétrie rompue};
\end{tikzpicture}
\legende{Le domaine $[-2\,;\,3]$ n'est pas symétrique par rapport à $0$ : le point $3$ est dans le domaine, mais son opposé $-3$ n'y est pas. La parité y est impossible.}
\end{popfigure}

\begin{attention}[La condition préalable à ne jamais oublier]
Avant même de calculer $f(-x)$, il faut \textbf{toujours} vérifier que le domaine est symétrique par rapport à zéro. Si ce n'est pas le cas, la fonction ne peut être \textbf{ni paire, ni impaire}, et la question est close. Au Probatoire, conclure à la parité sans cette vérification est une erreur de raisonnement sanctionnée.
\end{attention}

\courssub{Fonction paire}

Observe la courbe de $x \mapsto x^2$ : si tu plies la feuille le long de l'axe des ordonnées, les deux moitiés de la courbe se superposent exactement. L'image de $2$ et celle de $-2$ sont identiques : $4$. C'est exactement cela, la parité.

\begin{definition}[Fonction paire]
Soit $f$ une fonction définie sur un ensemble $D$. On dit que $f$ est \cle{paire} si :
\[ \text{pour tout } x \in D, \quad -x \in D \quad \text{et} \quad f(-x) = f(x). \]
\end{definition}

Autrement dit, deux nombres opposés ont la \emph{même image}. Le nom vient des puissances paires : $x^2$, $x^4$, $x^6$... vérifient toutes $(-x)^2 = x^2$, $(-x)^4 = x^4$, etc.

\begin{exemplebox}[Exemple chiffré détaillé]
Montrons que la fonction $f$ définie sur $[-4\,;\,4]$ par $f(x) = 3x^2 - 1$ est paire.
\begin{enumerate}
\item \textbf{Domaine} : $[-4\,;\,4]$ est symétrique par rapport à $0$ : si $x \in [-4\,;\,4]$, alors $-x \in [-4\,;\,4]$.
\item \textbf{Calcul de $f(-x)$} : pour tout $x \in [-4\,;\,4]$,
\[ f(-x) = 3(-x)^2 - 1 = 3x^2 - 1 = f(x). \]
\item \textbf{Conclusion} : $f$ est paire sur $[-4\,;\,4]$.
\end{enumerate}
Vérification numérique : $f(2) = 3\times 4 - 1 = 11$ et $f(-2) = 3\times 4 - 1 = 11$. On a bien $f(-2) = f(2)$.
\end{exemplebox}

\begin{propriete}[Interprétation graphique de la parité — admise]
Une fonction $f$ est \cle{paire} si et seulement si sa courbe représentative, dans un repère orthogonal, est \cle{symétrique par rapport à l'axe des ordonnées}.
\end{propriete}

Justifions graphiquement cette propriété. Dire que $f(-x) = f(x)$ signifie que les points $M(x\,;\,f(x))$ et $M'(-x\,;\,f(x))$ sont tous les deux sur la courbe. Or ces deux points ont la même ordonnée et des abscisses opposées : ils sont symétriques par rapport à l'axe des ordonnées. La courbe contient donc, avec chacun de ses points, son symétrique par rapport à cet axe.

\begin{popfigure}
\begin{tikzpicture}
\begin{axis}[
  width=0.82\linewidth, height=6.2cm,
  axis lines=middle, xlabel=$x$, ylabel=$y$,
  xmin=-2.6, xmax=2.6, ymin=-0.6, ymax=4.6,
  xtick={-2,-1,1,2}, ytick={1,2,3,4},
  grid=both, grid style={popcream},
  axis line style={popline}, tick label style={font=\small}
]
\addplot[popcurve,domain=-2:2,samples=200]{x^2};
% axe de symétrie mis en évidence
\draw[poporange,dashed,line width=1.1pt] (axis cs:0,-0.5) -- (axis cs:0,4.5) node[above,font=\small] {axe de symétrie};
% points symétriques
\fill[popink] (axis cs:1.5,2.25) circle (1.8pt) node[right,font=\small] {$M(1{,}5\,;\,2{,}25)$};
\fill[popink] (axis cs:-1.5,2.25) circle (1.8pt) node[left,font=\small] {$M'(-1{,}5\,;\,2{,}25)$};
\draw[popmuted,dashed] (axis cs:-1.5,2.25) -- (axis cs:1.5,2.25);
\end{axis}
\end{tikzpicture}
\legende{Courbe de $f(x)=x^2$ sur $[-2\,;\,2]$ : l'axe des ordonnées est axe de symétrie. Les points $M$ et $M'$ d'abscisses opposées ont la même ordonnée.}
\end{popfigure}

\courssub{Fonction impaire}

Observe maintenant la courbe de $x \mapsto x^3$ : le pliage le long de l'axe des ordonnées ne fonctionne plus, mais un demi-tour autour de l'origine superpose les deux moitiés. L'image de $2$ est $8$, celle de $-2$ est $-8$ : les images sont \emph{opposées}.

\begin{definition}[Fonction impaire]
Soit $f$ une fonction définie sur un ensemble $D$. On dit que $f$ est \cle{impaire} si :
\[ \text{pour tout } x \in D, \quad -x \in D \quad \text{et} \quad f(-x) = -f(x). \]
\end{definition}

Deux nombres opposés ont des images opposées. Le nom vient des puissances impaires : $x$, $x^3$, $x^5$... vérifient $(-x)^3 = -x^3$, $(-x)^5 = -x^5$, etc. La fonction inverse $x \mapsto \dfrac{1}{x}$ est aussi impaire sur $\mathbb{R}^*$.

\begin{propriete}[Interprétation graphique de l'imparité — admise]
Une fonction $f$ est \cle{impaire} si et seulement si sa courbe représentative est \cle{symétrique par rapport à l'origine} du repère.
\end{propriete}

Justifions graphiquement : dire que $f(-x) = -f(x)$ signifie que les points $M(x\,;\,f(x))$ et $M'(-x\,;\,-f(x))$ sont sur la courbe. Or le milieu du segment $[MM']$ est
\[ \left(\frac{x+(-x)}{2}\,;\,\frac{f(x)+(-f(x))}{2}\right) = (0\,;\,0), \]
c'est-à-dire l'origine $O$. Les points $M$ et $M'$ sont donc symétriques par rapport à $O$ : la courbe est invariante par la symétrie centrale de centre $O$.

\begin{popfigure}
\begin{tikzpicture}
\begin{axis}[
  width=0.82\linewidth, height=6.4cm,
  axis lines=middle, xlabel=$x$, ylabel=$y$,
  xmin=-2.1, xmax=2.1, ymin=-4.2, ymax=4.2,
  xtick={-1.5,-1,-0.5,0.5,1,1.5}, ytick={-3,-2,-1,1,2,3},
  grid=both, grid style={popcream},
  axis line style={popline}, tick label style={font=\small}
]
\addplot[popcurve,domain=-1.5:1.5,samples=200]{x^3};
\fill[popink] (axis cs:1.2,1.728) circle (1.8pt) node[above right,font=\small] {$M(x\,;\,g(x))$};
\fill[poppurple] (axis cs:-1.2,-1.728) circle (1.8pt) node[below left,font=\small] {$M'(-x\,;\,-g(x))$};
\draw[poporange,dashed,line width=1pt] (axis cs:-1.2,-1.728) -- (axis cs:0,0) -- (axis cs:1.2,1.728);
\fill[popdark] (axis cs:0,0) circle (1.5pt) node[below right,font=\small] {$O$};
\end{axis}
\end{tikzpicture}
\legende{Courbe de $g(x)=x^3$ sur $[-1{,}5\,;\,1{,}5]$ : les points $M$ et $M'$ sont symétriques par rapport à l'origine $O$, milieu du segment $[MM']$.}
\end{popfigure}

\courssub{Méthode de justification et exemples types}

\begin{methode}[Justifier la parité ou l'imparité d'une fonction]
Pour étudier la parité d'une fonction $f$ définie sur $D$ :
\begin{enumerate}
\item \textbf{Vérifier que le domaine est symétrique par rapport à $0$} : pour tout $x \in D$, a-t-on $-x \in D$ ? Si ce n'est pas le cas, conclure immédiatement : $f$ n'est ni paire ni impaire.
\item \textbf{Calculer $f(-x)$} pour un $x$ quelconque de $D$, en simplifiant soigneusement.
\item \textbf{Comparer} le résultat à $f(x)$ et à $-f(x)$.
\item \textbf{Conclure} :
\begin{itemize}
\item si $f(-x) = f(x)$ pour tout $x$, alors $f$ est \cle{paire} ;
\item si $f(-x) = -f(x)$ pour tout $x$, alors $f$ est \cle{impaire} ;
\item sinon, $f$ n'est ni paire ni impaire.
\end{itemize}
\end{enumerate}
\end{methode}

\begin{exoresolu}[Trois exemples de référence]
On considère les fonctions suivantes :
\begin{enumerate}
\item $f(x) = x^2$ sur $[-2\,;\,2]$ ;
\item $g(x) = x^3$ sur $[-1\,;\,1]$ ;
\item $h(x) = x^2 + x$ sur $[-3\,;\,3]$.
\end{enumerate}
Étudier la parité de chacune de ces fonctions.
\tcblower
\textbf{1. Fonction $f$.} Le domaine $[-2\,;\,2]$ est symétrique par rapport à $0$. Pour tout $x \in [-2\,;\,2]$ :
\[ f(-x) = (-x)^2 = x^2 = f(x). \]
Donc $f$ est \textbf{paire}. Sa courbe (une parabole) est symétrique par rapport à l'axe des ordonnées.

\textbf{2. Fonction $g$.} Le domaine $[-1\,;\,1]$ est symétrique par rapport à $0$. Pour tout $x \in [-1\,;\,1]$ :
\[ g(-x) = (-x)^3 = -x^3 = -g(x). \]
Donc $g$ est \textbf{impaire}. Sa courbe est symétrique par rapport à l'origine.

\textbf{3. Fonction $h$.} Le domaine $[-3\,;\,3]$ est symétrique par rapport à $0$. Pour tout $x \in [-3\,;\,3]$ :
\[ h(-x) = (-x)^2 + (-x) = x^2 - x. \]
Or $x^2 - x \neq h(x) = x^2 + x$ en général, et $x^2 - x \neq -h(x) = -x^2 - x$ en général. Pour le prouver proprement, un contre-exemple suffit : $h(1) = 2$ et $h(-1) = 0$. On n'a ni $h(-1) = h(1)$, ni $h(-1) = -h(1)$. Donc $h$ n'est \textbf{ni paire ni impaire}.
\end{exoresolu}

\begin{attention}[Deux pièges classiques]
\begin{itemize}
\item \textbf{Oublier la condition sur le domaine.} La fonction $x \mapsto x^2$ définie sur $[-2\,;\,3]$ n'est \emph{pas} paire, même si $(-x)^2 = x^2$ : le domaine n'est pas symétrique, car $3$ y est mais pas $-3$. Toujours vérifier le domaine \emph{en premier}.
\item \textbf{Conclure sur un seul calcul numérique.} Vérifier que $f(-1) = f(1)$ ne prouve pas que $f$ est paire : l'égalité doit valoir \emph{pour tout} $x$ du domaine. En revanche, un seul contre-exemple suffit pour prouver qu'une fonction n'est ni paire ni impaire.
\end{itemize}
\end{attention}

\courssub{Conjecture graphique de la parité sur un domaine borné}

Au Probatoire, on te donne souvent une portion de courbe tracée sur un domaine borné, et on te demande de \emph{conjecturer} la parité. La démarche est purement visuelle :

\begin{methode}[Conjecturer la parité à partir d'une courbe]
\begin{enumerate}
\item Repérer le domaine représenté et vérifier qu'il paraît symétrique par rapport à $0$.
\item Chercher un \cle{élément de symétrie} :
\begin{itemize}
\item si la courbe semble se superposer à elle-même par pliage le long de l'axe des ordonnées, on conjecture que $f$ est \textbf{paire} ;
\item si la courbe semble invariante par demi-tour autour de l'origine, on conjecture que $f$ est \textbf{impaire} ;
\item si aucune de ces deux symétries n'apparaît, on conjecture que $f$ n'est ni paire ni impaire.
\end{itemize}
\item Rédiger avec prudence : une lecture graphique donne une \emph{conjecture}, pas une preuve. Pour démontrer, il faut revenir au calcul de $f(-x)$.
\end{enumerate}
\end{methode}

\begin{exemplebox}[Lecture graphique]
Sur la figure de la fonction $f(x) = x^2$ tracée sur $[-2\,;\,2]$, on observe que la courbe se replie exactement sur elle-même le long de l'axe des ordonnées : on conjecture que $f$ est paire, ce que le calcul $f(-x) = (-x)^2 = x^2 = f(x)$ confirme. Sur la figure de $g(x) = x^3$, le demi-tour autour de $O$ laisse la courbe invariante : on conjecture que $g$ est impaire, confirmé par $g(-x) = -x^3 = -g(x)$.
\end{exemplebox}

\courssub{Intérêt pratique de la parité}

\begin{aretenir}[La parité divise le travail par deux]
Si $f$ est paire ou impaire sur un domaine symétrique $D$, il suffit d'étudier $f$ sur la \emph{moitié} du domaine (par exemple sur $D \cap [0\,;\,+\infty[$), puis de compléter par symétrie :
\begin{itemize}
\item fonction \cle{paire} : on complète la courbe par symétrie par rapport à l'axe des ordonnées ;
\item fonction \cle{impaire} : on complète la courbe par symétrie par rapport à l'origine.
\end{itemize}
Cela vaut pour le tableau de variations, le tableau de valeurs et le tracé de la courbe.
\end{aretenir}

Concrètement, pour tracer la courbe de $f(x) = 3x^2 - 1$ sur $[-4\,;\,4]$, il suffit de calculer $f(0)$, $f(1)$, $f(2)$, $f(3)$, $f(4)$, de tracer la moitié droite de la courbe, puis de la refléter dans l'axe des ordonnées. C'est ce genre d'économie que le commerçant de notre situation d'accroche exploite sans le savoir lorsqu'il constate que sa recette du matin et celle du soir « se ressemblent ».

\begin{exercice}[À toi de jouer]
\begin{enumerate}
\item La fonction $f$ définie sur $[-5\,;\,5]$ par $f(x) = x^4 - 2x^2$ est-elle paire ?
\item La fonction $g$ définie sur $[-2\,;\,2]$ par $g(x) = \dfrac{x}{1+x^2}$ est-elle impaire ?
\item La fonction $h$ définie sur $[-1\,;\,2]$ par $h(x) = x^2$ est-elle paire ? Justifier.
\end{enumerate}
\end{exercice}

\begin{corrige}[Exercice : À toi de jouer]
\textbf{1.} Le domaine $[-5\,;\,5]$ est symétrique par rapport à $0$. Pour tout $x$ :
\[ f(-x) = (-x)^4 - 2(-x)^2 = x^4 - 2x^2 = f(x). \]
Donc $f$ est paire.

\textbf{2.} Le domaine $[-2\,;\,2]$ est symétrique par rapport à $0$. Pour tout $x$ :
\[ g(-x) = \frac{-x}{1+(-x)^2} = \frac{-x}{1+x^2} = -g(x). \]
Donc $g$ est impaire.

\textbf{3.} Le domaine $[-1\,;\,2]$ n'est pas symétrique par rapport à $0$ : $2 \in [-1\,;\,2]$ mais $-2 \notin [-1\,;\,2]$. Donc $h$ n'est ni paire ni impaire, même si son expression ressemble à celle d'une fonction paire.
\end{corrige}

\begin{savaistu}[Pourquoi « paire » et « impaire » ?]
Le vocabulaire vient des polynômes : une fonction ne contenant que des puissances \emph{paires} de $x$ (comme $x^2$, $x^4$) est paire, et une fonction ne contenant que des puissances \emph{impaires} (comme $x$, $x^3$) est impaire. Mais attention : la notion dépasse largement les polynômes. La fonction cosinus, que tu rencontreras, est paire, et la fonction sinus est impaire. En physique, de nombreux phénomènes symétriques (diffusion de la chaleur, ondes) se modélisent par des fonctions paires ou impaires, ce qui simplifie énormément les calculs des ingénieurs, y compris ceux qui dimensionnent les réseaux électriques du Cameroun.
\end{savaistu}

\coursec{Approche intuitive de la limite en un réel}

Après avoir étudié la symétrie des courbes à travers la parité des fonctions, nous abordons maintenant une notion fondamentale de l'analyse : la \textbf{limite}. C'est l'outil qui permet de décrire le comportement d'une fonction au voisinage d'un point, même si la fonction n'est pas définie en ce point. Imagine que tu conduises sur la route Yaoundé--Douala : tu te rapproches de plus en plus de la ville d'Édéa sans forcément t'y arrêter. La limite, c'est la réponse à la question : « Quelle altitude (ou quelle valeur) mon altimètre affiche-t-il lorsque je suis \emph{arbitrairement proche} d'Édéa ? »

\courssub{L'idée directrice : se rapprocher sans atteindre}

Soit une fonction $f$ définie sur un ensemble $D \subset \mathbb{R}$ et un nombre réel $x_0$. L'idée centrale est d'observer les valeurs $f(x)$ lorsque $x$ prend des valeurs de plus en plus proches de $x_0$, \textbf{sans que $x$ soit nécessairement égal à $x_0$}. On dit que $x$ \emph{tend vers $x_0$} et on note $x \to x_0$.

\begin{definition}[Notion intuitive de limite finie]
On dit que la fonction $f$ \textbf{tend vers la limite finie $\ell$} (avec $\ell \in \mathbb{R}$) lorsque $x$ tend vers $x_0$ si, lorsque $x$ se rapproche de $x_0$ (sans l'atteindre), les valeurs $f(x)$ se rapprochent arbitrairement de $\ell$.
\end{definition}

\begin{aretenir}[Notation]
Si $f(x)$ se rapproche de $\ell$ quand $x \to x_0$, on écrit :
\[
\lim_{x \to x_0} f(x) = \ell \quad \text{ou} \quad f(x) \xrightarrow[x \to x_0]{} \ell
\]
On lit : « limite de $f(x)$ quand $x$ tend vers $x_0$ est égale à $\ell$ ».
\end{aretenir}

\courssub{Méthode 1 : La table de valeurs (approche numérique)}

Pour conjecturer (deviner) une limite, la méthode la plus concrète consiste à calculer $f(x)$ pour des valeurs de $x$ de plus en plus proches de $x_0$, \textbf{à la fois par valeurs inférieures ($x < x_0$) et par valeurs supérieures ($x > x_0$)}.

\begin{methode}[Conjecturer une limite par table de valeurs]
\begin{enumerate}
    \item Choisir des valeurs de $x$ proches de $x_0$ par défaut (ex: $x_0 - 0,1$, $x_0 - 0,01$, $x_0 - 0,001$, ...).
    \item Choisir des valeurs de $x$ proches de $x_0$ par excès (ex: $x_0 + 0,1$, $x_0 + 0,01$, $x_0 + 0,001$, ...).
    \item Calculer les images $f(x)$ correspondantes.
    \item Observer si les valeurs $f(x)$ se stabilisent vers un même nombre $\ell$ des deux côtés.
\end{enumerate}
\end{methode}

\begin{exemplebox}[Exemple : $f(x) = 2x + 1$ au voisinage de $x_0 = 3$]
Considérons la fonction affine $f(x) = 2x + 1$. Construisons une table de valeurs pour $x$ proche de $3$.

\begin{center}
\begin{tabularx}{\linewidth}{|c|*{6}{X|}}\hline
$x$ & $2,9$ & $2,99$ & $2,999$ & $3,001$ & $3,01$ & $3,1$ \\ \hline
$f(x) = 2x+1$ & $6,8$ & $6,98$ & $6,998$ & $7,002$ & $7,02$ & $7,2$ \\ \hline
\end{tabularx}
\end{center}

On observe que lorsque $x$ se rapproche de $3$ (que ce soit par valeurs inférieures en bleu ou supérieures en orange), $f(x)$ se rapproche de $7$. On conjecture donc :
\[
\lim_{x \to 3} (2x+1) = 7
\]
Remarque : ici $f(3) = 7$, la limite coïncide avec la valeur de la fonction. Ce ne sera pas toujours le cas !
\end{exemplebox}

\begin{popfigure}
\begin{tikzpicture}
\begin{axis}[
    width=\linewidth, height=7cm,
    axis lines=middle,
    xlabel=$x$, ylabel=$y$,
    xmin=2.5, xmax=3.5,
    ymin=6, ymax=8,
    xtick={2.5,2.8,2.9,3,3.1,3.2,3.5},
    ytick={6,6.5,6.8,7,7.2,7.5,8},
    grid=major, grid style={popline},
    tick label style={font=\footnotesize},
    label style={font=\small},
    clip=false
]
% Function
\addplot[popcurve, domain=2.5:3.5, samples=50, thick] {2*x+1};
% Point at x=3
\addplot[only marks, mark=*, mark size=2.5pt, popblue] coordinates {(3,7)};
% Dashed lines
\draw[popmuted, dashed] (axis cs:3,0) -- (axis cs:3,7);
\draw[popmuted, dashed] (axis cs:2.5,7) -- (axis cs:3,7);
% Zoom boxes visual representation (conceptual)
\draw[poporange, thick] (axis cs:2.9,6.8) rectangle (axis cs:3.1,7.2);
\node[poporange, anchor=south west, font=\small] at (axis cs:3.1,7.2) {Zoom $\times 10$};
\draw[poppurple, thick] (axis cs:2.99,6.98) rectangle (axis cs:3.01,7.02);
\node[poppurple, anchor=south west, font=\tiny] at (axis cs:3.01,7.02) {Zoom $\times 100$};
\end{axis}
\end{tikzpicture}
\legende{Représentation graphique de $y=2x+1$ : au voisinage de $x=3$, la courbe passe par $(3,7)$. Les rectangles illustrent les zooms successifs.}
\end{popfigure}

\courssub{Méthode 2 : Lecture graphique}

Si l'on dispose de la courbe représentative $\mathcal{C}_f$ de la fonction $f$ dans un repère orthonormé, la limite en $x_0$ se lit géométriquement.

\begin{propriete}[Lecture graphique de la limite]
Soit $\mathcal{C}_f$ la courbe de $f$. Pour connaître $\lim_{x \to x_0} f(x)$ :
\begin{enumerate}
    \item Repère l'abscisse $x_0$ sur l'axe des abscisses.
    \item Observe le comportement de la courbe lorsque l'abscisse $x$ se rapproche de $x_0$ \textbf{par la gauche} ($x < x_0$) et \textbf{par la droite} ($x > x_0$).
    \item Si la courbe se rapproche d'un même point de l'axe des ordonnées (de coordonnée $\ell$) des deux côtés, alors $\lim_{x \to x_0} f(x) = \ell$.
\end{enumerate}
\end{propriete}

\begin{attention}[Piège majeur : $f(x_0)$ vs $\lim_{x \to x_0} f(x)$]
\emph{Ne confonds jamais la valeur de la fonction en $x_0$ et la limite en $x_0$.}
\begin{itemize}
    \item La limite ne dépend \textbf{que} des valeurs de $f(x)$ pour $x \neq x_0$ proches de $x_0$.
    \item La fonction peut ne pas être définie en $x_0$ (trou dans le domaine) et pourtant avoir une limite en $x_0$.
    \item La fonction peut être définie en $x_0$ avec $f(x_0) \neq \ell$.
\end{itemize}
C'est l'erreur n°1 au Probatoire : écrire « comme $f(1)$ n'existe pas, la limite n'existe pas ». \textbf{FAUX !}
\end{attention}

\courssub{Le cas où la limite est infinie}

Il arrive que les valeurs $f(x)$ ne se rapprochent pas d'un nombre fini $\ell$, mais deviennent \emph{arbitrairement grandes} (en valeur absolue) lorsque $x \to x_0$.

\begin{definition}[Limite infinie]
\begin{itemize}
    \item On dit que $f(x)$ tend vers \textbf{$+\infty$} quand $x \to x_0$ si $f(x)$ devient aussi grand qu'on veut (positif) dès que $x$ est suffisamment proche de $x_0$. Notation : $\lim_{x \to x_0} f(x) = +\infty$.
    \item On dit que $f(x)$ tend vers \textbf{$-\infty$} quand $x \to x_0$ si $f(x)$ devient aussi grand qu'on veut en valeur absolue mais négatif. Notation : $\lim_{x \to x_0} f(x) = -\infty$.
\end{itemize}
\end{definition}

\begin{exemplebox}[Exemple classique : $f(x) = \frac{1}{x^2}$ au voisinage de $0$]
La fonction n'est pas définie en $0$ ($D_f = \mathbb{R}^*$). Calculons des valeurs proches de $0$ :

\begin{center}
\begin{tabular}{|c|cccc|cccc|}\hline
$x$ & $-0,1$ & $-0,01$ & $-0,001$ & $-0,0001$ & $0,0001$ & $0,001$ & $0,01$ & $0,1$ \\ \hline
$f(x)=\frac{1}{x^2}$ & $100$ & $10\,000$ & $10^6$ & $10^8$ & $10^8$ & $10^6$ & $10\,000$ & $100$ \\ \hline
\end{tabular}
\end{center}

Que $x$ approche $0$ par la gauche ($x<0$) ou par la droite ($x>0$), $x^2$ est positif et très petit, donc $\frac{1}{x^2}$ devient énorme. La limite est $+\infty$.

\begin{popfigure}
\begin{tikzpicture}
\begin{axis}[
    width=\linewidth, height=7cm,
    axis lines=middle,
    xlabel=$x$, ylabel=$y$,
    xmin=-1.5, xmax=1.5,
    ymin=0, ymax=12,
    xtick={-1,-0.5,0,0.5,1},
    ytick={0,2,4,6,8,10},
    grid=major, grid style={popline},
    tick label style={font=\footnotesize},
    label style={font=\small},
    clip=false,
    domain=-1.5:-0.3, samples=100
]
% Left branch
\addplot[popcurve, thick, domain=-1.5:-0.3] {1/x^2};
% Right branch
\addplot[popcurve, thick, domain=0.3:1.5] {1/x^2};
% Asymptote verticale
\draw[popred, dashed, thick] (axis cs:0,0) -- (axis cs:0,12) node[above, popred] {$x=0$ (asymptote)};
% Limit notation
\node[draw, fill=popblue, text=white, rounded corners, anchor=south] at (axis cs:0,12.5) {$\lim\limits_{x \to 0} \frac{1}{x^2} = +\infty$};
\end{axis}
\end{tikzpicture}
\legende{Courbe de $f(x)=\frac{1}{x^2}$ : les deux branches montent vers $+\infty$ au voisinage de $0$.}
\end{popfigure}
\end{exemplebox}

\begin{aretenir}[Distinction limite finie / limite infinie]
\begin{itemize}
    \item \textbf{Limite finie $\ell \in \mathbb{R}$} : $f(x)$ se stabilise près d'un nombre. La courbe a une \emph{asymptote horizontale} (ou passe par un point) au voisinage de $x_0$.
    \item \textbf{Limite $+\infty$} : $f(x)$ croît sans borne. La courbe monte verticalement (asymptote verticale).
    \item \textbf{Limite $-\infty$} : $f(x)$ décroît sans borne. La courbe descend verticalement.
\end{itemize}
\textbf{Attention :} Une limite infinie n'est \textbf{pas} une limite au sens strict (ce n'est pas un nombre réel). On dit que « la limite n'existe pas finie » ou « la fonction admet une limite infinie ».
\end{aretenir}

\courssub{Exemple fondamental : limite en un point hors du domaine}

C'est le cas le plus important pour la suite (dérivabilité, continuité). Une fonction peut avoir une limite en $x_0$ sans être définie en $x_0$.

\begin{exoresolu}[Conjecturer $\lim\limits_{x \to 1} \frac{x^2 - 1}{x - 1}$]
\emph{Énoncé :} Soit $f(x) = \frac{x^2 - 1}{x - 1}$. Déterminer $\lim_{x \to 1} f(x)$ par une table de valeurs.

\emph{Solution détaillée :}
\begin{enumerate}
    \item \textbf{Domaine :} $D_f = \mathbb{R} \setminus \{1\}$. La fonction n'est \textbf{pas définie en $1$} (dénominateur nul).
    \item \textbf{Table de valeurs :} On approche $1$ par valeurs inférieures et supérieures.
\end{enumerate}

\begin{center}
\begin{tabular}{|c|cccc|cccc|}\hline
$x$ & $0,9$ & $0,99$ & $0,999$ & $0,9999$ & $1,0001$ & $1,001$ & $1,01$ & $1,1$ \\ \hline
$x^2 - 1$ & $-0,19$ & $-0,0199$ & $-0,001999$ & $-0,00019999$ & $0,00020001$ & $0,002001$ & $0,0201$ & $0,21$ \\ \hline
$x - 1$ & $-0,1$ & $-0,01$ & $-0,001$ & $-0,0001$ & $0,0001$ & $0,001$ & $0,01$ & $0,1$ \\ \hline
$f(x)$ & $1,9$ & $1,99$ & $1,999$ & $1,9999$ & $2,0001$ & $2,001$ & $2,01$ & $2,1$ \\ \hline
\end{tabular}
\end{center}

\begin{popfigure}
\begin{tikzpicture}
\begin{axis}[
    width=\linewidth, height=7cm,
    axis lines=middle,
    xlabel=$x$, ylabel=$y$,
    xmin=0, xmax=2,
    ymin=0, ymax=3.5,
    xtick={0,0.5,1,1.5,2},
    ytick={0,1,2,3},
    grid=major, grid style={popline},
    tick label style={font=\footnotesize},
    label style={font=\small},
    clip=false
]
% Line y = x+1
\addplot[popcurve, thick, domain=0:2] {x+1};
% Hole at (1,2)
\addplot[only marks, mark=*, mark size=3pt, white] coordinates {(1,2)};
\addplot[only marks, mark=*, mark size=2.5pt, popred, thick] coordinates {(1,2)};
% Dashed lines to axes
\draw[popmuted, dashed] (axis cs:1,0) -- (axis cs:1,2);
\draw[popmuted, dashed] (axis cs:0,2) -- (axis cs:1,2);
% Label
\node[popred, anchor=south west, font=\small] at (axis cs:1.05,2.1) {Trou $(1,2)$};
\node[draw, fill=popblue, text=white, rounded corners, anchor=south west] at (axis cs:0.1,3.1) {$\lim\limits_{x \to 1} f(x) = 2$};
\end{axis}
\end{tikzpicture}
\legende{Graphique de $f(x)=\frac{x^2-1}{x-1}$ : la courbe est la droite $y=x+1$ avec un trou en $(1,2)$. La limite est l'ordonnée du trou.}
\end{popfigure}

\emph{Analyse :} Que $x$ approche $1$ par la gauche ($x<1$) ou par la droite ($x>1$), les valeurs de $f(x)$ se rapprochent de $2$.
\[
\boxed{\lim_{x \to 1} \frac{x^2 - 1}{x - 1} = 2}
\]

\emph{Note algébrique (pour comprendre) :} Pour $x \neq 1$, on peut factoriser : $x^2-1 = (x-1)(x+1)$. Donc $f(x) = \frac{(x-1)(x+1)}{x-1} = x+1$. La fonction coïncide avec la droite $y=x+1$ \emph{partout sauf en $1$} (où il y a un « trou »). La limite est l'ordonnée du point « manquant » sur la droite : $1+1=2$.
\end{exoresolu}

\begin{savaistu}[Histoire : La rigueur d'Epsilon-Delta]
L'intuition « se rapprocher autant qu'on veut » a satisfait les mathématiciens pendant des siècles (Newton, Leibniz, Euler). Mais au XIX\textsuperscript{e} siècle, des contre-exemples bizarres ont montré qu'il fallait une définition \emph{inattaquable}. C'est \textbf{Karl Weierstrass} (1815--1897), professeur à Berlin, qui a formalisé la définition $(\varepsilon, \delta)$ :
\[
\forall \varepsilon > 0,\ \exists \delta > 0,\ \text{tel que}\ 0 < |x - x_0| < \delta \implies |f(x) - \ell| < \varepsilon
\]
Cette définition rigoureuse est au programme de \textbf{Terminale C}. En Première, nous travaillons avec l'intuition, les tables de valeurs et la lecture graphique, qui sont les fondations indispensables avant la formalisation.
\end{savaistu}

\courssub{Synthèse de la section : Ce qu'il faut retenir}

\begin{aretenir}[Résumé : Approche intuitive de la limite en $x_0$]
\begin{enumerate}
    \item \textbf{Question posée :} Quelle valeur $f(x)$ approche-t-elle quand $x$ se rapproche de $x_0$ (sans l'atteindre) ?
    \item \textbf{Deux méthodes pour conjecturer :}
          \begin{itemize}
              \item \textbf{Numérique :} Table de valeurs ($x$ par défaut et par excès).
              \item \textbf{Graphique :} Lecture de l'ordonnée vers laquelle la courbe pointe.
          \end{itemize}
    \item \textbf{Notations :}
          \begin{itemize}
              \item $\lim_{x \to x_0} f(x) = \ell \in \mathbb{R}$ (limite finie).
              \item $\lim_{x \to x_0} f(x) = +\infty$ ou $-\infty$ (limite infinie).
          \end{itemize}
    \item \textbf{Point crucial :} L'existence de la limite en $x_0$ \textbf{ne nécessite pas} que $f$ soit définie en $x_0$, ni que $f(x_0) = \ell$.
\end{enumerate}
\end{aretenir}

\begin{exercice}[Entraînement : Conjecturer des limites]
\begin{enumerate}
    \item À l'aide d'une table de valeurs, conjecturer $\lim_{x \to 2} (3x - 5)$.
    \item Pour $f(x) = \frac{x^2 - 4}{x - 2}$, construire une table de valeurs au voisinage de $2$ et conjecturer la limite.
    \item On donne la courbe ci-contre (schéma mental) : une parabole $y = x^2$ avec un « trou » au point $(2, 4)$ et un point isolé tracé en $(2, 10)$. Quelle est $\lim_{x \to 2} f(x)$ ? Quelle est $f(2)$ ?
\end{enumerate}
\end{exercice}

\begin{corrige}[Exercice n°1]
\begin{enumerate}
    \item Table : $x=1,9 \to f=0,7$ ; $x=1,99 \to f=0,97$ ; $x=2,01 \to f=1,03$ ; $x=2,1 \to f=1,3$. Les valeurs tendent vers $1$. $\lim_{x \to 2} (3x-5) = 1$.
    \item $f(x) = \frac{(x-2)(x+2)}{x-2} = x+2$ pour $x \neq 2$. Table : $x=1,9 \to 3,9$ ; $x=1,99 \to 3,99$ ; $x=2,01 \to 4,01$. Limite $= 4$.
    \item La courbe suit la parabole $y=x^2$. Au voisinage de $2$, les points de la courbe ont pour ordonnée $\approx 4$. Le point isolé $(2,10)$ est $f(2)$. Donc $\lim_{x \to 2} f(x) = 4$ et $f(2) = 10$. La limite $\neq$ la valeur.
\end{enumerate}
\end{corrige}

\begin{attention}[Erreurs à ne pas commettre au Probatoire]
\begin{itemize}
    \item Écrire « $\lim_{x \to 1} \frac{x^2-1}{x-1} = \frac{0}{0} = 0$ » ou « n'existe pas car division par 0 ». \textbf{Interdit.} On ne substitue jamais $x_0$ directement si forme indéterminée.
    \item Confondre « limite infinie » et « pas de limite ». Dire « la limite est $+\infty$ » est une information précise (comportement connu), ce n'est pas « pas de limite ».
    \item Oublier d'approcher $x_0$ \textbf{des deux côtés} (gauche et droite) dans la table de valeurs. Si les comportements diffèrent, la limite n'existe pas (ce sera la section suivante : limites à droite et à gauche).
\end{itemize}
\end{attention}

Nous avons maintenant l'intuition de la limite « bilatère » (quand $x$ arrive de partout). Mais que se passe-t-il si la fonction se comporte différemment selon qu'on arrive par la gauche ou par la droite ? C'est l'objet de la section suivante : \textbf{Les limites à droite et à gauche}.

\coursec{Limites à droite et à gauche}

\courssub{Pourquoi distinguer l'approche par la gauche et par la droite ?}

Dans la section précédente, nous avons vu comment une fonction peut tendre vers une valeur finie $\ell$ lorsque $x$ se rapproche de $x_0$, sans distinction sur le « côté » par lequel $x$ arrive. Pourtant, prenons l'exemple concret d'un taxi à Douala : le tarif est souvent forfaitaire par palier de distance. Si tu parcours $2,9$ km, tu paies le tarif des $2$ km ; dès que tu franchis $3,0$ km, le tarif saute à la tranche supérieure. La fonction « prix en fonction de la distance » fait un \textbf{saut} en $x_0 = 3$. Si tu t'approches de $3$ km \emph{par en dessous} ($2,9 ; 2,99 ; 2,999$), le prix tend vers le tarif des $2$ km. Si tu t'approches \emph{par au-dessus} ($3,1 ; 3,01 ; 3,001$), le prix tend vers le tarif des $3$ km. Les deux valeurs limites sont \textbf{différentes}.

C'est exactement ce qui se passe pour la fonction inverse $x \mapsto \frac{1}{x}$ en $x_0 = 0$ : quand $x$ approche $0$ par valeurs négatives, $f(x)$ plonge vers $-\infty$ ; quand $x$ approche $0$ par valeurs positives, $f(x)$ s'envole vers $+\infty$. Il est donc indispensable de définir des limites \emph{unilatérales}.

\begin{popfigure}
\centering
\begin{tikzpicture}[scale=1.2, >=Stealth]
    % Axes
    \draw[->] (-3.5,0) -- (3.5,0) node[right] {$x$};
    \draw[->] (0,-3.5) -- (0,3.5) node[above] {$y$};
    % Graduations
    \foreach \x in {-3,-2,-1,1,2,3} \draw (\x,0.05) -- (\x,-0.05) node[below] {\small $\x$};
    \foreach \y in {-3,-2,-1,1,2,3} \draw (0.05,\y) -- (-0.05,\y) node[left] {\small $\y$};
    % Courbe y = 1/x (branche gauche)
    \draw[popblue, thick, domain=-3:-0.35, samples=100, smooth] plot (\x, {1/\x});
    % Courbe y = 1/x (branche droite)
    \draw[popblue, thick, domain=0.35:3, samples=100, smooth] plot (\x, {1/\x});
    % Flèches limites à gauche (x -> 0-)
    \draw[poporange, thick, ->] (-1.5, -0.67) -- (-0.5, -2);
    \draw[poporange, thick, ->] (-0.8, -1.25) -- (-0.4, -2.5);
    \node[poporange, anchor=north east] at (-0.35, -2.8) {$\lim\limits_{x\to 0^-} \frac{1}{x} = -\infty$};
    % Flèches limites à droite (x -> 0+)
    \draw[popgreen, thick, ->] (1.5, 0.67) -- (0.5, 2);
    \draw[popgreen, thick, ->] (0.8, 1.25) -- (0.4, 2.5);
    \node[popgreen, anchor=north west] at (0.35, 2.8) {$\lim\limits_{x\to 0^+} \frac{1}{x} = +\infty$};
    % Asymptote verticale pointillée
    \draw[popmuted, dashed] (0,-3.5) -- (0,3.5);
    \node[popmuted, anchor=north] at (0.2, -3.3) {$x=0$};
\end{tikzpicture}
\legende{Comportement de $x \mapsto 1/x$ au voisinage de $0$ : les deux branches divergent vers des infinis de signes opposés.}
\end{popfigure}

\courssub{Définitions rigoureuses : limites à gauche et à droite}

Soit $f$ une fonction définie sur un ensemble $D \subset \mathbb{R}$ et $x_0 \in \mathbb{R}$ une valeur d'adhérence de $D$ (c'est-à-dire qu'il existe des points de $D$ arbitrairement proches de $x_0$).

\begin{definition}[Limite à gauche en $x_0$]
On dit que $f$ admet la limite $\ell \in \mathbb{R} \cup \{-\infty, +\infty\}$ \textbf{à gauche} en $x_0$ si, lorsque $x$ tend vers $x_0$ en restant \textbf{strictement inférieur} à $x_0$ ($x < x_0$), $f(x)$ tend vers $\ell$.
On note :
\[
\lim_{x \to x_0,\ x < x_0} f(x) = \ell \qquad \text{ou plus simplement} \qquad \lim_{x \to x_0^-} f(x) = \ell.
\]
\end{definition}

\begin{definition}[Limite à droite en $x_0$]
On dit que $f$ admet la limite $\ell \in \mathbb{R} \cup \{-\infty, +\infty\}$ \textbf{à droite} en $x_0$ si, lorsque $x$ tend vers $x_0$ en restant \textbf{strictement supérieur} à $x_0$ ($x > x_0$), $f(x)$ tend vers $\ell$.
On note :
\[
\lim_{x \to x_0,\ x > x_0} f(x) = \ell \qquad \text{ou plus simplement} \qquad \lim_{x \to x_0^+} f(x) = \ell.
\]
\end{definition}

\begin{popfigure}
\centering
\begin{tikzpicture}[scale=1.5, >=Stealth]
    % Droite graduée
    \draw[thick, ->] (-3,0) -- (3,0);
    \foreach \x in {-2,-1,0,1,2} {
        \draw (\x, 0.1) -- (\x, -0.1) node[below] {\small $\x$};
    }
    \node[below] at (0, -0.4) {$x_0$};
    % Approche par la gauche
    \draw[popblue, thick, <->] (-2.5, 0.5) -- (-0.2, 0.5);
    \node[popblue, above] at (-1.35, 0.7) {Approche par la gauche ($x < x_0$)};
    \draw[popblue, ->] (-1.5, 0.3) -- (-1.5, 0.15);
    \draw[popblue, ->] (-1.0, 0.3) -- (-1.0, 0.15);
    \draw[popblue, ->] (-0.5, 0.3) -- (-0.5, 0.15);
    % Approche par la droite
    \draw[poporange, thick, <->] (0.2, -0.5) -- (2.5, -0.5);
    \node[poporange, below] at (1.35, -0.7) {Approche par la droite ($x > x_0$)};
    \draw[poporange, ->] (0.5, -0.3) -- (0.5, -0.15);
    \draw[poporange, ->] (1.0, -0.3) -- (1.0, -0.15);
    \draw[poporange, ->] (1.5, -0.3) -- (1.5, -0.15);
\end{tikzpicture}
\legende{Visualisation sur la droite réelle : $x \to x_0^-$ (par valeurs inférieures) et $x \to x_0^+$ (par valeurs supérieures).}
\end{popfigure}

\courssub{Lien avec la limite bilatère (usuelle)}

C'est le résultat fondamental qui relie les trois notions. Il est \textbf{admis} au programme de Première (sa démonstration rigoureuse $\varepsilon/\eta$ appartient au supérieur), mais tu dois savoir l'appliquer parfaitement.

\begin{propriete}[Caractérisation de la limite en un point]
Soit $f$ une fonction et $x_0 \in \mathbb{R}$. La fonction $f$ admet une limite $\ell$ (finie ou infinie) en $x_0$ \textbf{si et seulement si} :
\begin{enumerate}
    \item $f$ admet une limite à gauche en $x_0$,
    \item $f$ admet une limite à droite en $x_0$,
    \item ces deux limites sont \textbf{égales} : $\displaystyle \lim_{x \to x_0^-} f(x) = \lim_{x \to x_0^+} f(x) = \ell$.
\end{enumerate}
Dans ce cas, $\displaystyle \lim_{x \to x_0} f(x) = \ell$.
\end{propriete}

\begin{attention}[Piège classique au Probatoire]
Si $\lim_{x \to x_0^-} f(x) \neq \lim_{x \to x_0^+} f(x)$, alors \textbf{la limite en $x_0$ n'existe pas} (on note souvent « $\lim_{x \to x_0} f(x)$ n'existe pas » ou « pas de limite »).
Ne confonds jamais « limite infinie » ($+\infty$ ou $-\infty$) et « pas de limite ». Une limite infinie \textbf{existe} (dans $\overline{\mathbb{R}}$), elle décrit un comportement précis. L'absence de limite signifie un désaccord entre la gauche et la droite (ou un comportement oscillatoire).
\end{attention}

\courssub{Étude de la fonction inverse $x \mapsto \frac{a}{x}$ au voisinage de $0$}

C'est la « brique de base » pour toutes les fonctions homographiques. Rappelons le comportement de $x \mapsto \frac{1}{x}$ (déjà vu en Seconde, approfondi ici).

\begin{propriete}[Limites de la fonction inverse en $0$]
Soit $a \in \mathbb{R}^*$.
\begin{itemize}
    \item Si $a > 0$ :
    $\displaystyle \lim_{x \to 0^-} \frac{a}{x} = -\infty$ \quad et \quad $\displaystyle \lim_{x \to 0^+} \frac{a}{x} = +\infty$.
    \item Si $a < 0$ :
    $\displaystyle \lim_{x \to 0^-} \frac{a}{x} = +\infty$ \quad et \quad $\displaystyle \lim_{x \to 0^+} \frac{a}{x} = -\infty$.
\end{itemize}
Dans tous les cas, $\displaystyle \lim_{x \to 0} \frac{a}{x}$ \textbf{n'existe pas} (car les limites à gauche et à droite sont différentes).
\end{propriete}

\begin{exemplebox}[Application directe]
Déterminer les limites à gauche et à droite en $0$ de $f(x) = \frac{-3}{x}$.
\par\medskip
\textbf{Solution.} Ici $a = -3 < 0$.
\begin{itemize}
    \item Pour $x \to 0^-$ ($x$ négatif), $\frac{-3}{x} = \frac{\text{négatif}}{\text{négatif}} = \text{positif} \to +\infty$.
    \item Pour $x \to 0^+$ ($x$ positif), $\frac{-3}{x} = \frac{\text{négatif}}{\text{positif}} = \text{négatif} \to -\infty$.
\end{itemize}
Donc $\lim_{x \to 0^-} f(x) = +\infty$ et $\lim_{x \to 0^+} f(x) = -\infty$. La limite en $0$ n'existe pas.
\end{exemplebox}

\courssub{Application aux fonctions homographiques $x \mapsto \frac{ax+b}{cx+d}$}

C'est un savoir-faire \textbf{exigible au Probatoire} : déterminer les limites à gauche et à droite de la valeur interdite $x_0 = -\frac{d}{c}$ (le zéro du dénominateur).

\begin{methode}[Limites d'une fonction homographique en sa valeur interdite]
Soit $f(x) = \frac{ax+b}{cx+d}$ avec $c \neq 0$. La valeur interdite est $x_0 = -\frac{d}{c}$.
Pour trouver $\lim_{x \to x_0^-} f(x)$ et $\lim_{x \to x_0^+} f(x)$ :
\begin{enumerate}
    \item \textbf{Calculer le numérateur en $x_0$ :} $N_0 = a x_0 + b = a(-\frac{d}{c}) + b = \frac{bc - ad}{c}$.
    \item \textbf{Étudier le signe du dénominateur $cx+d$ de chaque côté de $x_0$ :}
    \begin{itemize}
        \item Si $c > 0$ : $cx+d < 0$ pour $x < x_0$ (gauche) et $cx+d > 0$ pour $x > x_0$ (droite).
        \item Si $c < 0$ : $cx+d > 0$ pour $x < x_0$ (gauche) et $cx+d < 0$ pour $x > x_0$ (droite).
    \end{itemize}
    \item \textbf{Conclure par la règle des signes :} $\frac{N_0}{\text{signe du dénominateur}}$ donne $+\infty$ ou $-\infty$.
\end{enumerate}
\emph{Astuce graphique :} La courbe est une hyperbole obtenue par translation/homothétie de $y = \frac{k}{x}$. La branche de gauche et la branche de droite sont dans des quadrants opposés par rapport au centre $(x_0, \frac{a}{c})$.
\end{methode}

\begin{exoresolu}[Exemple type Probatoire : $f(x) = \frac{2x+1}{x-3}$ en $x_0 = 3$]
\textbf{Énoncé.} Étudier les limites à gauche et à droite de $f(x) = \frac{2x+1}{x-3}$ en $x_0 = 3$.
\tcblower
\textbf{Solution détaillée.}

\textbf{1. Identifier la forme.}
$f$ est une fonction homographique. Le dénominateur s'annule pour $x-3=0 \iff x=3$. C'est bien $x_0 = 3$.

\textbf{2. Calculer la limite du numérateur en $3$.}
$N(x) = 2x+1$. $\lim_{x \to 3} N(x) = 2\times 3 + 1 = 7$.
Au voisinage de $3$, le numérateur est \textbf{strictement positif} (il tend vers $7 > 0$).

\textbf{3. Étudier le signe du dénominateur $D(x) = x-3$ de chaque côté.}
\begin{itemize}
    \item \textbf{À gauche ($x \to 3^-$) :} $x < 3 \implies x-3 < 0$. Le dénominateur tend vers $0$ \textbf{en restant négatif} ($0^-$).
    \item \textbf{À droite ($x \to 3^+$) :} $x > 3 \implies x-3 > 0$. Le dénominateur tend vers $0$ \textbf{en restant positif} ($0^+$).
\end{itemize}

\textbf{4. Conclure par la règle des signes.}
\begin{align*}
\lim_{x \to 3^-} f(x) &= \frac{7}{0^-} = -\infty. \\
\lim_{x \to 3^+} f(x) &= \frac{7}{0^+} = +\infty.
\end{align*}

\textbf{5. Conclusion sur la limite en $3$.}
Comme $\lim_{x \to 3^-} f(x) = -\infty \neq +\infty = \lim_{x \to 3^+} f(x)$, \textbf{la limite de $f$ en $3$ n'existe pas}.

\textbf{Interprétation graphique :} La droite $x=3$ est une \textbf{asymptote verticale}. La branche de gauche (pour $x<3$) descend vers $-\infty$, la branche de droite (pour $x>3$) monte vers $+\infty$.
\end{exoresolu}

\begin{popfigure}
\centering
\begin{tikzpicture}
\begin{axis}[
    width=\linewidth, height=8cm,
    axis lines=middle,
    xlabel=$x$, ylabel=$y$,
    xmin=-2, xmax=8,
    ymin=-15, ymax=15,
    xtick={-1,0,1,2,3,4,5,6,7},
    ytick={-10,-5,0,5,10},
    tick label style={font=\small},
    clip=false,
    domain=-2:2.8, samples=100, smooth,
]
% Branche gauche (x < 3)
\addplot[popblue, thick, domain=-2:2.9] {(2*x+1)/(x-3)};
% Branche droite (x > 3)
\addplot[popblue, thick, domain=3.1:8] {(2*x+1)/(x-3)};
% Asymptote verticale
\addplot[popmuted, dashed, domain=-15:15] (3, x);
\node[popmuted, anchor=south west] at (axis cs:3.1, -14) {$x=3$};
% Asymptote horizontale y=2
\addplot[popmuted, dashed, domain=-2:8] {2};
\node[popmuted, anchor=south west] at (axis cs:7.5, 2.2) {$y=2$};
% Flèches comportement
\draw[poporange, thick, ->] (axis cs:2.5, -5) -- (axis cs:2.9, -12);
\node[poporange, anchor=north] at (axis cs:2.5, -5.5) {$-\infty$};
\draw[popgreen, thick, ->] (axis cs:3.5, 10) -- (axis cs:3.1, 13);
\node[popgreen, anchor=south] at (axis cs:3.5, 10.5) {$+\infty$};
% Point centre symétrie (3,2)
\fill[poppink] (axis cs:3,2) circle (2pt);
\node[poppink, anchor=south east] at (axis cs:3,2) {$C(3,2)$};
\end{axis}
\end{tikzpicture}
\legende{Courbe de $f(x) = \frac{2x+1}{x-3}$. Asymptote verticale $x=3$ : branche gauche $\to -\infty$, branche droite $\to +\infty$. Asymptote horizontale $y=2$ (limite en $\pm\infty$). Centre de symétrie $C(3,2)$.}
\end{popfigure}

\begin{attention}[Notation proscrite : $\boldsymbol{\ell/0}$]
Au Probatoire (et en mathématiques supérieures), il est \textbf{strictement interdit} d'écrire :
\[
\lim_{x \to 3} \frac{2x+1}{x-3} = \frac{7}{0} \quad \text{ou} \quad \frac{7}{0^+}.
\]
C'est une \textbf{erreur de rédaction grave} qui coûte des points. On n'écrit \textbf{jamais} une fraction dont le dénominateur est $0$.
\par\medskip
\textbf{Rédige correctement :}
\begin{itemize}
    \item « Le numérateur tend vers $7 > 0$. »
    \item « Le dénominateur tend vers $0$ par valeurs positives (resp. négatives). »
    \item « Donc la limite est $+\infty$ (resp. $-\infty$). »
\end{itemize}
\end{attention}

\courssub{Exemple complet supplémentaire : cas où le numérateur est négatif}

\begin{exoresolu}[Étude de $g(x) = \frac{5-2x}{2x-4}$ en $x_0 = 2$]
\textbf{Énoncé.} Déterminer $\lim_{x \to 2^-} g(x)$ et $\lim_{x \to 2^+} g(x)$.
\tcblower
\textbf{Solution.}

\textbf{1. Valeur interdite :} $2x-4=0 \iff x=2$. C'est $x_0=2$.

\textbf{2. Numérateur en $2$ :} $N(x) = 5-2x$. $N(2) = 5-4 = 1 > 0$. Au voisinage de $2$, le numérateur est \textbf{positif}.

\textbf{3. Signe du dénominateur $D(x) = 2x-4 = 2(x-2)$.}
Le coefficient de $x$ est $2 > 0$.
\begin{itemize}
    \item $x \to 2^-$ ($x<2$) $\implies x-2 < 0 \implies D(x) < 0$. Dénominateur $\to 0^-$.
    \item $x \to 2^+$ ($x>2$) $\implies x-2 > 0 \implies D(x) > 0$. Dénominateur $\to 0^+$.
\end{itemize}

\textbf{4. Calcul des limites :}
\begin{align*}
\lim_{x \to 2^-} g(x) &= \frac{1}{0^-} = -\infty. \\
\lim_{x \to 2^+} g(x) &= \frac{1}{0^+} = +\infty.
\end{align*}
La limite en $2$ n'existe pas. Asymptote verticale $x=2$.
\end{exoresolu}

\courssub{Exercices d'entraînement}

\begin{exercice}[Limites unilatérales d'homographiques]
Déterminer les limites à gauche et à droite en la valeur interdite pour les fonctions suivantes :
\begin{enumerate}
    \item $h_1(x) = \frac{3x-2}{1-x}$ (valeur interdite $x_0 = 1$)
    \item $h_2(x) = \frac{x+4}{-2x-6}$ (valeur interdite $x_0 = -3$)
    \item $h_3(x) = \frac{-x+1}{x+2}$ (valeur interdite $x_0 = -2$)
\end{enumerate}
\end{exercice}

\begin{corrige}[Corrigé des exercices d'entraînement]
\textbf{1.} $h_1(x) = \frac{3x-2}{1-x}$, $x_0=1$.
Numérateur en $1$ : $3(1)-2 = 1 > 0$.
Dénominateur $1-x = -(x-1)$. Coefficient de $x$ : $-1 < 0$.
$x \to 1^- \implies x-1<0 \implies -(x-1)>0 \implies 0^+$. $\lim_{x\to 1^-} = \frac{1}{0^+} = +\infty$.
$x \to 1^+ \implies x-1>0 \implies -(x-1)<0 \implies 0^-$. $\lim_{x\to 1^+} = \frac{1}{0^-} = -\infty$.

\textbf{2.} $h_2(x) = \frac{x+4}{-2x-6}$, $x_0=-3$.
Numérateur en $-3$ : $-3+4 = 1 > 0$.
Dénominateur $-2x-6 = -2(x+3)$. Coefficient de $x$ : $-2 < 0$.
$x \to -3^- \implies x+3<0 \implies -2(x+3)>0 \implies 0^+$. $\lim_{x\to -3^-} = +\infty$.
$x \to -3^+ \implies x+3>0 \implies -2(x+3)<0 \implies 0^-$. $\lim_{x\to -3^+} = -\infty$.

\textbf{3.} $h_3(x) = \frac{-x+1}{x+2}$, $x_0=-2$.
Numérateur en $-2$ : $-(-2)+1 = 3 > 0$.
Dénominateur $x+2$. Coefficient de $x$ : $1 > 0$.
$x \to -2^- \implies x+2<0 \implies 0^-$. $\lim_{x\to -2^-} = \frac{3}{0^-} = -\infty$.
$x \to -2^+ \implies x+2>0 \implies 0^+$. $\lim_{x\to -2^+} = \frac{3}{0^+} = +\infty$.
\end{corrige}

\courssub{Synthèse de la méthode de rédaction (Check-list Probatoire)}

Quand tu dois calculer $\lim_{x \to x_0^\pm} \frac{ax+b}{cx+d}$ :

\begin{enumerate}
    \item \textbf{Écris la fonction} et identifie $x_0 = -\frac{d}{c}$.
    \item \textbf{Calcule la limite du numérateur} $N(x_0) = \frac{bc-ad}{c}$. Précise son signe (positif, négatif, ou nul — cas particulier vu en Terminale).
    \item \textbf{Étudie le signe du dénominateur} $cx+d$ pour $x < x_0$ et $x > x_0$ (regarde le signe de $c$).
    \item \textbf{Applique la règle des signes} : $\frac{+}{0^+}=+\infty$, $\frac{+}{0^-}=-\infty$, $\frac{-}{0^+}=-\infty$, $\frac{-}{0^-}=+\infty$.
    \item \textbf{Conclus distinctement} pour la gauche et la droite.
    \item \textbf{Réponds à la question posée} : « La limite en $x_0$ existe-t-elle ? » $\to$ Non, car limites à gauche et à droite différentes.
\end{enumerate}

\begin{savaistu}[Asymptotes verticales : vocabulaire de Terminale]
Le terme « asymptote verticale » est utilisé ici pour décrire la droite $x=x_0$ vers laquelle la courbe « s'éloigne à l'infini ». En Première, on parle de \textbf{comportement au voisinage de la valeur interdite}. L'étude formelle des asymptotes (verticales, horizontales, obliques) et leur tracé précis font partie du programme de Terminale. Pour l'instant, sache simplement que si $\lim_{x \to x_0^\pm} f(x) = \pm\infty$, la droite $x=x_0$ est une asymptote verticale.
\end{savaistu}

\begin{aretenir}[Résumé : Limites à droite et à gauche]
\begin{itemize}
    \item $\lim_{x \to x_0^-} f(x)$ : limite quand $x$ approche $x_0$ par valeurs $< x_0$.
    \item $\lim_{x \to x_0^+} f(x)$ : limite quand $x$ approche $x_0$ par valeurs $> x_0$.
    \item $\lim_{x \to x_0} f(x) = \ell \iff \lim_{x \to x_0^-} f(x) = \lim_{x \to x_0^+} f(x) = \ell$.
    \item Pour $f(x) = \frac{ax+b}{cx+d}$ en $x_0 = -\frac{d}{c}$ :
    \begin{itemize}
        \item Calculer $N_0 = ax_0+b$ (signe).
        \item Étudier signe de $cx+d$ à gauche et à droite (signe de $c$).
        \item Conclure $\pm\infty$.
    \end{itemize}
    \item \textbf{Interdit :} notation $\frac{\ell}{0}$. \textbf{Obligatoire :} raisonnement par signes ($0^+$, $0^-$).
\end{itemize}
\end{aretenir}

\coursec{Continuité en un réel et sur un intervalle}

Après avoir étudié les limites — ce qui se passe « aux abords » d'un point — nous abordons maintenant la \textbf{continuité}. C'est la propriété qui garantit qu'une fonction ne présente ni « saut » brutal, ni « trou » invisible, ni « asymptote » verticale en un point. En termes concrets : si tu traces la courbe de la fonction sur une feuille, tu ne lèves pas le crayon.

\courssub{Intuition : tracer sans lever le crayon}

Imagine que tu suis le tracé de la route nationale N3 entre Yaoundé et Bertoua. Si la route est \emph{continue}, tu roule sans interruption : pas de pont effondré (trou), pas de marche d'escalier brutale (saut), pas de ravin infranchissable (asymptote). En mathématiques, une fonction est \textbf{continue sur un intervalle} si sa courbe représentative se trace d'un seul tenant, sans lever le crayon, sur cet intervalle.

\begin{popfigure}
\begin{tikzpicture}[scale=0.9]
\begin{axis}[
    axis lines=middle,
    xlabel=$x$, ylabel=$y$,
    xmin=-2.5, xmax=3.5,
    ymin=-1.5, ymax=5.5,
    xtick={-2,-1,0,1,2,3},
    ytick={-1,0,1,2,3,4,5},
    grid=major, grid style={popline},
    width=\linewidth, height=0.6\linewidth,
    clip=false
]
\addplot[popcurve, domain=-2:3, samples=200, thick] {x^2 - 2*x + 1};
\addplot[only marks, mark=*, mark size=2, popblue] coordinates {(-2,9) (3,4)};
\node[popblue, anchor=south west] at (axis cs:-2,9) {$A(-2;9)$};
\node[popblue, anchor=north west] at (axis cs:3,4) {$B(3;4)$};
\node[popdark, align=center, font=\small] at (axis cs:0.5,5) {Courbe tracée d'un seul tenant $\rightarrow$ fonction continue sur $[-2;3]$};
\end{axis}
\end{tikzpicture}
\legende{Courbe de $f(x)=x^2-2x+1$ sur $[-2;3]$ : pas de saut, pas de trou.}
\end{popfigure}

À l'inverse, la fonction \textbf{partie entière} $E(x)$ (ou $\lfloor x \rfloor$) fait des « marches » à chaque entier : on est obligé de lever le crayon pour passer d'un palier au suivant.

\begin{popfigure}
\begin{tikzpicture}[scale=0.9]
\begin{axis}[
    axis lines=middle,
    xlabel=$x$, ylabel=$y$,
    xmin=-3.5, xmax=3.5,
    ymin=-3.5, ymax=3.5,
    xtick={-3,-2,-1,0,1,2,3},
    ytick={-3,-2,-1,0,1,2,3},
    grid=major, grid style={popline},
    width=\linewidth, height=0.6\linewidth,
    clip=false
]
\foreach \k in {-3,-2,-1,0,1,2,3} {
    \addplot[popcurve, domain=\k:\k+0.999, samples=2, thick] {\k};
    \addplot[only marks, mark=*, mark size=2, popblue, fill=popblue] coordinates {(\k,\k)};
    \addplot[only marks, mark=*, mark size=2, popblue, fill=white] coordinates {(\k+1,\k)};
}
\node[popdark, align=center, font=\small] at (axis cs:0,-3) {Sauts aux entiers $\rightarrow$ discontinuité};
\end{axis}
\end{tikzpicture}
\legende{Courbe de la fonction partie entière $x \mapsto \lfloor x \rfloor$ : sauts aux entiers.}
\end{popfigure}

\courssub{Définition rigoureuse : continuité en un point}

L'intuition du crayon est précieuse, mais elle ne suffit pas pour une démonstration. La définition formelle utilise la notion de limite que nous venons d'étudier.

\begin{definition}[Continuité en un réel]
Soit $f$ une fonction définie sur un ensemble $D \subset \mathbb{R}$ et $x_0 \in D$.
On dit que $f$ est \textbf{continue en $x_0$} si et seulement si :
\[
\lim_{x \to x_0} f(x) = f(x_0).
\]
Cette égalité suppose trois conditions \textbf{simultanées} :
\begin{enumerate}
    \item $f$ est \textbf{définie} en $x_0$ (donc $x_0 \in D$).
    \item $f$ admet une \textbf{limite finie} en $x_0$ (les limites à gauche et à droite existent et sont égales).
    \item Cette limite \textbf{coïncide} avec la valeur $f(x_0)$.
\end{enumerate}
\end{definition}

\begin{attention}[Les trois piliers de la continuité — Piège fréquent]
Ne vérifie \textbf{jamais} seulement la limite. Si $f$ n'est pas définie en $x_0$, elle n'est \textbf{pas} continue en $x_0$, même si la limite existe.
\textbf{Exemple classique :} $f(x) = \dfrac{x^2-1}{x-1}$ en $x_0=1$.
$\lim_{x \to 1} f(x) = \lim_{x \to 1} (x+1) = 2$, mais $f(1)$ n'existe pas (division par zéro).
$\rightarrow$ $f$ n'est \textbf{pas} continue en $1$ (c'est un « trou » qu'on peut boucher en prolongeant la fonction).
\end{attention}

\begin{popfigure}
\begin{tikzpicture}[scale=0.9]
\begin{axis}[
    axis lines=middle,
    xlabel=$x$, ylabel=$y$,
    xmin=-0.5, xmax=2.5,
    ymin=-0.5, ymax=3.5,
    xtick={0,1,2},
    ytick={0,1,2,3},
    grid=major, grid style={popline},
    width=\linewidth, height=0.6\linewidth,
    clip=false
]
% Branche gauche
\addplot[popcurve, domain=-0.5:0.99, samples=100, thick] {x+1};
% Branche droite
\addplot[popcurve, domain=1.01:2.5, samples=100, thick] {x+1};
% Trou en (1,2)
\addplot[only marks, mark=*, mark size=3, poporange, fill=white, thick] coordinates {(1,2)};
\node[poporange, anchor=south west] at (axis cs:1,2) {$\,(1;2)$ : trou (image manquante)};
% Point défini ailleurs si fonction par morceaux (ex: f(1)=0)
\addplot[only marks, mark=*, mark size=3, popgreen, fill=popgreen] coordinates {(1,0)};
\node[popgreen, anchor=north west] at (axis cs:1,0) {$\,(1;0)$ : valeur réelle $f(1)$};
\node[popdark, align=center, font=\small] at (axis cs:0.5,3.2) {$\lim_{x\to 1}f(x)=2 \neq f(1)=0$ \\ ou $f(1)$ non définie $\rightarrow$ Discontinuité};
\end{axis}
\end{tikzpicture}
\legende{Discontinuité en $x_0=1$ : un « trou » (limite $\neq$ valeur) ou point isolé.}
\end{popfigure}

\courssub{Continuité sur un intervalle}

\begin{definition}[Continuité sur un intervalle]
Soit $I$ un intervalle de $\mathbb{R}$ (ouvert, fermé, semi-ouvert, borné ou non).
Une fonction $f$ est \textbf{continue sur $I$} si et seulement si elle est continue en \textbf{tout point} de $I$.
\end{definition}

\begin{aretenir}[Lecture graphique de la continuité sur un intervalle]
Sur la courbe représentative $C_f$ :
\begin{itemize}
    \item $f$ est continue sur $]a;b[$ $\iff$ l'arc de courbe au-dessus de $]a;b[$ se trace sans lever le crayon (ni saut, ni trou, ni branche infinie).
    \item Un \textbf{saut} (limites latérales différentes) $\rightarrow$ discontinuité.
    \item Un \textbf{trou} (limite finie $\neq$ valeur, ou valeur non définie) $\rightarrow$ discontinuité.
    \item Une \textbf{branche infinie} (limite infinie) $\rightarrow$ discontinuité (asymptote verticale).
\end{itemize}
\end{aretenir}

\courssub{Fonctions continues de référence (admises au Probatoire)}

Tu dois connaître par cœur le tableau suivant. Ces fonctions sont continues \textbf{sur leur ensemble de définition}. Au Probatoire, tu peux écrire « $f$ est continue sur $\mathbb{R}$ car c'est une fonction polynôme » sans redémontrer.

\begin{propriete}[Fonctions continues usuelles — Admises]
\begin{center}
\begin{tabularx}{\linewidth}{|l|X|}\hline
\rowcolor{popblueL}
\textbf{Fonction} & \textbf{Ensemble de continuité ( = ensemble de définition)} \\ \hline
Fonction affine $x \mapsto ax+b$ & $\mathbb{R}$ \\ \hline
Fonction polynôme $x \mapsto a_nx^n+\dots+a_0$ & $\mathbb{R}$ \\ \hline
Fonction carré $x \mapsto x^2$ & $\mathbb{R}$ \\ \hline
Fonction valeur absolue $x \mapsto |x|$ & $\mathbb{R}$ \\ \hline
Fonction inverse $x \mapsto \dfrac{1}{x}$ & $]-\infty;0[ \cup ]0;+\infty[$ (continue sur \emph{chacun} de ces deux intervalles) \\ \hline
Fonction racine carrée $x \mapsto \sqrt{x}$ & $[0;+\infty[$ \\ \hline
Fonctions trigonométriques $\sin, \cos$ & $\mathbb{R}$ \\ \hline
Fonction tangente $\tan$ & $\mathbb{R} \setminus \{\frac{\pi}{2}+k\pi,\ k\in\mathbb{Z}\}$ \\ \hline
\end{tabularx}
\end{center}
\end{propriete}

\begin{attention}[Piège mortel : la fonction inverse n'est PAS continue sur $\mathbb{R}$ !]
C'est l'erreur n°1 au Probatoire.
$\bullet$ $x \mapsto \frac{1}{x}$ n'est \textbf{pas définie} en $0$, donc \textbf{pas continue} en $0$.
$\bullet$ Elle est continue sur $]-\infty;0[$ \textbf{et} sur $]0;+\infty[$ (deux intervalles \emph{séparés}).
$\bullet$ Ne jamais écrire : « $x \mapsto 1/x$ est continue sur $\mathbb{R}^*$ » (ce n'est pas un intervalle). Écrire : « continue sur $]-\infty;0[$ et sur $]0;+\infty[$ ».
\end{attention}

\courssub{Méthode : Justifier la continuité sur un intervalle}

\begin{methode}[Justifier que $f$ est continue sur un intervalle $I$]
\begin{enumerate}
    \item \textbf{Identifier la nature de $f$ :} Est-ce une fonction de référence (affine, polynôme, $1/x$, $\sqrt{x}$, $|x|$, $\sin$, $\cos$...) ? Une somme, produit, quotient, composée de telles fonctions ?
    \item \textbf{Appliquer les règles de calcul sur les fonctions continues (admises) :}
          \begin{itemize}
              \item Somme, différence, produit de fonctions continues $\rightarrow$ continue.
              \item Quotient $\frac{u}{v}$ $\rightarrow$ continue \textbf{là où $v(x) \neq 0$}.
              \item Composée $g \circ f$ $\rightarrow$ continue si $f$ continue en $x_0$ et $g$ continue en $f(x_0)$.
          \end{itemize}
    \item \textbf{Conclure :} Préciser l'intervalle exact de continuité (exclure les zéros des dénominateurs, les arguments négatifs de $\sqrt{\cdot}$, etc.).
    \item \textbf{Alternative graphique :} Si la courbe est donnée, décrire : « La courbe ne présente ni saut, ni trou, ni asymptote verticale sur $I$, donc $f$ est continue sur $I$. »
\end{enumerate}
\end{methode}

\courssub{Exemples résolus}

\begin{exoresolu}[Continuité d'une fonction rationnelle]
\textbf{Énoncé :} Étudier la continuité de la fonction $f(x) = \dfrac{2x+1}{x-3}$ sur $\mathbb{R}$.

\textbf{Solution détaillée :}
\begin{enumerate}
    \item $f$ est un quotient de deux fonctions polynomiales (donc continues sur $\mathbb{R}$) : $u(x)=2x+1$ et $v(x)=x-3$.
    \item La règle du quotient dit : $f = u/v$ est continue en tout point où le dénominateur $v(x) \neq 0$.
    \item $v(x)=0 \iff x-3=0 \iff x=3$.
    \item En $x=3$, $f$ n'est pas définie (division par zéro), donc $f$ n'est pas continue en $3$.
    \item Sur tout intervalle \textbf{ne contenant pas $3$}, $f$ est continue.
\end{enumerate}
\textbf{Conclusion :} $f$ est continue sur $]-\infty;3[$ et sur $]3;+\infty[$. Elle n'est pas continue en $3$ car non définie en $3$.
\tcblower
\textbf{Rédaction attendue au Probatoire :}
« $f$ est le quotient de deux fonctions polynomiales, donc continues sur $\mathbb{R}$.
Le dénominateur $x-3$ s'annule pour $x=3$.
Donc $f$ est continue sur $]-\infty;3[$ et sur $]3;+\infty[$.
$f$ n'est pas définie en $3$, donc $f$ n'est pas continue en $3$. »
\end{exoresolu}

\begin{exoresolu}[Continuité d'une fonction définie par morceaux]
\textbf{Énoncé :} Soit $f$ définie sur $\mathbb{R}$ par :
\[
f(x) = \begin{cases}
x^2 + 1 & \text{si } x \leq 1 \\
2x + 1 & \text{si } x > 1
\end{cases}
\]
Étudier la continuité de $f$ en $1$ et sur $\mathbb{R}$.

\textbf{Solution détaillée :}
\begin{enumerate}
    \item \textbf{Sur $]-\infty;1[$ et $]1;+\infty[$ :} $f$ coïncide avec des fonctions polynomiales ($x^2+1$ et $2x+1$), donc $f$ est continue sur ces deux intervalles ouverts.
    \item \textbf{En $x_0 = 1$ (point de raccord) :} Il faut vérifier les 3 conditions.
          \begin{itemize}
              \item $f(1) = 1^2+1 = 2$ (définie, 1\textsuperscript{ère} condition OK).
              \item Limite à gauche : $\lim_{x \to 1^-} f(x) = \lim_{x \to 1^-} (x^2+1) = 2$.
              \item Limite à droite : $\lim_{x \to 1^+} f(x) = \lim_{x \to 1^+} (2x+1) = 3$.
              \item Les limites latérales sont \textbf{différentes} ($2 \neq 3$), donc la limite $\lim_{x \to 1} f(x)$ \textbf{n'existe pas} (2\textsuperscript{ème} condition non remplie).
          \end{itemize}
    \item \textbf{Conclusion :} $f$ n'est \textbf{pas continue en $1$} (discontinuité de première espèce, ou « saut »).
          $f$ est continue sur $]-\infty;1[$ et sur $]1;+\infty[$, mais \textbf{pas sur $\mathbb{R}$}.
\end{enumerate}
\tcblower
\textbf{Remarque graphique :} La courbe présente un « saut » en $x=1$ : le trait de gauche arrive à $(1;2)$ (point plein), le trait de droite démarre de $(1;3)$ (point creux). Impossible de tracer sans lever le crayon.
\end{exoresolu}

\begin{exoresolu}[Continuité par composition : racine carrée]
\textbf{Énoncé :} Étudier la continuité de $f(x) = \sqrt{x^2 - 4}$ sur $\mathbb{R}$.

\textbf{Solution détaillée :}
\begin{enumerate}
    \item $f = g \circ h$ avec $h(x) = x^2-4$ (polynôme, continue sur $\mathbb{R}$) et $g(u) = \sqrt{u}$ (continue sur $[0;+\infty[$).
    \item $f$ est définie $\iff h(x) \geq 0 \iff x^2-4 \geq 0 \iff x \in ]-\infty;-2] \cup [2;+\infty[$.
    \item Pour la composition $g \circ h$ être continue en $x_0$, il faut $h$ continue en $x_0$ (vrai partout) et $g$ continue en $h(x_0)$.
    \item $g$ est continue sur $[0;+\infty[$. Or $h(x) \in [0;+\infty[$ exactement sur l'ensemble de définition de $f$.
    \item \textbf{Attention aux bornes $-2$ et $2$ :} $h(-2)=0$, $h(2)=0$. $g$ est continue \textbf{à droite} en $0$ (car définie sur $[0;+\infty[$). $h$ est continue partout. Donc $f$ est continue en $-2$ (continuité à droite) et en $2$ (continuité à gauche).
\end{enumerate}
\textbf{Conclusion :} $f$ est continue sur $]-\infty;-2]$ et sur $[2;+\infty[$.
\end{exoresolu}

\courssub{Lien limite — continuité : Synthèse}

\begin{aretenir}[Résumé : Limite vs Continuité]
\begin{center}
\begin{tabularx}{\linewidth}{|l|X|X|}\hline
\rowcolor{popblueL}
\textbf{Situation en $x_0$} & \textbf{Limite $\lim_{x\to x_0}f(x)$} & \textbf{Continuité en $x_0$} \\ \hline
Fonction « normale », courbe lisse & Existe, finie, $= f(x_0)$ & \textbf{OUI} \\ \hline
Trou : $f(x_0)$ non définie ou $\neq$ limite & Existe, finie ($\ell$) & \textbf{NON} (mais prolongeable par continuité si on pose $f(x_0)=\ell$) \\ \hline
Saut : limites latérales différentes & N'existe pas (finie) & \textbf{NON} \\ \hline
Asymptote verticale : limite infinie & $\pm\infty$ (n'existe pas au sens fini) & \textbf{NON} \\ \hline
\end{tabularx}
\end{center}
\end{aretenir}

\begin{savaistu}[Continuité et vie réelle : le signal MTN/Orange]
Ton téléphone affiche « 4G » ou « H+ ». Le débit de données $D(t)$ en fonction du temps $t$ est \emph{idéalement} continu : tu regardes une vidéo, le flux ne coupe pas. Mais dans la réalité (tunnel, zone d'ombre, saturation antenne), $D(t)$ fait des \textbf{sauts} (chute brutale à 0) ou présente des \textbf{trous} (micro-coupures). Les ingénieurs télécoms modélisent la \textbf{qualité de service (QoS)} en cherchant à maximiser les intervalles de continuité du débit. Une fonction « continue par morceaux » avec peu de sauts = bonne couverture réseau !
\end{savaistu}

\courssub{Exercices d'entraînement (style Probatoire)}

\begin{exercice}[Exercice 1 — Continuité d'une fonction rationnelle]
Soit $f(x) = \dfrac{x^2 - 4}{x - 2}$.
\begin{enumerate}
    \item Déterminer l'ensemble de définition $D_f$.
    \item Calculer $\lim_{x \to 2} f(x)$.
    \item $f$ est-elle continue en $2$ ? Justifier.
    \item Peut-on définir une fonction $\tilde{f}$ continue sur $\mathbb{R}$ qui coïncide avec $f$ sur $D_f$ ?
\end{enumerate}
\end{exercice}

\begin{exercice}[Exercice 2 — Fonction définie par morceaux et paramètre]
Soit $f_a$ définie sur $\mathbb{R}$ par :
\[
f_a(x) = \begin{cases}
x^2 + a & \text{si } x \leq 1 \\
3x - 1 & \text{si } x > 1
\end{cases}
\]
Déterminer la valeur du réel $a$ pour que $f_a$ soit continue sur $\mathbb{R}$.
\end{exercice}

\begin{exercice}[Exercice 3 — Lecture graphique]
On donne ci-dessous la courbe $C_f$ d'une fonction $f$ définie sur $[-3;4]$.
\begin{center}
\begin{tikzpicture}[scale=0.7]
\begin{axis}[
    axis lines=middle, xlabel=$x$, ylabel=$y$,
    xmin=-3.5, xmax=4.5, ymin=-2.5, ymax=3.5,
    xtick={-3,-2,-1,0,1,2,3,4}, ytick={-2,-1,0,1,2,3},
    grid=major, grid style={popline}, width=\linewidth, height=0.5\linewidth,
    clip=false
]
% Branche 1 : continue sur [-3; -1[ , f(x)=x. Limite en -1 = -1.
\addplot[popcurve, domain=-3:-1, samples=50, thick] {x};
\addplot[only marks, mark=*, fill=popblue, mark size=2] coordinates {(-3,-3)};
\addplot[only marks, mark=*, fill=white, mark size=2, thick, popblue] coordinates {(-1,-1)};
% Point isolé en (-1; 2)
\addplot[only marks, mark=*, fill=poporange, mark size=2] coordinates {(-1,2)};
% Branche 2 : continue sur ]-1; 2[ , f(x) = (1/3)x - 2/3. Limite en -1 = -1, en 2 = 0.
\addplot[popcurve, domain=-1:2, samples=50, thick] {x/3 - 2/3};
\addplot[only marks, mark=*, fill=white, mark size=2, thick, popblue] coordinates {(-1,-1) (2,0)};
% Point plein en (2; -1) : valeur f(2)
\addplot[only marks, mark=*, fill=popblue, mark size=2] coordinates {(2,-1)};
% Branche 3 : continue sur ]2; 4] , f(x) = 0.5x - 2. f(2)=-1, f(4)=0.
\addplot[popcurve, domain=2:4, samples=50, thick] {0.5*x - 2};
\addplot[only marks, mark=*, fill=white, mark size=2, thick, popblue] coordinates {(2,-1)};
\addplot[only marks, mark=*, fill=popblue, mark size=2] coordinates {(4,0)};
\end{axis}
\end{tikzpicture}
\end{center}
\begin{enumerate}
    \item Dire si $f$ est continue en $-1$, en $2$, en $3$. Justifier par lecture graphique.
    \item Sur quels intervalles $f$ est-elle continue ?
\end{enumerate}
\end{exercice}

\begin{corrige}[Exercice 1]
\begin{enumerate}
    \item $D_f = \mathbb{R} \setminus \{2\} = ]-\infty;2[ \cup ]2;+\infty[$.
    \item Pour $x \neq 2$, $f(x) = \dfrac{(x-2)(x+2)}{x-2} = x+2$.
          $\lim_{x \to 2} f(x) = \lim_{x \to 2} (x+2) = 4$.
    \item $f$ n'est \textbf{pas continue en $2$} car $2 \notin D_f$ ($f(2)$ n'existe pas). La première condition de la définition n'est pas remplie.
    \item Oui. On pose $\tilde{f}(x) = \begin{cases} x+2 & \text{si } x \neq 2 \\ 4 & \text{si } x = 2 \end{cases}$.
          Alors $\tilde{f}(x) = x+2$ sur $\mathbb{R}$, fonction polynôme, donc continue sur $\mathbb{R}$.
          C'est le \textbf{prolongement par continuité} de $f$ en $2$.
\end{enumerate}
\end{corrige}

\begin{corrige}[Exercice 2]
$f_a$ est continue sur $]-\infty;1[$ (polynôme $x^2+a$) et sur $]1;+\infty[$ (polynôme $3x-1$).
Pour la continuité en $1$ :
$f_a(1) = 1^2 + a = 1+a$.
$\lim_{x \to 1^-} f_a(x) = 1+a$.
$\lim_{x \to 1^+} f_a(x) = 3(1)-1 = 2$.
Continuité en $1 \iff 1+a = 2 \iff a = 1$.
\textbf{Réponse :} Pour $a=1$, $f_1$ est continue sur $\mathbb{R}$.
\end{corrige}

\begin{corrige}[Exercice 3]
\begin{enumerate}
    \item \textbf{En $-1$ :} $\lim_{x \to -1^-} f(x) = -1$ (branche gauche, point creux), $f(-1)=2$ (point plein orange), $\lim_{x \to -1^+} f(x) = -1$ (branche droite, point creux).
          La limite existe (égale à $-1$) mais est différente de $f(-1)=2$. \textbf{Discontinuité (trou « rebouché » ailleurs)}.
    \item \textbf{En $2$ :} $\lim_{x \to 2^-} f(x) = 0$ (point creux), $f(2) = -1$ (point plein bleu), $\lim_{x \to 2^+} f(x) = -1$.
          La limite à gauche ($0$) est différente de la limite à droite ($-1$). La limite n'existe pas. \textbf{Discontinuité (saut)}.
    \item \textbf{En $3$ :} La courbe est lisse, pas de trou, pas de saut, pas d'asymptote. \textbf{Continue en $3$}.
    \item $f$ est continue sur $[-3;-1[$, $]-1;2[$, $]2;4]$.
\end{enumerate}
\end{corrige}

\coursec{Synthèse : lien entre limite et continuité, méthodes de rédaction}

Dans la section précédente, tu as appris à reconnaître la continuité d'une fonction en un réel ou sur un intervalle, d'abord par lecture graphique (une courbe tracée « sans lever le crayon »), puis en identifiant la fonction à une fonction de référence continue du programme. Nous avons aussi découvert, dans les sections 2 et 3, la notion intuitive de limite en un réel, y compris les limites à gauche et à droite. Il est temps de rassembler toutes ces pièces du puzzle : tu vas voir que \cle{limite} et \cle{continuité} sont deux visages d'une même idée, et tu vas apprendre à rédiger une étude complète de fonction comme on l'exige au Probatoire.

\courssub{Le bilan fondamental : continuité et limite}

\courssub{Une intuition pour commencer}

Imagine que tu suives en taxi-moto le compteur de prix d'une course à Douala : le prix affiché est une fonction de la distance parcourue. Si, lorsque la distance s'approche de \SI{5}{km}, le prix affiché se rapproche de \num{1500} FCFA, et qu'\emph{exactement} à \SI{5}{km} le compteur affiche bien \num{1500} FCFA, tout est « normal » : le prix varie de façon continue. En revanche, si une taxe de péage de 200 FCFA s'ajoute brutalement au passage d'un point précis, le prix fait un \emph{saut} : la fonction n'est plus continue en ce point. Mathématiquement, la continuité en $x_0$, c'est exactement cela : ce que la fonction « annonce » quand $x$ s'approche de $x_0$ (la limite) coïncide avec ce qu'elle « fait réellement » en $x_0$ (l'image $f(x_0)$).

\begin{propriete}[Caractérisation de la continuité par la limite]
Soit $f$ une fonction et $x_0$ un réel. On a l'équivalence :
\[ f \text{ est continue en } x_0 \iff \begin{cases} f \text{ est définie en } x_0,\\ \lim\limits_{x \to x_0} f(x) \text{ existe (et est un réel fini)},\\ \lim\limits_{x \to x_0} f(x) = f(x_0). \end{cases} \]
Autrement dit, la continuité en $x_0$ exige la réunion de \textbf{trois conditions} : la fonction est définie en $x_0$, la limite en $x_0$ existe (les limites à gauche et à droite sont égales et finies), et cette limite vaut exactement $f(x_0)$.
\end{propriete}

\begin{aretenir}[Les trois conditions à vérifier]
Pour affirmer que $f$ est continue en $x_0$, il faut vérifier, dans l'ordre :
\begin{enumerate}
\item $x_0 \in D_f$ (la fonction est définie en $x_0$) ;
\item $\lim\limits_{x \to x_0^-} f(x) = \lim\limits_{x \to x_0^+} f(x) = \ell$ avec $\ell$ réel fini (la limite existe) ;
\item $\ell = f(x_0)$ (la limite coïncide avec l'image).
\end{enumerate}
Si l'une de ces trois conditions fait défaut, $f$ est \cle{discontinue} en $x_0$.
\end{aretenir}

\courssub{Les visages de la discontinuité}

La propriété précédente nous donne un moyen commode de \emph{classifier} les situations où la continuité échoue. Trois cas typiques se présentent.

\begin{definition}[Principaux cas de discontinuité en $x_0$]
Soit $f$ une fonction et $x_0$ un réel. On dit que $f$ est discontinue en $x_0$ dans chacune des situations suivantes :
\begin{itemize}
\item \textbf{Discontinuité par saut} : les limites à gauche et à droite existent, sont finies, mais sont différentes :
\[ \lim_{x \to x_0^-} f(x) \neq \lim_{x \to x_0^+} f(x). \]
La courbe « saute » d'une valeur à une autre en $x_0$.
\item \textbf{Discontinuité par limite infinie} : la limite à gauche ou à droite (ou les deux) est infinie. La droite d'équation $x = x_0$ est alors une \cle{asymptote verticale} à la courbe.
\item \textbf{Fonction non définie en $x_0$} : $x_0 \notin D_f$, donc $f(x_0)$ n'existe pas et la troisième condition n'a aucun sens.
\end{itemize}
\end{definition}

\begin{attention}[La notation $\dfrac{\ell}{0}$ est proscrite]
Lorsqu'un dénominateur s'annule en $x_0$ sans que le numérateur s'annule, on n'écrit \textbf{jamais} $\lim\limits_{x \to x_0} f(x) = \dfrac{\ell}{0}$. On raisonne sur le \emph{signe} du dénominateur à gauche et à droite de $x_0$, et on conclut par une limite infinie : $+\infty$ ou $-\infty$. Par exemple, au voisinage de $1$, l'expression $\dfrac{3}{x-1}$ tend vers $-\infty$ à gauche de $1$ (car $x - 1 < 0$) et vers $+\infty$ à droite de $1$ (car $x - 1 > 0$).
\end{attention}

\courssub{Méthodes de rédaction pour le Probatoire}

Au Probatoire, les questions sur les fonctions suivent presque toujours le même scénario : on te donne une fonction, on te demande son domaine, sa parité éventuelle, ses limites en un point remarquable, puis sa continuité. Voici la grille de rédaction complète.

\begin{methode}[Grille de rédaction : étude d'une fonction en $x_0$]
\begin{enumerate}
\item \textbf{Domaine de définition.} Résous les conditions d'existence (dénominateur non nul, quantité sous une racine positive ou nulle). Écris $D_f$ sous forme d'union d'intervalles.
\item \textbf{Limites à gauche et à droite en $x_0$.} Pour chaque côté, étudie le signe du dénominateur et conclus : $\lim\limits_{x \to x_0^-} f(x) = \ldots$ et $\lim\limits_{x \to x_0^+} f(x) = \ldots$
\item \textbf{Comparaison.} Compare les deux limites : si elles sont égales et finies, la limite en $x_0$ existe ; si elles diffèrent ou sont infinies, la limite en $x_0$ n'existe pas (au sens fini).
\item \textbf{Conclusion.} Conclus sur la continuité en $x_0$ grâce aux trois conditions, puis sur la continuité sur chaque intervalle du domaine en identifiant $f$ à une fonction continue du programme (polynôme, fonction rationnelle sur son domaine, etc.).
\end{enumerate}
\end{methode}

\begin{methode}[Rédaction type : justifier une parité]
Pour justifier que $f$ est \textbf{paire} sur un domaine borné symétrique par rapport à zéro (par exemple $[-3 ; 3]$) :
\begin{enumerate}
\item Vérifie la symétrie du domaine : pour tout $x \in D_f$, on a $-x \in D_f$.
\item Calcule $f(-x)$ et montre que $f(-x) = f(x)$ pour tout $x \in D_f$.
\item Conclus : $f$ est paire, donc sa courbe est symétrique par rapport à l'axe des ordonnées.
\end{enumerate}
Pour une fonction \textbf{impaire}, montre de même que $f(-x) = -f(x)$ et conclus par la symétrie de la courbe par rapport à l'origine du repère. Si tu trouves \emph{un seul} contre-exemple (par exemple $f(-1) \neq f(1)$ et $f(-1) \neq -f(1)$), conclus que $f$ n'est ni paire ni impaire.
\end{methode}

\begin{attention}[Erreurs fréquentes à bannir]
\begin{itemize}
\item Écrire « la fonction est continue car la courbe est jolie » ou « car ça se voit » : toute justification doit citer la \textbf{définition} (les trois conditions) ou identifier $f$ à une \textbf{fonction continue du programme} (polynôme, fonction rationnelle sur son domaine, valeur absolue, racine carrée sur son domaine).
\item Conclure sur la limite en $x_0$ après avoir calculé seulement la limite à droite : il faut \textbf{les deux} limites latérales.
\item Confondre « $f$ n'est pas définie en $x_0$ » et « $f$ n'est pas continue sur $]x_0 ; +\infty[$ » : une fonction peut être non définie en $x_0$ et pourtant parfaitement continue sur chaque intervalle de son domaine.
\item Oublier de vérifier la symétrie du domaine avant de conclure à la parité : sur $[-2 ; 3]$, la condition $f(-x) = f(x)$ ne peut pas être vérifiée pour $x = 2{,}5$ car $-2{,}5 \notin D_f$.
\end{itemize}
\end{attention}

\courssub{Tableau récapitulatif de la leçon}

\begin{popfigure}
\begin{tikzpicture}
\fill[popblueL, rounded corners=2pt] (0,4.6) rectangle (4.6,5.6);
\fill[poporangeL, rounded corners=2pt] (4.9,4.6) rectangle (9.5,5.6);
\fill[popgreenL, rounded corners=2pt] (9.8,4.6) rectangle (14.4,5.6);
\node[font=\bfseries, text=popdark] at (2.3,5.1) {Parité};
\node[font=\bfseries, text=popdark] at (7.2,5.1) {Limites};
\node[font=\bfseries, text=popdark] at (12.1,5.1) {Continuité};
\draw[popline, rounded corners=2pt] (0,0) rectangle (4.6,4.6);
\draw[popline, rounded corners=2pt] (4.9,0) rectangle (9.5,4.6);
\draw[popline, rounded corners=2pt] (9.8,0) rectangle (14.4,4.6);
\node[anchor=north west, text width=4.2cm, font=\small, align=left] at (0.15,4.45)
{Paire : $f(-x)=f(x)$\\[2pt] Impaire : $f(-x)=-f(x)$\\[2pt] Condition : $D_f$ symétrique par rapport \`a $0$.\\[2pt] Courbe : sym\'etrie / axe $(Oy)$ (paire) ou / origine $O$ (impaire).};
\node[anchor=north west, text width=4.2cm, font=\small, align=left] at (5.05,4.45)
{$\lim\limits_{x\to x_0} f(x)=\ell$\\[2pt] \`A gauche : $x\to x_0^-$\\ \`A droite : $x\to x_0^+$\\[2pt] La limite existe ssi les deux limites lat\'erales sont \'egales et finies.\\[2pt] Limite infinie : asymptote verticale $x=x_0$.};
\node[anchor=north west, text width=4.2cm, font=\small, align=left] at (9.95,4.45)
{$f$ continue en $x_0$ ssi :\\ d\'efinie en $x_0$, limite existante, $\lim = f(x_0)$.\\[2pt] R\'ef\'erences : polyn\^omes, fonctions rationnelles sur leur domaine.\\[2pt] Lecture : courbe sans saut.};
\end{tikzpicture}
\legende{Tableau récapitulatif : parité, limites et continuité — les trois piliers de la leçon.}
\end{popfigure}

\courssub{Exercice de synthèse guidé}

Nous allons maintenant mener ensemble l'étude complète d'une fonction rationnelle, en appliquant pas à pas la grille de rédaction. La figure suivante servira de support visuel à tout le raisonnement.

\begin{popfigure}
\begin{center}
\begin{tikzpicture}
\begin{axis}[
width=0.85\linewidth, height=7cm,
axis lines=middle,
xlabel={$x$}, ylabel={$y$},
xmin=-5, xmax=7, ymin=-6, ymax=8,
xtick={-4,-2,1,2,4,6}, ytick={-4,-2,2,4,6},
grid=both, grid style={popline!40},
legend style={at={(0.02,0.98)}, anchor=north west, font=\small}
]
\addplot[popcurve, domain=-5:0.7, samples=150] {(x+2)/(x-1)};
\addlegendentry{$f(x)=\dfrac{x+2}{x-1}$}
\addplot[popcurve, domain=1.35:7, samples=150] {(x+2)/(x-1)};
\addplot[poporange, dashed, thick] coordinates {(1,-6) (1,8)};
\addlegendentry{asymptote $x=1$}
\addplot[popgreen, dashed] coordinates {(-5,1) (7,1)};
\addplot[only marks, mark=*, popdark] coordinates {(0,-2) (-2,0) (4,2)};
\end{axis}
\end{tikzpicture}
\end{center}
\legende{Courbe de $f(x)=\dfrac{x+2}{x-1}$ : la droite $x=1$ est asymptote verticale ; la courbe « plonge » vers $-\infty$ à gauche de $1$ et « monte » depuis $+\infty$ à droite.}
\end{popfigure}

\begin{exoresolu}[Étude complète de $f(x) = \dfrac{x+2}{x-1}$ en $1$]
On considère la fonction $f$ définie par $f(x) = \dfrac{x+2}{x-1}$.
\begin{enumerate}
\item Détermine le domaine de définition de $f$.
\item La fonction $f$ est-elle paire ? impaire ?
\item Détermine les limites de $f$ à gauche et à droite de $1$. La limite de $f$ en $1$ existe-t-elle ?
\item $f$ est-elle continue en $1$ ? sur $]1 ; +\infty[$ ?
\end{enumerate}
\tcblower
\textbf{1. Domaine de définition.} $f(x)$ existe si et seulement si $x - 1 \neq 0$, c'est-à-dire $x \neq 1$. Donc :
\[ D_f = \mathbb{R} \setminus \{1\} = \left]-\infty ; 1\right[ \cup \left]1 ; +\infty\right[. \]

\textbf{2. Parité.} Le domaine $D_f$ n'est pas symétrique par rapport à $0$ : en effet, $1 \notin D_f$ alors que $-1 \in D_f$. La condition « pour tout $x \in D_f$, $-x \in D_f$ » n'est donc pas satisfaite. De plus, $f(-1) = \dfrac{1}{-2} = -\dfrac{1}{2}$, tandis que $f(1)$ n'existe pas : on ne peut avoir ni $f(-x) = f(x)$ ni $f(-x) = -f(x)$ pour tout $x$. \textbf{Conclusion : $f$ n'est ni paire ni impaire.}

\textbf{3. Limites en $1$.} Le numérateur $x + 2$ tend vers $3$ quand $x$ tend vers $1$ ; le dénominateur $x - 1$ tend vers $0$. Étudions le signe de $x - 1$ :
\begin{itemize}
\item si $x < 1$ (à gauche), alors $x - 1 < 0$, donc le quotient $\dfrac{x+2}{x-1}$ est négatif et son module devient arbitrairement grand :
\[ \lim_{x \to 1^-} f(x) = -\infty ; \]
\item si $x > 1$ (à droite), alors $x - 1 > 0$, donc :
\[ \lim_{x \to 1^+} f(x) = +\infty. \]
\end{itemize}
Les deux limites latérales sont différentes (et infinies) : \textbf{la limite de $f$ en $1$ n'existe pas}. Graphiquement, la droite d'équation $x = 1$ est une asymptote verticale à la courbe.

\textbf{4. Continuité.} En $1$ : $f$ n'est pas définie en $1$ (première condition en échec) et sa limite en $1$ n'existe pas. \textbf{$f$ est discontinue en $1$} (discontinuité par limite infinie).

Sur $]1 ; +\infty[$ : $f$ est une fonction rationnelle, quotient des polynômes $x + 2$ et $x - 1$, dont le dénominateur ne s'annule pas sur cet intervalle. Or toute fonction rationnelle est continue sur chaque intervalle de son domaine. \textbf{Donc $f$ est continue sur $]1 ; +\infty[$} (et de même sur $]-\infty ; 1[$).
\end{exoresolu}

\begin{exemplebox}[Vérification numérique par une table de valeurs]
Confirmons la limite à droite de $1$ par des calculs d'images, comme le recommande le programme :
\begin{center}
\begin{tabular}{|c|ccccc|}\hline
$x$ & $1{,}5$ & $1{,}1$ & $1{,}01$ & $1{,}001$ & $1{,}0001$ \\ \hline
$f(x)$ & $7$ & $31$ & $301$ & $\num{3001}$ & $\num{30001}$ \\ \hline
\end{tabular}
\end{center}
Quand $x$ se rapproche de $1$ par valeurs supérieures, $f(x)$ devient arbitrairement grand : cela illustre $\lim\limits_{x \to 1^+} f(x) = +\infty$. À gauche ($x = 0{,}9$ donne $f(x) = -29$ ; $x = 0{,}99$ donne $-299$), on lit de même $\lim\limits_{x \to 1^-} f(x) = -\infty$.
\end{exemplebox}

\begin{exercice}[À toi de jouer]
Soit $g$ la fonction définie par $g(x) = \dfrac{2x - 1}{x + 3}$.
\begin{enumerate}
\item Détermine $D_g$.
\item Étudie les limites de $g$ à gauche et à droite de $-3$.
\item $g$ est-elle continue en $-3$ ? sur $]-3 ; +\infty[$ ? Justifie.
\item La fonction $h$ définie sur $[-4 ; 4]$ par $h(x) = \dfrac{x}{x^2 + 1}$ est-elle paire, impaire, ou ni l'un ni l'autre ?
\end{enumerate}
\end{exercice}

\begin{corrige}[Exercice : fonction $g$ et fonction $h$]
\textbf{1.} $g(x)$ existe ssi $x + 3 \neq 0$, donc $D_g = \mathbb{R} \setminus \{-3\} = ]-\infty ; -3[ \cup ]-3 ; +\infty[$.

\textbf{2.} Le numérateur tend vers $-7$ en $-3$. Si $x \to -3^-$, alors $x + 3 < 0$ et $\lim\limits_{x \to -3^-} g(x) = +\infty$ (numérateur négatif, dénominateur négatif). Si $x \to -3^+$, alors $x + 3 > 0$ et $\lim\limits_{x \to -3^+} g(x) = -\infty$.

\textbf{3.} $g$ n'est pas définie en $-3$ et ses limites latérales sont infinies et distinctes : $g$ est discontinue en $-3$. En revanche, $g$ est une fonction rationnelle dont le dénominateur ne s'annule pas sur $]-3 ; +\infty[$ : elle y est continue.

\textbf{4.} Le domaine $[-4 ; 4]$ est symétrique par rapport à $0$. Pour tout $x \in [-4 ; 4]$ :
\[ h(-x) = \frac{-x}{(-x)^2 + 1} = \frac{-x}{x^2 + 1} = -h(x). \]
Donc $h$ est impaire ; sa courbe est symétrique par rapport à l'origine du repère.
\end{corrige}

\courssub{Pour aller plus loin : un aperçu de la Terminale}

\begin{savaistu}[Ce que la continuité te réserve en Terminale]
Tout ce que tu viens d'apprendre sera approfondi l'an prochain. Tu étudieras les \textbf{limites en $+\infty$ et en $-\infty$} (que fait $f(x)$ quand $x$ devient très grand ?), qui décrivent les branches infinies des courbes et les asymptotes horizontales — par exemple, notre fonction $f(x) = \dfrac{x+2}{x-1}$ tend vers $1$ quand $x$ tend vers $+\infty$, ce que la droite horizontale en pointillés verts de la figure laisse deviner. Surtout, tu rencontreras le célèbre \textbf{théorème des valeurs intermédiaires} : si $f$ est continue sur $[a ; b]$ et change de signe entre $a$ et $b$, alors l'équation $f(x) = 0$ admet au moins une solution dans $]a ; b[$. Ce théorème — conséquence directe de la continuité — permettra de prouver qu'une équation comme $x^3 + x - 1 = 0$ admet une solution, puis de l'encadrer avec la précision voulue. Retiens l'idée : \emph{une fonction continue ne peut pas passer d'une valeur positive à une valeur négative sans franchir zéro}, tout comme la montée du mont Cameroun ne peut pas te faire passer de \SI{1000}{m} à \SI{3000}{m} d'altitude sans passer par \SI{2000}{m}.
\end{savaistu}

\begin{aretenir}[L'essentiel de la synthèse]
\begin{itemize}
\item $f$ continue en $x_0$ $\iff$ $f$ définie en $x_0$, limite existante en $x_0$, et $\lim\limits_{x \to x_0} f(x) = f(x_0)$.
\item Discontinuité : saut (limites latérales finies distinctes), limite infinie (asymptote verticale), ou fonction non définie en $x_0$.
\item Rédaction type : domaine $\to$ limites à gauche et à droite $\to$ comparaison $\to$ conclusion sur la limite et la continuité.
\item Parité : toujours vérifier d'abord la symétrie du domaine, puis calculer $f(-x)$.
\item Toute justification de continuité s'appuie sur la définition ou sur une fonction continue du programme — jamais sur l'aspect de la courbe seul.
\end{itemize}
\end{aretenir}

\newpage

\coursec{Activité d'intégration}

\coursec{Fonctions : généralités, limites et continuité}

\courssub{Objectifs de l'activité}

À l'issue de cette activité d'intégration, tu seras capable de :
\begin{itemize}
\item justifier qu'une fonction est \cle{paire} sur un domaine borné symétrique par rapport à zéro et interpréter concrètement la symétrie de sa courbe ;
\item construire une \cle{table de valeurs} pour conjecturer la limite d'une fonction en un réel (ici, une limite à droite) ;
\item exploiter la courbe représentative d'une fonction pour justifier qu'elle est \cle{continue} ou non en un réel, et continue sur un intervalle ;
\item rédiger un argumentaire mathématique court et rigoureux à destination d'un non-spécialiste (le gérant d'une entreprise), compétence clé de l'approche par les compétences.
\end{itemize}

\courssub{Situation}

La société \textbf{CAMTEL-Mobile}, opérateur de téléphonie basé à Yaoundé, propose à ses abonnés un forfait international. Le service commercial a confié à son ingénieur la modélisation de deux quantités.

\textbf{Premier modèle : le coût d'une communication internationale.} Pour un appel international, le coût total d'une communication (en FCFA) dépend de sa durée $x$ exprimée en minutes. Il comprend un frais fixe de mise en relation de $50$ FCFA, auquel s'ajoute un coût d'interconnexion qui diminue lorsque la communication dure plus longtemps. L'ingénieur a retenu la fonction $f$ définie sur $]0\,;\,10]$ par :
\[ f(x) = 50 + \frac{100}{x} \quad \text{(en FCFA)}. \]
Ainsi, une communication de $1$ minute coûte $f(1) = 150$ FCFA, et une communication de $10$ minutes coûte $f(10) = 60$ FCFA.

\textbf{Second modèle : la recette horaire d'une agence.} L'agence CAMTEL-Mobile d'Akwa (Douala) enregistre ses recettes de vente de forfaits tout au long de la journée. En notant $x$ le nombre d'heures écoulées par rapport à midi ($x = -5$ correspond à 7 h du matin, $x = 0$ à midi, $x = 5$ à 17 h), la recette horaire (en milliers de FCFA) est modélisée par la fonction $g$ définie sur $[-5\,;\,5]$ par :
\[ g(x) = 25x^2 + 100. \]

Le gérant de l'agence s'étonne de deux choses : d'une part, la recette du matin semble « en miroir » de celle du soir ; d'autre part, les clients qui passent des appels internationaux très courts (quelques secondes) paient des montants qu'il juge « anormalement élevés ». Il te sollicite, toi et ton groupe, pour analyser mathématiquement ces deux phénomènes et lui rédiger un avis clair.

\courssub{Consignes}

Travail en groupes de quatre élèves. Durée : 55 minutes, suivies d'une restitution collective au tableau.

\begin{enumerate}
\item \textbf{Étude de la recette horaire.}
\begin{enumerate}
\item Vérifie que l'ensemble de définition de $g$ est symétrique par rapport à zéro.
\item Calcule $g(-x)$ pour tout $x \in [-5\,;\,5]$ et justifie que $g$ est une fonction paire.
\item Calcule $g(-4)$ et $g(4)$, puis $g(-2)$ et $g(2)$. Interprète la symétrie observée en termes de recette du matin et du soir. Quel élément de symétrie la courbe de $g$ possède-t-elle ?
\end{enumerate}
\item \textbf{Comportement du coût pour les communications très courtes.}
\begin{enumerate}
\item Complète la table de valeurs suivante (arrondir au FCFA près) :
\begin{center}
\begin{tabular}{|c|ccc|}\hline
\rowcolor{popblueL} $x$ (min) & $0{,}5$ & $0{,}1$ & $0{,}01$ \\ \hline
$f(x)$ (FCFA) & & & \\ \hline
\end{tabular}
\end{center}
\item Que devient $f(x)$ lorsque $x$ se rapproche de $0$ par valeurs supérieures ? Conjecture la limite de $f$ à droite en $0$, notée $\lim\limits_{\substack{x \to 0 \\ x > 0}} f(x)$.
\item La fonction $f$ est-elle définie en $0$ ? Peut-on parler de limite à gauche en $0$ pour $f$ ? Justifie.
\end{enumerate}
\item \textbf{Continuité.}
\begin{enumerate}
\item Trace la courbe représentative de $f$ sur $]0\,;\,10]$ dans un repère orthogonal (unités : 1 cm pour 1 minute en abscisse, 1 cm pour 50 FCFA en ordonnée), en utilisant les valeurs $f(0{,}5)$, $f(1)$, $f(2)$, $f(5)$, $f(10)$.
\item Peut-on prolonger cette courbe jusqu'au point d'abscisse $0$ « sans lever le crayon » ? Justifie que $f$ n'est pas continue en $0$.
\item Justifie que $f$ est continue sur l'intervalle $]0\,;\,10]$.
\end{enumerate}
\item \textbf{Avis au gérant.} Rédige, en cinq à huit lignes, un avis destiné au gérant expliquant, à l'aide de tes résultats sur la limite, pourquoi les communications très courtes coûtent si cher, et proposant une mesure commerciale raisonnable (par exemple une durée minimale facturée).
\end{enumerate}

\courssub{Ressources à mobiliser}

\begin{definition}[Fonction paire]
Soit $f$ une fonction définie sur un ensemble $D \subset \mathbb{R}$. On dit que $f$ est \cle{paire} sur $D$ si :
\begin{enumerate}
\item $D$ est symétrique par rapport à $0$ : $\forall x \in D,\; -x \in D$ ;
\item $\forall x \in D,\; f(-x) = f(x)$.
\end{enumerate}
La courbe représentative de $f$ est alors symétrique par rapport à l'axe des ordonnées (axe $Oy$).
\end{definition}

\begin{definition}[Fonction impaire]
Soit $f$ une fonction définie sur un ensemble $D \subset \mathbb{R}$. On dit que $f$ est \cle{impaire} sur $D$ si :
\begin{enumerate}
\item $D$ est symétrique par rapport à $0$ : $\forall x \in D,\; -x \in D$ ;
\item $\forall x \in D,\; f(-x) = -f(x)$.
\end{enumerate}
La courbe représentative de $f$ est alors symétrique par rapport à l'origine $O$ (centre de symétrie).
\end{definition}

\begin{propriete}[Limite à droite et à gauche en un réel]
Soit $f$ une fonction définie sur un intervalle et $x_0$ un réel.
\begin{itemize}
\item \textbf{Limite à droite en $x_0$ :} Si $f$ est définie sur un intervalle $]x_0\,;\,x_0+\alpha[$ ($\alpha>0$), on dit que $f$ admet la limite $\ell \in \mathbb{R}$ à droite en $x_0$ si, lorsque $x$ tend vers $x_0$ par valeurs $> x_0$, $f(x)$ se rapproche arbitrairement de $\ell$. On note $\lim\limits_{\substack{x \to x_0 \\ x > x_0}} f(x) = \ell$. Si $f(x)$ devient aussi grand que l'on veut, on note $\lim\limits_{\substack{x \to x_0 \\ x > x_0}} f(x) = +\infty$ ; s'il devient aussi petit que l'on veut (négatif), on note $-\infty$.
\item \textbf{Limite à gauche en $x_0$ :} Définie de manière analogue pour $x < x_0$, notée $\lim\limits_{\substack{x \to x_0 \\ x < x_0}} f(x)$.
\item \textbf{Limite bilatère :} $\lim\limits_{x \to x_0} f(x) = \ell$ si et seulement si les limites à droite et à gauche existent et sont égales à $\ell$.
\end{itemize}
\emph{Attention :} La notation $\frac{\ell}{0}$ ($\ell \in \mathbb{R}^*$) est \textbf{proscrite} par le programme. On raisonne toujours par tables de valeurs ou par opérations sur les limites.
\end{propriete}

\begin{propriete}[Limites de référence (fonction inverse)]
Pour la fonction $h(x) = \frac{1}{x}$ définie sur $\mathbb{R}^*$ :
\[ \lim\limits_{\substack{x \to 0 \\ x > 0}} \frac{1}{x} = +\infty, \qquad \lim\limits_{\substack{x \to 0 \\ x < 0}} \frac{1}{x} = -\infty, \qquad \lim\limits_{x \to +\infty} \frac{1}{x} = 0, \qquad \lim\limits_{x \to -\infty} \frac{1}{x} = 0. \]
\end{propriete}

\begin{definition}[Continuité en un point]
Soit $f$ une fonction définie sur un intervalle $I$ et $a \in I$. On dit que $f$ est \cle{continue en $a$} si :
\[ \lim\limits_{x \to a} f(x) = f(a). \]
Cela suppose que $f$ est définie en $a$ et que la limite (bilatère) existe et est finie.
Graphiquement : la courbe se trace « sans lever le crayon » au voisinage de $a$.
\end{definition}

\begin{propriete}[Continuité sur un intervalle / Fonctions continues de référence]
Une fonction est \cle{continue sur un intervalle} $I$ si elle est continue en tout point de $I$.
Les fonctions suivantes sont continues sur tout intervalle inclus dans leur ensemble de définition :
\begin{itemize}
\item les fonctions polynômes (donc sur $\mathbb{R}$) ;
\item la fonction inverse $x \mapsto \frac{1}{x}$ (sur $]-\infty\,;\,0[$ et $]0\,;\,+\infty[$) ;
\item la fonction racine carrée $x \mapsto \sqrt{x}$ (sur $[0\,;\,+\infty[$) ;
\item les fonctions trigonométriques $\sin$, $\cos$ (sur $\mathbb{R}$).
\end{itemize}
Toute somme, produit, quotient (dénominateur non nul), composée de fonctions continues sur un intervalle $I$ est continue sur $I$.
\end{propriete}

\begin{methode}[Montrer qu'une fonction est paire sur un domaine borné symétrique]
\begin{enumerate}
\item Vérifier que l'ensemble de définition $D$ est symétrique par rapport à $0$ : montrer que $x \in D \implies -x \in D$.
\item Pour tout $x \in D$, calculer $f(-x)$ et simplifier l'expression.
\item Conclure : si $f(-x) = f(x)$, alors $f$ est paire sur $D$ ; sa courbe est symétrique par rapport à l'axe des ordonnées.
\end{enumerate}
\end{methode}

\begin{methode}[Conjecturer une limite à l'aide d'un tableau de valeurs]
\begin{enumerate}
\item Choisir des valeurs de $x$ de plus en plus proches de $x_0$ (par la droite : $x > x_0$ ; par la gauche : $x < x_0$).
\item Calculer les images $f(x)$ correspondantes.
\item Observer la tendance des valeurs de $f(x)$ : se rapprochent-elles d'un nombre $\ell$ ? Deviennent-elles arbitrairement grandes ($+\infty$) ou petites ($-\infty$) ?
\item Formuler la conjecture : $\lim\limits_{\substack{x \to x_0 \\ x > x_0}} f(x) = \ell$ (ou $+\infty$, $-\infty$).
\end{enumerate}
\end{methode}

\begin{methode}[Étudier la continuité d'une fonction en un point et sur un intervalle]
\begin{enumerate}
\item \textbf{En un point $a$ :} Vérifier que $f$ est définie en $a$. Calculer $\lim\limits_{x \to a} f(x)$ (ou limites latérales si $a$ est au bord du domaine). Comparer la limite et $f(a)$. Si égales, $f$ est continue en $a$.
\item \textbf{Sur un intervalle $I$ :} Identifier $f$ comme combinaison (somme, produit, quotient, composée) de fonctions continues de référence sur $I$. Conclure par les opérations sur les fonctions continues.
\item \textbf{Graphiquement :} La courbe ne présente ni « trou », ni « saut », ni « branche infinie » (asymptote verticale) sur $I$.
\end{enumerate}
\end{methode}

\courssub{Démarche et corrigé commenté}

\begin{exoresolu}[Corrigé complet de l'activité CAMTEL-Mobile]
\textbf{Question 1.} Montrer que $g$ est paire sur $[-5\,;\,5]$ et interpréter.

\textbf{Question 2.} Construire la table de valeurs de $f$ près de $0^+$ et conjecturer la limite à droite en $0$.

\textbf{Question 3.} Étudier la continuité de $f$ en $0$ et sur $]0\,;\,10]$.

\textbf{Question 4.} Rédiger l'avis au gérant.
\tcblower
\textbf{1.a) Symétrie du domaine.} L'ensemble de définition est $D_g = [-5\,;\,5]$. Soit $x \in [-5\,;\,5]$. Alors $-5 \le x \le 5$. En multipliant par $-1$ (ce qui inverse l'ordre des inégalités), on obtient $-5 \le -x \le 5$, c'est-à-dire $-x \in [-5\,;\,5]$. Donc $D_g$ est bien symétrique par rapport à zéro. La première condition de la définition est remplie.

\textbf{1.b) Calcul de $g(-x)$ et parité.} Pour tout $x \in [-5\,;\,5]$ :
\[ g(-x) = 25(-x)^2 + 100 = 25x^2 + 100 = g(x). \]
Les deux conditions de la définition d'une fonction paire sont vérifiées : $g$ est une \cle{fonction paire} sur $[-5\,;\,5]$. Sa courbe $\mathcal{C}_g$ admet l'axe des ordonnées ($Oy$) comme axe de symétrie.

\begin{popfigure}
\begin{tikzpicture}[scale=0.6, y=0.0075cm, >=Stealth]
\def\xmin{-5.5} \def\xmax{5.5} \def\ymin{0} \def\ymax{800}
\draw[popline, thin] (\xmin,0) -- (\xmax,0) node[right] {$x$ (heures)};
\draw[popline, thin] (0,\ymin) -- (0,\ymax) node[above] {$g(x)$ (kFCFA)};
\foreach \x in {-5,-4,...,5} \draw (\x,2pt) -- (\x,-2pt) node[below, font=\tiny] {\x};
\foreach \y in {100,200,...,700} \draw (2pt,\y) -- (-2pt,\y) node[left, font=\tiny] {\y};
\draw[popblue, thick, domain=-5:5, samples=200] plot (\x, {25*\x*\x + 100});
\draw[poporange, dashed] (-4,500) -- (-4,0) node[below, font=\tiny] {-4};
\draw[poporange, dashed] (4,500) -- (4,0) node[below, font=\tiny] {4};
\draw[poporange, dashed] (-2,200) -- (-2,0) node[below, font=\tiny] {-2};
\draw[poporange, dashed] (2,200) -- (2,0) node[below, font=\tiny] {2};
\node[popblue, right] at (5,725) {$g(x)=25x^2+100$};
\node[poporange, anchor=south] at (0,500) {Symétrie par rapport à $Oy$};
\end{tikzpicture}
\legende{Courbe de la fonction $g$ (recette horaire) sur $[-5\,;\,5]$.}
\end{popfigure}

\textbf{1.c) Interprétation concrète.}
$g(-4) = 25 \times 16 + 100 = 500$ et $g(4) = 500$ ; $g(-2) = 25 \times 4 + 100 = 200$ et $g(2) = 200$.
$x = -4$ correspond à 8 h du matin, $x = 4$ à 16 h (4 h de l'après-midi). La recette est identique : \num{500} milliers de FCFA (soit \num{500000} FCFA).
$x = -2$ correspond à 10 h du matin, $x = 2$ à 14 h. La recette est identique : \num{200} milliers de FCFA.
L'activité de l'agence est symétrique par rapport à midi ($x=0$). La recette est minimale à midi ($g(0)=100$) et maximale aux heures d'ouverture/fermeture ($g(\pm 5) = 725$).

\textbf{2.a) Table de valeurs de $f$ près de $0^+$.}
\begin{align*}
f(0{,}5) &= 50 + \frac{100}{0{,}5} = 50 + 200 = 250 \text{ FCFA} ;\\
f(0{,}1) &= 50 + \frac{100}{0{,}1} = 50 + 1000 = 1050 \text{ FCFA} ;\\
f(0{,}01) &= 50 + \frac{100}{0{,}01} = 50 + 10000 = 10050 \text{ FCFA}.
\end{align*}
\begin{center}
\begin{tabular}{|c|ccc|}\hline
\rowcolor{popblueL} $x$ (min) & $0{,}5$ & $0{,}1$ & $0{,}01$ \\ \hline
$f(x)$ (FCFA) & $250$ & $1050$ & $10050$ \\ \hline
\end{tabular}
\end{center}

\textbf{2.b) Conjecture de la limite à droite en $0$.}
Lorsque $x$ se rapproche de $0$ par valeurs strictement positives ($x \to 0$, $x>0$), le quotient $\frac{100}{x}$ devient aussi grand que l'on veut (car le numérateur est constant positif et le dénominateur tend vers $0^+$). Par conséquent, $f(x) = 50 + \frac{100}{x}$ devient aussi grand que l'on veut. On conjecture :
\[ \lim\limits_{\substack{x \to 0 \\ x > 0}} f(x) = +\infty. \]
Ceci est cohérent avec le comportement de la fonction inverse $x \mapsto \frac{1}{x}$ dont la courbe « monte » indéfiniment au voisinage de $0$ par la droite.

\textbf{2.c) Limite à gauche en $0$ ?}
La fonction $f$ n'est pas définie en $0$ (division par zéro impossible). De plus, son ensemble de définition est $]0\,;\,10]$ : il ne contient \textbf{aucun} nombre strictement inférieur à $0$. La notion de limite à gauche en $0$ (qui nécessiterait des valeurs $x<0$ dans le domaine) \textbf{n'a pas de sens} pour $f$. Seule la limite à droite est pertinente.

\textbf{3.a) Tracé de la courbe $\mathcal{C}_f$ de $f$ sur $]0\,;\,10]$.}
On utilise les points calculés et le tableau ci-dessous :
\begin{center}
\begin{tabular}{|c|ccccc|}\hline
\rowcolor{popblueL} $x$ & $0{,}5$ & $1$ & $2$ & $5$ & $10$ \\ \hline
$f(x)$ & $250$ & $150$ & $100$ & $70$ & $60$ \\ \hline
\end{tabular}
\end{center}

\begin{popfigure}
\begin{tikzpicture}[scale=0.7, y=0.02cm, >=Stealth]
\def\xmin{-0.5} \def\xmax{10.5} \def\ymin{0} \def\ymax{300}
\draw[popline, thin] (\xmin,0) -- (\xmax,0) node[right] {$x$ (min)};
\draw[popline, thin] (0,\ymin) -- (0,\ymax) node[above] {$f(x)$ (FCFA)};
\foreach \x in {1,2,...,10} \draw (\x,2pt) -- (\x,-2pt) node[below, font=\tiny] {\x};
\foreach \y in {50,100,...,300} \draw (2pt,\y) -- (-2pt,\y) node[left, font=\tiny] {\y};
% Asymptote verticale x=0
\draw[popmuted, dashed] (0,\ymin) -- (0,\ymax) node[above, font=\tiny, popmuted] {Asymptote $x=0$};
% Courbe
\draw[popblue, thick, domain=0.4:10, samples=300] plot (\x, {50 + 100/\x});
% Points remarquables
\foreach \x/\y in {0.5/250, 1/150, 2/100, 5/70, 10/60} {
    \fill[popblue] (\x,\y) circle (2pt);
}
\node[popblue, above right] at (10,60) {$f(x)=50+\frac{100}{x}$};
\node[poporange, anchor=south west] at (1,250) {Croissance vers $+\infty$ quand $x\to 0^+$};
\end{tikzpicture}
\legende{Courbe représentative de $f$ (coût d'appel) sur $]0\,;\,10]$.}
\end{popfigure}

\textbf{3.b) Continuité en $0$.}
La courbe $\mathcal{C}_f$ possède une branche infinie (asymptote verticale l'axe des ordonnées) au voisinage de $0$. Il est impossible de prolonger le tracé jusqu'à l'abscisse $0$ « sans lever le crayon » : la courbe s'envole vers $+\infty$.
De plus, $f$ n'est pas définie en $0$. La condition « $f$ définie en $a$ » de la définition de la continuité n'est pas satisfaite. Par conséquent, $f$ \textbf{n'est pas continue en $0$}.

\textbf{3.c) Continuité sur $]0\,;\,10]$.}
Sur l'intervalle $]0\,;\,10]$, la fonction inverse $x \mapsto \frac{1}{x}$ est continue (c'est une fonction continue de référence sur $]0\,;\,+\infty[$).
La fonction $f$ s'écrit $f(x) = 50 + 100 \times \frac{1}{x}$. C'est une somme de la fonction constante $x \mapsto 50$ (continue sur $\mathbb{R}$) et du produit de la constante $100$ par la fonction inverse (continue sur $]0\,;\,10]$).
D'après les opérations sur les fonctions continues, $f$ est \cle{continue} sur l'intervalle $]0\,;\,10]$. Graphiquement, sa courbe se trace d'un seul trait sur cet intervalle, sans interruption.

\textbf{4) Exemple d'avis au gérant.}
« Monsieur le Gérant,
Le coût d'un appel international est modélisé par $f(x) = 50 + \frac{100}{x}$ pour une durée $x$ en minutes. Lorsque la durée $x$ se rapproche de zéro minute, le terme $\frac{100}{x}$ devient arbitrairement grand : une communication de $0{,}01$ minute (soit $0{,}6$ seconde) coûterait déjà \num{10050} FCFA. Mathématiquement, la limite de $f$ à droite en $0$ est $+\infty$, et la fonction n'est pas continue en $0$ (elle n'est même pas définie).
Nous vous recommandons vivement de fixer une durée minimale facturée, par exemple $1$ minute (coût \num{150} FCFA), voire $30$ secondes (coût \num{250} FCFA). Cette mesure commerciale simple évite des facturations aberrantes, protège les clients contre des frais disproportionnés pour des appels avortés, et assure une tarification continue et prévisible sur l'intervalle des durées réellement facturées. »
\end{exoresolu}

\courssub{Pièges à éviter}

\begin{attention}[Erreurs fréquentes au Probatoire]
\begin{itemize}
\item \textbf{Division par zéro interdite :} Ne jamais écrire $f(0) = 50 + \frac{100}{0}$ ni $\frac{100}{0} = +\infty$. La notation $\frac{\ell}{0}$ ($\ell \in \mathbb{R}^*$) est \textbf{proscrite} par le programme officiel. On écrit : « quand $x \to 0^+$, $\frac{100}{x} \to +\infty$, donc $f(x) \to +\infty$ ».
\item \textbf{Parité : deux conditions obligatoires.} Oublier de vérifier la symétrie du domaine ($x \in D \implies -x \in D$) est une erreur de rédaction majeure. Il faut l'écrire explicitement : « $D_f$ est symétrique car... ».
\item \textbf{Limite à gauche inexistante :} Ne pas confondre « la limite à gauche n'existe pas » (comportement oscillant ou infini) et « la limite à gauche n'a pas de sens » (domaine ne contenant pas de valeurs $< x_0$). Ici, pour $f$ en $0$, c'est le second cas.
\item \textbf{Continuité en un point :} Trois conditions : (1) $f$ définie en $a$, (2) $\lim_{x\to a} f(x)$ existe et est finie, (3) égalité des deux. L'absence de l'une suffit pour la non-continuité.
\item \textbf{Unités et arrondis :} Respecter les unités du problème (FCFA, minutes, heures) et les consignes d'arrondi (ici « au FCFA près »).
\end{itemize}
\end{attention}

\courssub{Bilan de l'activité}

Cette activité t'a permis de mobiliser, dans un contexte réaliste de tarification téléphonique, l'ensemble des savoir-faire de la leçon :
\begin{itemize}
\item la justification rigoureuse de la \cle{parité} d'une fonction sur un domaine borné symétrique, avec interprétation concrète de la symétrie de la courbe (recette matin/soir) ;
\item la \cle{conjecture d'une limite} à droite en un réel à l'aide d'une table de valeurs, en lien avec la courbe de la fonction inverse ;
\item la distinction entre \cle{limite à droite} et \cle{limite à gauche}, et la reconnaissance des situations où l'une des deux n'a pas de sens (domaine non étendu de ce côté) ;
\item la justification de la \cle{continuité} sur un intervalle par identification à des fonctions continues du programme (somme, produit, inverse), et de la non-continuité en un point par lecture graphique (« lever le crayon ») et analyse algébrique (non-définition, limite infinie) ;
\item enfin, la \cle{communication mathématique} : transformer des résultats techniques (limites, continuité) en un avis compréhensible, argumenté et utile pour un décideur non mathématicien.
\end{itemize}

\begin{savaistu}[Pourquoi les opérateurs facturent-ils une durée minimale ?]
Les opérateurs camerounais comme MTN, Orange ou CAMTEL facturent souvent la première minute indivisible (ou par paliers de 30 secondes) pour les appels internationaux. Ce n'est pas un hasard commercial : comme tu viens de le découvrir, les modèles de coût intègrent des frais fixes et des termes en $1/x$ qui « explosent » près de zéro. Fixer un seuil minimal $x_{\min} > 0$ revient mathématiquement à restreindre l'ensemble de définition à un intervalle $[x_{\min}\,;\,+\infty[$ sur lequel la fonction de coût est continue, strictement décroissante et à valeurs raisonnables. Les mathématiques protègent ainsi à la fois l'opérateur (modèle viable) et le client (facturation prévisible et équitable) !
\end{savaistu}

\newpage

\coursec{Méthodes et exercices résolus}

\coursec{Fonctions : généralités, limites et continuité}

\courssub{Étudier la parité d'une fonction}

\begin{definition}[Fonction paire et fonction impaire]
Soit $f$ une fonction définie sur un ensemble $D \subset \mathbb{R}$.
\begin{itemize}
\item $f$ est \cle{paire} sur $D$ si $D$ est symétrique par rapport à $0$ (c'est-à-dire $\forall x \in D,\ -x \in D$) et si $\forall x \in D,\ f(-x) = f(x)$.
\item $f$ est \cle{impaire} sur $D$ si $D$ est symétrique par rapport à $0$ et si $\forall x \in D,\ f(-x) = -f(x)$.
\end{itemize}
\end{definition}

\begin{propriete}[Conséquences graphiques]
\begin{itemize}
\item Si $f$ est paire, sa courbe représentative $\mathscr{C}_f$ admet l'axe des ordonnées $(Oy)$ comme \textbf{axe de symétrie}.
\item Si $f$ est impaire, sa courbe représentative $\mathscr{C}_f$ admet l'origine $O$ du repère comme \textbf{centre de symétrie}.
\end{itemize}
\end{propriete}

\begin{methode}[Justifier qu'une fonction est paire ou impaire]
Pour étudier la parité d'une fonction $f$ définie sur un domaine $D$ (souvent un intervalle borné $[-a\,;\,a]$ ou $[-a\,;\,a]\setminus\{0\}$) :
\begin{enumerate}
\item \textbf{Vérifier la symétrie du domaine} : montrer que pour tout $x\in D$, on a aussi $-x\in D$ (le domaine est \cle{symétrique par rapport à zéro}). Si ce n'est pas le cas, $f$ n'est ni paire ni impaire : on s'arrête là.
\item \textbf{Calculer $f(-x)$} pour un $x$ quelconque de $D$, en simplifiant soigneusement.
\item \textbf{Comparer avec $f(x)$ et $-f(x)$} :
\begin{itemize}
\item si $f(-x)=f(x)$ pour tout $x\in D$, alors $f$ est \cle{paire} ;
\item si $f(-x)=-f(x)$ pour tout $x\in D$, alors $f$ est \cle{impaire} ;
\item sinon, $f$ n'est ni paire ni impaire.
\end{itemize}
\item \textbf{Conclure graphiquement} : si $f$ est paire, sa courbe admet l'axe des ordonnées comme axe de symétrie ; si $f$ est impaire, sa courbe admet l'origine $O$ comme centre de symétrie.
\end{enumerate}
\textbf{Réflexes de calcul} : $(-x)^2=x^2$, $(-x)^3=-x^3$, $|-x|=|x|$, $\dfrac{1}{-x}=-\dfrac{1}{x}$.
\end{methode}

\begin{exoresolu}[Parité sur un domaine borné]
On considère les fonctions définies sur $[-3\,;\,3]$ par :
\[ f(x)=x^2-2|x|+1 \qquad \text{et} \qquad g(x)=\dfrac{x^3-4x}{2}. \]
\begin{enumerate}
\item Justifier que le domaine $[-3\,;\,3]$ est symétrique par rapport à zéro.
\item Étudier la parité de $f$, puis celle de $g$.
\item En déduire les éléments de symétrie de leurs courbes représentatives.
\end{enumerate}
\tcblower
\textbf{1. Symétrie du domaine.} Soit $x\in[-3\,;\,3]$. Alors $-3\le x\le 3$, donc $-3\le -x\le 3$, c'est-à-dire $-x\in[-3\,;\,3]$. Le domaine est bien symétrique par rapport à zéro : l'étude de parité a un sens.

\textbf{2. Parité de $f$.} Pour tout $x\in[-3\,;\,3]$ :
\begin{align*}
f(-x) &= (-x)^2-2|-x|+1 \\
&= x^2-2|x|+1 \quad \text{car } (-x)^2=x^2 \text{ et } |-x|=|x|\\
&= f(x).
\end{align*}
Donc pour tout $x\in[-3\,;\,3]$, $-x\in[-3\,;\,3]$ et $f(-x)=f(x)$ : $f$ est \textbf{paire}.

\textbf{Parité de $g$ :} pour tout $x\in[-3\,;\,3]$ :
\begin{align*}
g(-x) &= \dfrac{(-x)^3-4(-x)}{2} = \dfrac{-x^3+4x}{2} \\
&= -\dfrac{x^3-4x}{2} = -g(x).
\end{align*}
Donc pour tout $x\in[-3\,;\,3]$, $g(-x)=-g(x)$ : $g$ est \textbf{impaire}.

\textbf{3. Éléments de symétrie.} La courbe de $f$ admet l'axe des ordonnées comme \textbf{axe de symétrie} ; la courbe de $g$ admet l'origine $O$ du repère comme \textbf{centre de symétrie}. Concrètement, il suffit de tracer chaque courbe sur $[0\,;\,3]$, puis de compléter par symétrie.
\end{exoresolu}

\begin{methode}[Conjecturer la parité par lecture graphique]
À partir de la courbe représentative d'une fonction définie sur un domaine borné :
\begin{enumerate}
\item Observer si le domaine représenté est symétrique par rapport à zéro (sinon, pas de parité possible).
\item Tester visuellement une \textbf{symétrie par rapport à l'axe des ordonnées} : si plier la feuille le long de cet axe superpose les deux moitiés de la courbe, on conjecture que $f$ est \cle{paire}.
\item Tester visuellement une \textbf{symétrie par rapport à l'origine} : si un demi-tour autour de $O$ laisse la courbe invariante, on conjecture que $f$ est \cle{impaire}.
\item Vérifier sur des points : relever $f(1)$ et $f(-1)$, $f(2)$ et $f(-2)$. Si $f(-a)=f(a)$, indice de parité ; si $f(-a)=-f(a)$, indice d'imparité ; un seul contre-exemple suffit à rejeter les deux.
\end{enumerate}
\textbf{Attention} : une conjecture graphique n'est pas une preuve ; la justification exige le calcul de $f(-x)$.
\end{methode}

\begin{exoresolu}[Lecture graphique puis justification]
La courbe ci-dessous représente une fonction $h$ définie sur $[-2\,;\,2]$.
\begin{popfigure}
\begin{tikzpicture}[scale=1.1]
\draw[popline] (-2.6,0) -- (2.6,0) node[below left]{\small $x$};
\draw[popline] (0,-0.6) -- (0,4.4) node[below left]{\small $y$};
\draw[popline] (0,0) node[below left]{\small $O$};
\foreach \x in {-2,-1,1,2} \draw[popline] (\x,0.08) -- (\x,-0.08) node[below]{\small $\x$};
\foreach \y in {1,2,3,4} \draw[popline] (0.08,\y) -- (-0.08,\y) node[left]{\small $\y$};
\draw[popblue,very thick,domain=-2:2,samples=100] plot (\x,{\x*\x});
\draw[popink] (1.5,2.25) node[right]{\small $\mathscr{C}_h$};
\end{tikzpicture}
\legende{Courbe de la fonction $h(x)=x^2$ sur $[-2;2]$.}
\end{popfigure}
\begin{enumerate}
\item Conjecturer la parité de $h$ par lecture graphique.
\item Sachant que $h(x)=x^2$ sur $[-2\,;\,2]$, justifier cette conjecture.
\end{enumerate}
\tcblower
\textbf{1. Conjecture.} La courbe est tracée sur un domaine symétrique par rapport à zéro, et elle semble invariante par symétrie par rapport à l'axe des ordonnées : les points d'abscisses $1$ et $-1$ ont la même image ($1$), de même pour $2$ et $-2$ (image $4$). On conjecture que $h$ est \textbf{paire}.

\textbf{2. Justification.} Pour tout $x\in[-2\,;\,2]$, on a $-x\in[-2\,;\,2]$ (domaine symétrique) et :
\[ h(-x)=(-x)^2=x^2=h(x). \]
Donc $h$ est bien \textbf{paire}, ce qui confirme la conjecture : l'axe des ordonnées est axe de symétrie de $\mathscr{C}_h$.
\end{exoresolu}

\courssub{Approche intuitive de la limite en un réel}

\begin{methode}[Conjecturer une limite par table de valeurs ou graphiquement]
Pour conjecturer la limite de $f$ en un réel $x_0$ (où $f$ n'est pas forcément définie) :
\begin{enumerate}
\item \textbf{Table de valeurs} : calculer $f(x)$ pour des valeurs de $x$ de plus en plus proches de $x_0$, par valeurs inférieures ($x_0-0{,}1$ ; $x_0-0{,}01$ ; \ldots) et par valeurs supérieures ($x_0+0{,}1$ ; $x_0+0{,}01$ ; \ldots).
\item Observer vers quel nombre $\ell$ les images se rapprochent : on conjecture que $\lim\limits_{x\to x_0}f(x)=\ell$.
\item \textbf{Lecture graphique} : suivre la courbe des deux côtés de $x_0$ et lire l'ordonnée vers laquelle elle semble tendre.
\end{enumerate}
\textbf{Réflexe} : si $f$ est une fonction « usuelle » définie en $x_0$ (polynôme, quotient défini en $x_0$, racine carrée, valeur absolue), la limite en $x_0$ est simplement $f(x_0)$.
\end{methode}

\begin{exoresolu}[Limite conjecturée par une table de valeurs]
Soit $f$ la fonction définie sur $\mathbb{R}\setminus\{2\}$ par $f(x)=\dfrac{x^2-4}{x-2}$.
\begin{enumerate}
\item Compléter la table de valeurs suivante (arrondir au centième) :
\begin{center}
\begin{tabular}{|c|cccccc|}\hline
$x$ & $1{,}9$ & $1{,}99$ & $1{,}999$ & $2{,}001$ & $2{,}01$ & $2{,}1$ \\ \hline
$f(x)$ & & & & & & \\ \hline
\end{tabular}
\end{center}
\item Conjecturer la limite de $f$ en $2$.
\item Justifier cette conjecture par le calcul.
\end{enumerate}
\tcblower
\textbf{1. Table de valeurs.} En calculant chaque image :
\begin{center}
\begin{tabular}{|c|cccccc|}\hline
$x$ & $1{,}9$ & $1{,}99$ & $1{,}999$ & $2{,}001$ & $2{,}01$ & $2{,}1$ \\ \hline
$f(x)$ & $3{,}9$ & $3{,}99$ & $3{,}999$ & $4{,}001$ & $4{,}01$ & $4{,}1$ \\ \hline
\end{tabular}
\end{center}
Par exemple $f(1{,}9)=\dfrac{1{,}9^2-4}{1{,}9-2}=\dfrac{3{,}61-4}{-0{,}1}=\dfrac{-0{,}39}{-0{,}1}=3{,}9$.

\textbf{2. Conjecture.} Quand $x$ se rapproche de $2$ (par valeurs inférieures ou supérieures), $f(x)$ se rapproche de $4$. On conjecture :
\[ \lim_{x\to 2}f(x)=4. \]

\textbf{3. Justification.} Pour $x\neq 2$, on factorise le numérateur (identité remarquable) :
\[ f(x)=\dfrac{x^2-4}{x-2}=\dfrac{(x-2)(x+2)}{x-2}=x+2. \]
Ainsi, pour $x\neq 2$, $f(x)=x+2$. Quand $x$ tend vers $2$, $x+2$ tend vers $4$. Donc :
\[ \lim_{x\to 2}f(x)=\lim_{x\to 2}(x+2)=4. \]
La conjecture est confirmée. Remarque : $f$ n'est pas définie en $2$, mais sa limite en $2$ existe et vaut $4$.
\end{exoresolu}

\courssub{Limites à droite et à gauche}

\begin{methode}[Déterminer une limite à droite et à gauche en $x_0$]
Pour une fonction du type $x\mapsto \dfrac{ax+b}{cx+d}$ au voisinage de $x_0=-\dfrac{d}{c}$ (valeur interdite) :
\begin{enumerate}
\item \textbf{Identifier la valeur interdite} $x_0=-\dfrac{d}{c}$ (celle qui annule le dénominateur).
\item \textbf{Étudier le signe du dénominateur} $cx+d$ à gauche et à droite de $x_0$ (tableau de signes).
\item \textbf{Exploiter la fonction inverse de référence} : quand $x\to 0$ avec $x>0$, $\dfrac{1}{x}$ devient très grand ($+\infty$) ; quand $x\to 0$ avec $x<0$, $\dfrac{1}{x}$ devient très grand en valeur absolue mais négatif ($-\infty$).
\item \textbf{Conclure} : le numérateur tend vers un nombre non nul, le dénominateur tend vers $0$ en gardant un signe constant de chaque côté ; le quotient tend donc vers $+\infty$ ou $-\infty$ selon la règle des signes.
\item Comparer les deux limites : si elles diffèrent, la fonction \textbf{n'a pas de limite} en $x_0$ (la droite $x=x_0$ est asymptote verticale à la courbe).
\end{enumerate}
\textbf{Rappel du programme} : la notation $\dfrac{\ell}{0}$ ($\ell\in\mathbb{R}^*$) est \textbf{proscrite} ; on rédige avec le signe du dénominateur.
\end{methode}

\begin{exoresolu}[Limites à gauche et à droite d'une fonction homographique]
Soit $g$ la fonction définie sur $\mathbb{R}\setminus\{1\}$ par $g(x)=\dfrac{2x+1}{x-1}$.
\begin{enumerate}
\item Étudier le signe de $x-1$ sur $\mathbb{R}$.
\item Déterminer la limite de $g$ à gauche en $1$, puis à droite en $1$.
\item $g$ admet-elle une limite en $1$ ? Interpréter graphiquement.
\end{enumerate}
\tcblower
\textbf{1. Signe du dénominateur.} $x-1>0 \iff x>1$ et $x-1<0 \iff x<1$.
\begin{center}
\begin{tabular}{|c|ccc|}\hline
$x$ & $-\infty$ & $1$ & $+\infty$ \\ \hline
$x-1$ & $-$ & $0$ & $+$ \\ \hline
\end{tabular}
\end{center}

\textbf{2. Limites.} Quand $x$ tend vers $1$, le numérateur $2x+1$ tend vers $2\times 1+1=3>0$.

\textbf{À gauche en $1$} ($x<1$) : le dénominateur $x-1$ tend vers $0$ en restant \textbf{négatif}. Le quotient $\dfrac{\text{proche de }3}{\text{proche de }0^-}$ devient très grand en valeur absolue, de signe négatif :
\[ \lim_{\substack{x\to 1\\ x<1}} g(x)=-\infty. \]

\textbf{À droite en $1$} ($x>1$) : le dénominateur $x-1$ tend vers $0$ en restant \textbf{positif}. Le quotient devient très grand, de signe positif :
\[ \lim_{\substack{x\to 1\\ x>1}} g(x)=+\infty. \]

\textbf{3. Conclusion.} Les limites à gauche et à droite sont différentes ($-\infty\neq +\infty$), donc $g$ \textbf{n'admet pas de limite en $1$}. Graphiquement, la droite d'équation $x=1$ est \textbf{asymptote verticale} à la courbe de $g$ : la courbe « descend » vers $-\infty$ à gauche de $1$ et « monte » vers $+\infty$ à droite de $1$.
\end{exoresolu}

\courssub{Continuité en un réel et sur un intervalle}

\begin{definition}[Continuité en un point]
Soit $f$ une fonction définie sur un intervalle $I$ et $x_0 \in I$. On dit que $f$ est \cle{continue en $x_0$} si et seulement si :
\[ \lim_{x\to x_0}f(x)=f(x_0). \]
Cela suppose trois conditions :
\begin{enumerate}
\item $f$ est définie en $x_0$ ;
\item $f$ admet une limite finie $\ell$ en $x_0$ (les limites à gauche et à droite existent et sont égales) ;
\item cette limite est égale à l'image : $\ell = f(x_0)$.
\end{enumerate}
\end{definition}

\begin{propriete}[Fonctions continues de référence (programme de Première)]
Les fonctions suivantes sont continues sur tout intervalle inclus dans leur ensemble de définition :
\begin{itemize}
\item les fonctions polynômes (sur $\mathbb{R}$) ;
\item la fonction valeur absolue $x \mapsto |x|$ (sur $\mathbb{R}$) ;
\item les fonctions rationnelles (sur tout intervalle ne contenant pas de valeur interdite) ;
\item la fonction racine carrée $x \mapsto \sqrt{x}$ (sur $[0\,;\,+\infty[$).
\end{itemize}
On peut justifier la continuité d'une fonction \textbf{en l'identifiant} à l'une de ces fonctions sur l'intervalle considéré.
\end{propriete}

\begin{methode}[Justifier la continuité en un réel]
Pour justifier qu'une fonction $f$ est \cle{continue en} $x_0$ :
\begin{enumerate}
\item Vérifier que $f$ est \textbf{définie en $x_0$} (sinon, $f$ ne peut pas être continue en $x_0$).
\item Déterminer la limite de $f$ en $x_0$ (éventuellement les limites à gauche et à droite, qui doivent être égales).
\item Comparer avec $f(x_0)$ : $f$ est continue en $x_0$ si et seulement si
\[ \lim_{x\to x_0}f(x)=f(x_0). \]
\end{enumerate}
\end{methode}

\begin{exoresolu}[Continuité en un réel]
Soit $f$ la fonction définie sur $\mathbb{R}$ par :
\[ f(x)=\begin{cases} \dfrac{x^2-4}{x-2} & \text{si } x\neq 2 \\ 4 & \text{si } x=2 \end{cases} \]
La fonction $f$ est-elle continue en $2$ ?
\tcblower
\textbf{Étape 1 :} $f$ est définie en $2$ et $f(2)=4$.

\textbf{Étape 2 :} déterminons la limite de $f$ en $2$. Pour $x\neq 2$ :
\[ f(x)=\dfrac{x^2-4}{x-2}=\dfrac{(x-2)(x+2)}{x-2}=x+2. \]
Donc $\lim\limits_{x\to 2}f(x)=\lim\limits_{x\to 2}(x+2)=4$.

\textbf{Étape 3 :} on compare : $\lim\limits_{x\to 2}f(x)=4=f(2)$.

\textbf{Conclusion :} la limite de $f$ en $2$ existe et est égale à $f(2)$, donc $f$ est \textbf{continue en $2$}. On dit qu'on a « prolongé par continuité » la fonction $x\mapsto\dfrac{x^2-4}{x-2}$ en $2$.
\end{exoresolu}

\begin{methode}[Justifier la continuité sur un intervalle par lecture graphique]
Pour justifier graphiquement qu'une fonction est continue sur un intervalle $]a\,;\,b[$ :
\begin{enumerate}
\item Observer la courbe sur tout l'intervalle $]a\,;\,b[$.
\item Vérifier qu'on peut la tracer \textbf{sans lever le crayon} : aucun « saut », aucun « trou », aucune asymptote verticale entre $a$ et $b$.
\item Conclure : la courbe étant en un seul morceau sur $]a\,;\,b[$, la fonction est continue sur $]a\,;\,b[$.
\end{enumerate}
\textbf{Réflexe} : si la courbe présente une coupure en un point $x_0$ de l'intervalle (saut ou valeur interdite), la fonction n'est pas continue en $x_0$, donc pas continue sur l'intervalle.
\end{methode}

\begin{exoresolu}[Continuité sur un intervalle par lecture graphique]
On donne ci-dessous les courbes de deux fonctions $u$ et $v$ sur $]-3\,;\,3[$.
\begin{popfigure}
\begin{tikzpicture}[scale=0.85]
% Courbe de u
\begin{scope}
\draw[popline] (-3.3,0) -- (3.3,0) node[below left]{\small $x$};
\draw[popline] (0,-1.2) -- (0,3.2) node[below left]{\small $y$};
\draw[popline] (0,0) node[below left]{\small $O$};
\foreach \x in {-3,-2,-1,1,2,3} \draw[popline] (\x,0.06) -- (\x,-0.06) node[below]{\scriptsize $\x$};
\draw[popblue,very thick,domain=-2.9:2.9,samples=100] plot (\x,{0.25*\x*\x});
\draw[popink] (0,3) node{\small $\mathscr{C}_u$};
\end{scope}
% Courbe de v
\begin{scope}[xshift=9cm]
\draw[popline] (-3.3,0) -- (3.3,0) node[below left]{\small $x$};
\draw[popline] (0,-1.2) -- (0,3.2) node[below left]{\small $y$};
\draw[popline] (0,0) node[below left]{\small $O$};
\foreach \x in {-3,-2,-1,1,2,3} \draw[popline] (\x,0.06) -- (\x,-0.06) node[below]{\scriptsize $\x$};
\draw[poporange,very thick,domain=-2.9:0.95,samples=60] plot (\x,{0.5*\x+1});
\draw[poporange,very thick,domain=1.05:2.9,samples=60] plot (\x,{0.5*\x+2});
\draw[poporange,fill=white] (1,1.5) circle (0.09);
\draw[poporange,fill=poporange] (1,2.5) circle (0.09);
\draw[popink] (0,3) node{\small $\mathscr{C}_v$};
\end{scope}
\end{tikzpicture}
\legende{À gauche : courbe de $u$ (parabole). À droite : courbe de $v$ (affine par morceaux avec un saut en $x=1$).}
\end{popfigure}
Justifier, par lecture graphique, la continuité ou la non-continuité de $u$ et de $v$ sur $]-3\,;\,3[$.
\tcblower
\textbf{Fonction $u$.} Sur $]-3\,;\,3[$, la courbe $\mathscr{C}_u$ est tracée d'un seul tenant : on peut la parcourir sans lever le crayon, elle ne présente ni saut ni trou. Donc $u$ est \textbf{continue sur $]-3\,;\,3[$}.

\textbf{Fonction $v$.} La courbe $\mathscr{C}_v$ présente une \textbf{coupure en $x=1$} : à gauche de $1$, la courbe s'approche du point $(1\,;\,1{,}5)$ (point vide), tandis que $v(1)=2{,}5$ (point plein). On ne peut pas tracer la courbe sans lever le crayon au voisinage de $1$. Donc $v$ \textbf{n'est pas continue en $1$}, et par conséquent $v$ \textbf{n'est pas continue sur $]-3\,;\,3[$}. En revanche, $v$ est continue sur $]-3\,;\,1[$ et sur $[1\,;\,3[$ séparément.
\end{exoresolu}

\courssub{Synthèse : lien entre limite et continuité, méthodes de rédaction}

\begin{aretenir}[L'essentiel à retenir pour le Probatoire]
\begin{itemize}
\item \textbf{Parité} : Vérifier \textbf{d'abord} la symétrie du domaine ($x\in D \Rightarrow -x\in D$). Ensuite calculer $f(-x)$ et comparer à $f(x)$ (paire) ou $-f(x)$ (impaire). Conclure sur l'axe ou le centre de symétrie de la courbe.
\item \textbf{Limite en $x_0$} : Si $f$ est « usuelle » et définie en $x_0$, $\lim\limits_{x\to x_0}f(x)=f(x_0)$. Sinon, simplifier l'expression (factorisation, conjugaison) pour lever l'indétermination, ou utiliser une table de valeurs / lecture graphique pour conjecturer.
\item \textbf{Limites à droite/gauche} : Pour $x_0$ valeur interdite d'une fonction homographique $\frac{ax+b}{cx+d}$, étudier le signe du dénominateur de part et d'autre de $x_0$. Rédiger : « numérateur tend vers $\ell \neq 0$, dénominateur tend vers $0$ en restant $>0$ (ou $<0$), donc quotient tend vers $+\infty$ (ou $-\infty$) ». \textbf{Jamais} écrire $\frac{\ell}{0}$.
\item \textbf{Continuité en $x_0$} : Trois conditions obligatoires : (1) $f(x_0)$ existe, (2) $\lim\limits_{x\to x_0}f(x)$ existe (finie), (3) égalité des deux. Pour une fonction définie par morceaux, vérifier l'égalité des limites à gauche et à droite en le point de jonction.
\item \textbf{Continuité sur un intervalle} : Fonction continue en tout point de l'intervalle. Graphiquement : courbe tracée sans lever le crayon.
\end{itemize}
\end{aretenir}

\begin{methode}[Rédiger une justification de continuité (modèle Probatoire)]
\textbf{Exemple type :} « Montrer que $f$ est continue sur $[-2\,;\,3]$. »
\begin{enumerate}
\item \textbf{Sur $[-2\,;\,1[ \cup ]1\,;\,3]$} : $f$ coïncide avec la fonction [polynôme / rationnelle / racine / valeur absolue] qui est continue sur cet intervalle (fonction de référence). Donc $f$ est continue sur $[-2\,;\,1[ \cup ]1\,;\,3]$.
\item \textbf{En $1$ (point de jonction ou valeur suspecte)} :
\begin{itemize}
\item $f(1) = \dots$ (calculer l'image).
\item $\lim\limits_{x\to 1^-}f(x) = \dots$ (calculer limite à gauche).
\item $\lim\limits_{x\to 1^+}f(x) = \dots$ (calculer limite à droite).
\item Comme les deux limites sont égales à $f(1)$, $f$ est continue en $1$.
\end{itemize}
\item \textbf{Conclusion} : $f$ est continue sur $[-2\,;\,3]$.
\end{enumerate}
\end{methode}

\begin{attention}[Les 5 erreurs qui coûtent des points au Probatoire]
\begin{enumerate}
\item \textbf{Oublier de vérifier la symétrie du domaine avant la parité.} Écrire « $f(-x)=f(x)$ donc $f$ est paire » sans avoir justifié que pour tout $x\in D$, $-x\in D$, est une démonstration incomplète : la condition sur le domaine est exigée et rapporte un point.
\item \textbf{Conclure sur la parité avec un seul exemple numérique.} Vérifier que $f(-2)=f(2)$ ne prouve rien : il faut le calcul littéral de $f(-x)$ pour un $x$ \emph{quelconque}. Un exemple numérique ne sert qu'à \emph{infirmer} une conjecture.
\item \textbf{Écrire la notation proscrite $\dfrac{\ell}{0}$.} Le programme l'interdit explicitement. Rédige : « quand $x$ tend vers $1$ avec $x<1$, le numérateur tend vers $3$ et le dénominateur tend vers $0$ en restant négatif, donc le quotient tend vers $-\infty$ ».
\item \textbf{Confondre « limite en $x_0$ » et « image $f(x_0)$ ».} Une fonction peut avoir une limite en $x_0$ sans y être définie (cas de $\dfrac{x^2-4}{x-2}$ en $2$). Ne jamais écrire $f(2)$ quand $2$ est une valeur interdite : calculer la limite par simplification.
\item \textbf{Déclarer une fonction continue en $x_0$ sans comparer limite et image.} La continuité exige l'égalité $\lim\limits_{x\to x_0}f(x)=f(x_0)$ : il faut vérifier les \emph{trois} conditions (définie en $x_0$, limite existante — à gauche et à droite égales —, égalité avec l'image). Oublier l'une d'elles rend la justification fausse.
\end{enumerate}
\end{attention}

\newpage

\coursec{Exercices}

\coursec{Fonctions : généralités, limites et continuité}

\courssub{Parité d'une fonction}

\begin{definition}[Domaine symétrique]
Soit $D$ une partie de $\mathbb{R}$. On dit que $D$ est \cle{symétrique par rapport à $0$} (ou \cle{symétrique}) si pour tout $x\in D$, son opposé $-x$ appartient aussi à $D$.
\end{definition}

\begin{exemplebox}[Domaines symétriques ou non]
\begin{itemize}
\item $[-3\,;\,3]$, $]-2\,;\,2[$, $\mathbb{R}$, $\mathbb{R}^*$ sont symétriques.
\item $[0\,;\,5]$, $[-1\,;\,2]$, $]-\infty\,;\,3]$ ne le sont pas (ex~: $2\in[-1\,;\,2]$ mais $-2\notin[-1\,;\,2]$).
\end{itemize}
\end{exemplebox}

\begin{definition}[Fonction paire, fonction impaire]
Soit $f$ une fonction définie sur un domaine $D$ \textbf{symétrique}.
\begin{itemize}
\item $f$ est \cle{paire} si pour tout $x\in D$, $f(-x)=f(x)$.
\item $f$ est \cle{impaire} si pour tout $x\in D$, $f(-x)=-f(x)$.
\end{itemize}
Si $D$ n'est pas symétrique, $f$ n'est \textbf{ni paire ni impaire}.
\end{definition}

\begin{propriete}[Interprétation graphique]
Soit $f$ une fonction définie sur un domaine symétrique $D$ et $\mathcal{C}_f$ sa courbe représentative dans un repère orthonormé.
\begin{itemize}
\item $f$ est paire $\iff$ $\mathcal{C}_f$ est symétrique par rapport à l'axe des ordonnées (axe $Oy$).
\item $f$ est impaire $\iff$ $\mathcal{C}_f$ est symétrique par rapport à l'origine $O$ (rotation d'angle $\pi$).
\end{itemize}
\end{propriete}

\begin{methode}[Étudier la parité d'une fonction]
\begin{enumerate}
\item Vérifier que le domaine de définition $D_f$ est symétrique par rapport à $0$. Sinon, conclure : « $f$ n'est ni paire ni impaire car son domaine n'est pas symétrique ».
\item Pour $x\in D_f$, calculer $f(-x)$.
\item Comparer $f(-x)$ à $f(x)$ :
\begin{itemize}
\item si $f(-x)=f(x)$, $f$ est paire ;
\item si $f(-x)=-f(x)$, $f$ est impaire ;
\item sinon, $f$ n'est ni paire ni impaire.
\end{itemize}
\end{enumerate}
\end{methode}

\begin{exemplebox}[Exemples classiques]
\begin{itemize}
\item $x\mapsto x^2$ sur $\mathbb{R}$ : paire.
\item $x\mapsto x^3$ sur $\mathbb{R}$ : impaire.
\item $x\mapsto \dfrac{1}{x}$ sur $\mathbb{R}^*$ : impaire.
\item $x\mapsto x^2+1$ sur $[-2\,;\,2]$ : paire.
\item $x\mapsto x^3-3x$ sur $[-2\,;\,2]$ : impaire.
\item $x\mapsto x^2$ sur $[-1\,;\,2]$ : \textbf{ni paire ni impaire} (domaine non symétrique).
\end{itemize}
\end{exemplebox}

\begin{attention}[Piège fréquent : le domaine]
Au Probatoire, l'étude de la parité \textbf{commence toujours par la vérification de la symétrie du domaine}. Une fonction définie sur $[-1\,;\,2]$ ne peut \textbf{jamais} être paire ni impaire, quelle que soit sa formule. Oublier cette étape fait perdre la majorité des points.
\end{attention}

\courssub{Approche intuitive de la limite en un réel}

\begin{experience}[Approche numérique : $f(x)=\dfrac{x^2-4}{x-2}$ près de $2$]
La fonction $f$ n'est pas définie en $2$. Calculons des images pour des valeurs proches de $2$ :
\begin{center}
\begin{tabular}{|c|cccccc|}\hline
$x$ & $1,9$ & $1,99$ & $1,999$ & $2,001$ & $2,01$ & $2,1$ \\ \hline
$f(x)$ & $3,9$ & $3,99$ & $3,999$ & $4,001$ & $4,01$ & $4,1$ \\ \hline
\end{tabular}
\end{center}
On observe que $f(x)$ se rapproche de $4$ quand $x$ se rapproche de $2$. On \cle{conjecture} $\lim_{x\to 2}f(x)=4$.
\end{experience}

\begin{definition}[Limite finie en un réel]
Soit $f$ une fonction définie sur un intervalle (sauf peut-être en $x_0$) et $L\in\mathbb{R}$.
On dit que $f$ admet la \cle{limite $L$ en $x_0$} et on note $\displaystyle\lim_{x\to x_0}f(x)=L$ si, lorsque $x$ se rapproche de $x_0$ (sans être égal à $x_0$), $f(x)$ se rapproche de $L$ de façon arbitrairement précise.
\end{definition}

\begin{propriete}[Limites usuelles et opérations (rappels)]
Soient $f,g$ admettant des limites finies en $x_0$.
\begin{itemize}
\item $\lim_{x\to x_0}(f+g)(x)=\lim f + \lim g$
\item $\lim_{x\to x_0}(f\times g)(x)=\lim f \times \lim g$
\item $\lim_{x\to x_0}\dfrac{f}{g}(x)=\dfrac{\lim f}{\lim g}$ si $\lim g\neq 0$
\item Polynôme : $\lim_{x\to x_0}P(x)=P(x_0)$
\item Fonction rationnelle : si $Q(x_0)\neq 0$, $\lim\frac{P}{Q}=\frac{P(x_0)}{Q(x_0)}$ ; si $Q(x_0)=0$, factoriser et simplifier.
\end{itemize}
\end{propriete}

\begin{definition}[Limites infinies]
Si $f(x)$ devient arbitrairement grand (positif) quand $x\to x_0$, on note $\lim_{x\to x_0}f(x)=+\infty$.
Si $f(x)$ devient arbitrairement grand en valeur absolue mais négatif, on note $\lim_{x\to x_0}f(x)=-\infty$.
\textbf{Attention :} Une limite infinie n'est \textbf{pas} une limite (au sens fini). On dit « $f$ n'admet pas de limite finie » ou « la limite est infinie ».
\end{definition}

\begin{aretenir}[Notation interdite]
La notation $l_0$ (pour $l\in\mathbb{R}^*$) est \textbf{proscrite} au Probatoire. On écrit $\lim_{x\to x_0}f(x)=+\infty$ ou $-\infty$.
\end{aretenir}

\begin{exemplebox}[Calculs de limites]
\begin{align*}
\lim_{x\to 1}(2x^2-5x+3) &= 2-5+3 = 0. \\
\lim_{x\to 2}\frac{x^2-4}{x-2} &= \lim_{x\to 2}\frac{(x-2)(x+2)}{x-2} = \lim_{x\to 2}(x+2) = 4. \\
\lim_{x\to 0^+}\frac{1}{x} &= +\infty,\quad \lim_{x\to 0^-}\frac{1}{x} = -\infty.
\end{align*}
\end{exemplebox}

\courssub{Limites à droite et à gauche}

\begin{definition}[Limite à droite et à gauche]
Soit $f$ définie autour de $x_0$ (sauf peut-être en $x_0$).
\begin{itemize}
\item \cle{Limite à droite} en $x_0$ : $x$ tend vers $x_0$ par valeurs \textbf{supérieures} à $x_0$. Notation : $\lim_{x\to x_0^+}f(x)$.
\item \cle{Limite à gauche} en $x_0$ : $x$ tend vers $x_0$ par valeurs \textbf{inférieures} à $x_0$. Notation : $\lim_{x\to x_0^-}f(x)$.
\end{itemize}
\end{definition}

\begin{propriete}[Théorème de la limite bilatère]
$f$ admet une limite \textbf{finie} $L$ en $x_0$ \textbf{si et seulement si} les limites à droite et à gauche en $x_0$ existent, sont finies et \textbf{égales} à $L$.
\[ \lim_{x\to x_0}f(x)=L \iff \lim_{x\to x_0^-}f(x)=\lim_{x\to x_0^+}f(x)=L. \]
\end{propriete}

\begin{methode}[Limite d'une fonction homographique $\frac{ax+b}{cx+d}$ en $x_0=-\frac{d}{c}$]
\begin{enumerate}
\item Calculer le numérateur en $x_0$ : $N_0 = a x_0 + b$. (Si $N_0=0$, forme indéterminée, simplifier d'abord).
\item Étudier le signe du dénominateur $cx+d$ :
\begin{itemize}
\item pour $x\to x_0^-$ (valeurs $<x_0$) ;
\item pour $x\to x_0^+$ (valeurs $>x_0$).
\end{itemize}
\item Conclure par la règle des signes : $\frac{N_0}{0^+}=\operatorname{signe}(N_0)\times(+\infty)$, $\frac{N_0}{0^-}=\operatorname{signe}(N_0)\times(-\infty)$.
\end{enumerate}
\end{methode}

\begin{exemplebox}[Application : $f(x)=\frac{3x-1}{x+2}$ en $-2$]
$N_0 = 3(-2)-1 = -7 < 0$.
\begin{itemize}
\item $x\to -2^- \implies x<-2 \implies x+2<0 \implies \frac{-7}{0^-}=+\infty$.
\item $x\to -2^+ \implies x>-2 \implies x+2>0 \implies \frac{-7}{0^+}=-\infty$.
\end{itemize}
Limites à gauche et à droite différentes $\implies$ pas de limite en $-2$.
\end{exemplebox}

\courssub{Continuité en un réel et sur un intervalle}

\begin{definition}[Continuité en un point]
Soit $f$ définie en $x_0$. $f$ est \cle{continue en $x_0$} si :
\[ \lim_{x\to x_0}f(x) = f(x_0). \]
Cela suppose trois conditions :
\begin{enumerate}
\item $f$ est définie en $x_0$ ($x_0\in D_f$) ;
\item $f$ admet une limite \textbf{finie} en $x_0$ ;
\item Cette limite est égale à $f(x_0)$.
\end{enumerate}
Si l'une de ces conditions manque, $f$ est \cle{discontinue} en $x_0$.
\end{definition}

\begin{propriete}[Continuité sur un intervalle]
$f$ est \cle{continue sur un intervalle $I$} si elle est continue en tout point de $I$.
Graphiquement : on peut tracer la courbe sur $I$ sans lever le crayon.
\end{propriete}

\begin{propriete}[Fonctions continues sur leur domaine (référentiel)]
Les fonctions suivantes sont continues sur leur domaine de définition :
\begin{itemize}
\item Fonctions polynomiales (sur $\mathbb{R}$).
\item Fonctions rationnelles (sur $\mathbb{R}\setminus\{\text{zéros du dénominateur}\}$).
\item Fonction racine carrée $x\mapsto\sqrt{x}$ (sur $[0\,;\,+\infty[$).
\item Fonctions trigonométriques $\sin$, $\cos$ (sur $\mathbb{R}$).
\end{itemize}
Toute combinaison (somme, produit, quotient, composition) de fonctions continues est continue sur son domaine.
\end{propriete}

\begin{methode}[Étudier la continuité d'une fonction définie par morceaux en $x_0$ (frontière)]
\begin{enumerate}
\item Calculer $f(x_0)$ (valeur imposée par la définition).
\item Calculer $\lim_{x\to x_0^-}f(x)$ (expression de gauche).
\item Calculer $\lim_{x\to x_0^+}f(x)$ (expression de droite).
\item Comparer :
\begin{itemize}
\item Si les trois valeurs sont égales $\implies$ continue en $x_0$.
\item Sinon $\implies$ discontinue en $x_0$ (saut, asymptote, trou).
\end{itemize}
\end{enumerate}
\end{methode}

\courssub{Synthèse : lien entre limite et continuité, méthodes de rédaction}

\begin{aretenir}[Résumé : Limite vs Continuité]
\begin{center}
\begin{tabularx}{\linewidth}{|l|X|X|}\hline
\textbf{Situation} & \textbf{Limite en $x_0$} & \textbf{Continuité en $x_0$} \\ \hline
$f$ définie en $x_0$, $\lim f = f(x_0)$ & Existe (finie) & \textbf{Oui} \\ \hline
$f$ définie en $x_0$, $\lim f \neq f(x_0)$ & Existe (finie) & Non (trou) \\ \hline
$f$ non définie en $x_0$, $\lim f = L$ & Existe (finie) & Non (non définie) \\ \hline
$\lim_{x\to x_0^-} \neq \lim_{x\to x_0^+}$ (finis) & \textbf{Non} & Non (saut) \\ \hline
Limite infinie ($\pm\infty$) & \textbf{Non (pas de limite finie)} & Non (asymptote verticale) \\ \hline
\end{tabularx}
\end{center}
\end{aretenir}

\begin{methode}[Rédaction type au Probatoire]
\begin{enumerate}
\item \textbf{Limite :} « On a $\lim_{x\to x_0} f(x) = \lim_{x\to x_0} [\text{expression simplifiée}] = L$ » (ou $+\infty/-\infty$).
\item \textbf{Limites unilatérales :} Préciser « Pour $x\to x_0^-$, $x<x_0$ donc [signe dénominateur]... » puis « Pour $x\to x_0^+$... ».
\item \textbf{Continuité en $x_0$ :} « $f$ est définie en $x_0$ car $x_0\in D_f$. $\lim_{x\to x_0}f(x)=L$ et $f(x_0)=L$, donc $f$ est continue en $x_0$. »
\item \textbf{Continuité sur $I$ :} « $f$ est continue sur $I$ car c'est une fonction [polynomiale / rationnelle / etc.] continue sur son domaine qui contient $I$. » Ou vérifier point par point aux frontières.
\end{enumerate}
\end{methode}

\courssub{Je vérifie mes connaissances}

\begin{exercice}[QCM : parité]
Pour chaque affirmation, choisir la seule bonne réponse.
\begin{enumerate}
\item La fonction $f$ définie sur $[-2\,;\,2]$ par $f(x)=x^3-3x$ est :
\begin{enumerate}
\item paire \quad \item impaire \quad \item ni paire ni impaire
\end{enumerate}
\item La fonction $g$ définie sur $[-3\,;\,3]$ par $g(x)=x^2+1$ est :
\begin{enumerate}
\item paire \quad \item impaire \quad \item ni paire ni impaire
\end{enumerate}
\item La fonction $h$ définie sur $[-1\,;\,2]$ par $h(x)=x^2$ est :
\begin{enumerate}
\item paire \quad \item impaire \quad \item ni paire ni impaire
\end{enumerate}
\item Si $f$ est impaire sur un domaine symétrique contenant $1$ et $f(1)=-4$, alors $f(-1)$ vaut :
\begin{enumerate}
\item $-4$ \quad \item $4$ \quad \item on ne peut pas savoir
\end{enumerate}
\end{enumerate}
\end{exercice}

\begin{exercice}[Vrai ou faux, en justifiant]
Répondre par Vrai ou Faux et justifier chaque réponse par une phrase ou un calcul.
\begin{enumerate}
\item Si $f$ est une fonction paire définie sur $[-5\,;\,5]$, alors sa courbe représentative est symétrique par rapport à l'axe des ordonnées.
\item Toute fonction définie sur $[-2\,;\,2]$ est paire ou impaire.
\item Si $\displaystyle\lim_{x\to 3^-}f(x)=2$ et $\displaystyle\lim_{x\to 3^+}f(x)=2$, alors $f$ admet une limite en $3$.
\item Si $\displaystyle\lim_{x\to x_0}f(x)=f(x_0)$, alors $f$ est continue en $x_0$.
\end{enumerate}
\end{exercice}

\begin{exercice}[Définitions à compléter]
Recopier et compléter les phrases suivantes.
\begin{enumerate}
\item Soit $f$ une fonction définie sur $D$. $f$ est \cle{paire} si pour tout $x\in D$, $-x\in D$ et $f(-x)=\ldots$
\item Soit $f$ une fonction définie sur $D$. $f$ est \cle{impaire} si pour tout $x\in D$, $-x\in D$ et $f(-x)=\ldots$
\item $f$ admet une limite en $x_0$ si et seulement si la limite à gauche et la limite à droite en $x_0$ existent et sont $\ldots$
\item $f$ est \cle{continue} en $x_0$ si $f$ est définie en $x_0$ et si $\displaystyle\lim_{x\to x_0}f(x)=\ldots$
\end{enumerate}
\end{exercice}

\begin{exercice}[Calculs directs de limites]
Déterminer, lorsqu'elle existe, la limite demandée.
\begin{enumerate}
\item $\displaystyle\lim_{x\to 1}(2x^2-5x+3)$ ;
\item $\displaystyle\lim_{x\to 2}\dfrac{x^2-4}{x-2}$ (on pourra factoriser le numérateur) ;
\item $\displaystyle\lim_{x\to 0^+}\dfrac{1}{x}$ et $\displaystyle\lim_{x\to 0^-}\dfrac{1}{x}$ ; que conclure sur la limite en $0$ ?
\end{enumerate}
\end{exercice}

\courssub{Je m'entraîne}

\begin{exercice}[Parité : le piège du domaine]
On considère les trois fonctions suivantes :
\[ f(x)=x^3-3x \text{ sur } [-2\,;\,2], \qquad g(x)=x^2+1 \text{ sur } [-3\,;\,3], \qquad h(x)=x^2 \text{ sur } [-1\,;\,2]. \]
\begin{enumerate}
\item Pour chaque fonction, vérifier d'abord si le domaine est symétrique par rapport à $0$.
\item Étudier la parité de $f$ et de $g$ en calculant $f(-x)$ et $g(-x)$.
\item La fonction $h$ peut-elle être paire ? Justifier soigneusement.
\end{enumerate}
\end{exercice}

\begin{exercice}[Conjecture par table de valeurs, puis confirmation]
Soit $f$ la fonction définie sur $\mathbb{R}\setminus\{2\}$ par $f(x)=\dfrac{x^2-4}{x-2}$.
\begin{enumerate}
\item Compléter la table de valeurs suivante (on donnera les résultats à $10^{-3}$ près) :
\begin{center}
\begin{tabular}{|c|cccccc|}\hline
$x$ & $1,9$ & $1,99$ & $1,999$ & $2,001$ & $2,01$ & $2,1$ \\ \hline
$f(x)$ & & & & & & \\ \hline
\end{tabular}
\end{center}
\item Conjecturer la limite de $f$ en $2$.
\item Confirmer cette conjecture en simplifiant l'écriture de $f(x)$ pour $x\neq 2$.
\end{enumerate}
\end{exercice}

\begin{exercice}[Limites à gauche et à droite]
Soit $f$ la fonction définie sur $\mathbb{R}\setminus\{-2\}$ par $f(x)=\dfrac{3x-1}{x+2}$.
\begin{enumerate}
\item Étudier le signe du dénominateur $x+2$ selon que $x<-2$ ou $x>-2$.
\item En déduire $\displaystyle\lim_{x\to -2^-}f(x)$ et $\displaystyle\lim_{x\to -2^+}f(x)$.
\item La fonction $f$ admet-elle une limite en $-2$ ? Justifier.
\end{enumerate}
\end{exercice}

\begin{exercice}[Continuité d'une fonction définie par morceaux]
Soit $f$ la fonction définie sur $\mathbb{R}$ par :
\[ f(x)=\begin{cases} x+1 & \text{si } x<1 \\ 3-x & \text{si } x\geq 1 \end{cases} \]
\begin{enumerate}
\item Calculer $f(1)$.
\item Déterminer $\displaystyle\lim_{x\to 1^-}f(x)$ et $\displaystyle\lim_{x\to 1^+}f(x)$.
\item Comparer ces résultats à $f(1)$ et conclure sur la continuité de $f$ en $1$.
\item Représenter graphiquement $f$ dans un repère orthonormé et vérifier la conclusion sur le graphique.
\end{enumerate}
\end{exercice}

\courssub{Je me prépare au Probatoire}

\begin{exercice}[Lecture graphique complète]
La courbe ci-dessous représente une fonction $f$ définie sur $[-4\,;\,4]$.
\begin{popfigure}
\begin{tikzpicture}[scale=0.85]
\draw[->,thick] (-4.6,0) -- (4.6,0) node[below left]{$x$};
\draw[->,thick] (0,-3.2) -- (0,3.4) node[below left]{$y$};
\foreach \x in {-4,-3,-2,-1,1,2,3,4} \draw (\x,0.08) -- (\x,-0.08) node[below]{\footnotesize $\x$};
\foreach \y in {-3,-2,-1,1,2,3} \draw (0.08,\y) -- (-0.08,\y) node[left]{\footnotesize $\y$};
% Branche 1 : 1/(x+1) sur [-4, -1.15]
\draw[popblue,very thick,domain=-4:-1.15,samples=100] plot (\x,{1/(\x+1)});
% Branche 2 : x^2-2 sur [-0.85, 1.95]
\draw[popblue,very thick,domain=-0.85:1.95,samples=100] plot (\x,{(\x)^2-2});
% Branche 3 : 1/(x-2) sur [2.05, 4]
\draw[popblue,very thick,domain=2.05:4,samples=100] plot (\x,{1/(\x-2)});
% Points aux bornes du domaine
\fill[popblue] (-4,-0.333) circle (1.6pt);
\fill[popblue] (4,0.5) circle (1.6pt);
% Asymptotes verticales pointillées
\draw[dashed,popmuted] (-1,-3.1) -- (-1,3.3);
\draw[dashed,popmuted] (2,-3.1) -- (2,3.3);
\end{tikzpicture}
\legende{Courbe représentative de $f$ sur $[-4\,;\,4]$}
\end{popfigure}
\emph{Toutes les réponses aux questions 1 à 4 seront justifiées par une lecture graphique.}
\begin{enumerate}
\item La fonction $f$ est-elle paire ? impaire ? ni l'une ni l'autre ?
\item Conjecturer $\displaystyle\lim_{x\to -1^-}f(x)$ et $\displaystyle\lim_{x\to -1^+}f(x)$. La fonction $f$ admet-elle une limite en $-1$ ?
\item Conjecturer de même les limites à gauche et à droite de $f$ en $2$.
\item Donner, sous forme de réunion d'intervalles, l'ensemble des intervalles de $[-4\,;\,4]$ sur lesquels $f$ semble continue.
\item Dresser le tableau de variations de $f$ sur $[-4\,;\,4]$.
\end{enumerate}
\end{exercice}

\begin{exercice}[Le stand de beignets de Mme Etoundi]
Mme Etoundi tient un stand de beignets au marché de Mokolo. Le bénéfice journalier, en centaines de francs CFA, réalisé pour $x$ dizaines de beignets vendus est modélisé par la fonction $B$ définie par :
\[ B(x)=\begin{cases} 2x+4 & \text{si } 0\leq x<5 \\ 24-x & \text{si } x\geq 5 \end{cases} \]
\begin{enumerate}
\item Quel est le bénéfice réalisé pour $30$ beignets vendus ? pour $50$ beignets vendus ?
\item On considère la restriction de $B$ à l'intervalle $[-5\,;\,5]$. Rappel : le domaine de définition de $B$ est $[0\,;\,+\infty[$. Quel est le domaine de cette restriction ? Peut-on étudier sa parité ? Justifier.
\item Déterminer $\displaystyle\lim_{x\to 5^-}B(x)$ et $\displaystyle\lim_{x\to 5^+}B(x)$.
\item La fonction $B$ est-elle continue en $5$ ? Interpréter ce résultat pour Mme Etoundi : que se passe-t-il autour de la vente de $50$ beignets ?
\item Justifier que la fonction $x\mapsto 2x+4$ est continue sur $[0\,;\,5[$ et que la fonction $x\mapsto 24-x$ est continue sur $]5\,;\,+\infty[$.
\item Représenter graphiquement $B$ dans un repère orthogonal (unités : \SI{1}{cm} pour une dizaine de beignets en abscisse, \SI{1}{cm} pour $500$ FCFA en ordonnée).
\end{enumerate}
\end{exercice}

\begin{exercice}[Étude complète d'une fonction homographique]
On considère la fonction $f$ définie par $f(x)=\dfrac{x-3}{x+1}$ et on note $\mathcal{C}_f$ sa courbe représentative dans un repère orthonormé $(O\,;\vec{i},\vec{j})$.
\begin{enumerate}
\item Déterminer le domaine de définition $D_f$ de $f$.
\item Le domaine $D_f$ est-il symétrique par rapport à $0$ ? La fonction $f$ peut-elle être paire ou impaire ? Justifier.
\item \textbf{Étude des limites en $-1$.}
\begin{enumerate}
\item Étudier le signe de $x+1$ selon les valeurs de $x$.
\item En déduire $\displaystyle\lim_{x\to -1^-}f(x)$ et $\displaystyle\lim_{x\to -1^+}f(x)$.
\item La fonction $f$ admet-elle une limite en $-1$ ? Justifier.
\end{enumerate}
\item \textbf{Continuité.}
\begin{enumerate}
\item Justifier que $f$ est continue sur $]-1\,;\,+\infty[$.
\item $f$ est-elle continue sur $[-2\,;\,0]$ ? Justifier.
\end{enumerate}
\item \textbf{Approche numérique.} Compléter la table de valeurs ci-dessous (résultats à $10^{-2}$ près) et vérifier la cohérence avec les résultats de la question 3 :
\begin{center}
\begin{tabular}{|c|cccc|}\hline
$x$ & $-1,1$ & $-1,01$ & $-0,99$ & $-0,9$ \\ \hline
$f(x)$ & & & & \\ \hline
\end{tabular}
\end{center}
\item Tracer l'allure de $\mathcal{C}_f$ au voisinage de la droite d'équation $x=-1$, en utilisant les résultats précédents.
\end{enumerate}
\end{exercice}

\newpage

\coursec{Corrigés des exercices}

\begin{corrige}[Exercice 1 — QCM : parité]
\begin{enumerate}
\item \textbf{Réponse b : impaire.} Le domaine $[-2\,;\,2]$ est symétrique par rapport à $0$. Pour tout $x\in[-2\,;\,2]$ :
\[ f(-x)=(-x)^3-3(-x)=-x^3+3x=-(x^3-3x)=-f(x). \]
Donc $f$ est impaire.
\item \textbf{Réponse a : paire.} Le domaine $[-3\,;\,3]$ est symétrique. Pour tout $x\in[-3\,;\,3]$ :
\[ g(-x)=(-x)^2+1=x^2+1=g(x). \]
Donc $g$ est paire.
\item \textbf{Réponse c : ni paire ni impaire.} Le domaine $[-1\,;\,2]$ n'est \textbf{pas} symétrique par rapport à $0$ : par exemple $2\in[-1\,;\,2]$ mais $-2\notin[-1\,;\,2]$. La condition « pour tout $x\in D$, $-x\in D$ » n'est pas satisfaite, donc $h$ n'est ni paire ni impaire.
\item \textbf{Réponse b : $4$.} Comme $f$ est impaire, $f(-1)=-f(1)=-(-4)=4$.
\end{enumerate}
\end{corrige}

\begin{attention}[Point de vigilance]
Pour la question 3, la tentation est grande de répondre « paire » car $x^2$ est une expression paire. Mais la parité concerne \textbf{la fonction}, c'est-à-dire la formule \textbf{et} le domaine. Un domaine non symétrique disqualifie immédiatement toute parité.
\end{attention}

\begin{corrige}[Exercice 2 — Vrai ou faux, en justifiant]
\begin{enumerate}
\item \textbf{Vrai.} C'est l'interprétation graphique de la parité : $f$ paire sur un domaine symétrique équivaut à « la courbe $\mathcal{C}_f$ est symétrique par rapport à l'axe des ordonnées ». Ici $[-5\,;\,5]$ est bien symétrique.
\item \textbf{Faux.} La symétrie du domaine est une condition \textbf{nécessaire} mais pas suffisante. Contre-exemple : $f(x)=x^2+x$ sur $[-2\,;\,2]$. On a $f(-x)=x^2-x$, qui n'est égal ni à $f(x)$ ni à $-f(x)$ (par exemple $f(1)=2$ et $f(-1)=0$). Donc $f$ n'est ni paire ni impaire.
\item \textbf{Vrai.} D'après le théorème de la limite bilatère, si les limites à gauche et à droite en $3$ existent, sont finies et égales à $2$, alors $f$ admet une limite en $3$ et $\lim_{x\to 3}f(x)=2$. (Attention : cela ne dit rien sur la continuité, car on ne connaît pas $f(3)$.)
\item \textbf{Vrai.} C'est exactement la définition de la continuité en $x_0$ : $f$ est définie en $x_0$ (puisque $f(x_0)$ existe dans l'égalité) et $\lim_{x\to x_0}f(x)=f(x_0)$, donc $f$ est continue en $x_0$.
\end{enumerate}
\end{corrige}

\begin{corrige}[Exercice 3 — Définitions à compléter]
\begin{enumerate}
\item $f$ est paire si pour tout $x\in D$, $-x\in D$ et $f(-x)=\mathbf{f(x)}$.
\item $f$ est impaire si pour tout $x\in D$, $-x\in D$ et $f(-x)=\mathbf{-f(x)}$.
\item $f$ admet une limite en $x_0$ si et seulement si la limite à gauche et la limite à droite en $x_0$ existent et sont \textbf{égales} (et finies).
\item $f$ est continue en $x_0$ si $f$ est définie en $x_0$ et si $\displaystyle\lim_{x\to x_0}f(x)=\mathbf{f(x_0)}$.
\end{enumerate}
\end{corrige}

\begin{corrige}[Exercice 4 — Calculs directs de limites]
\begin{enumerate}
\item $x\mapsto 2x^2-5x+3$ est une fonction polynomiale, donc :
\[ \lim_{x\to 1}(2x^2-5x+3)=2(1)^2-5(1)+3=2-5+3=0. \]
\item Pour $x\neq 2$, on factorise le numérateur (identité remarquable $a^2-b^2=(a-b)(a+b)$) :
\[ \frac{x^2-4}{x-2}=\frac{(x-2)(x+2)}{x-2}=x+2. \]
Donc $\displaystyle\lim_{x\to 2}\frac{x^2-4}{x-2}=\lim_{x\to 2}(x+2)=4$.
\item Pour $x>0$, $\dfrac{1}{x}>0$ et devient arbitrairement grand quand $x$ se rapproche de $0$ : $\displaystyle\lim_{x\to 0^+}\frac{1}{x}=+\infty$.
Pour $x<0$, $\dfrac{1}{x}<0$ et devient arbitrairement grand en valeur absolue : $\displaystyle\lim_{x\to 0^-}\frac{1}{x}=-\infty$.
\textbf{Conclusion :} les limites à gauche et à droite sont différentes (et infinies), donc $x\mapsto\frac1x$ \textbf{n'admet pas de limite} en $0$.
\end{enumerate}
\end{corrige}

\begin{attention}[Point de vigilance]
La notation $\frac{1}{0}$ est \textbf{interdite} dans la rédaction. On écrit le raisonnement sur le signe et on conclut par une limite infinie, jamais par une division par zéro.
\end{attention}

\begin{corrige}[Exercice 5 — Parité : le piège du domaine]
\begin{enumerate}
\item $[-2\,;\,2]$ est symétrique : si $x\in[-2\,;\,2]$, alors $-x\in[-2\,;\,2]$. De même $[-3\,;\,3]$ est symétrique. En revanche $[-1\,;\,2]$ n'est \textbf{pas} symétrique : $2\in[-1\,;\,2]$ mais $-2\notin[-1\,;\,2]$.
\item \textbf{Parité de $f$ :} pour tout $x\in[-2\,;\,2]$,
\[ f(-x)=(-x)^3-3(-x)=-x^3+3x=-(x^3-3x)=-f(x). \]
Donc $f$ est \textbf{impaire}.
\textbf{Parité de $g$ :} pour tout $x\in[-3\,;\,3]$,
\[ g(-x)=(-x)^2+1=x^2+1=g(x). \]
Donc $g$ est \textbf{paire}.
\item \textbf{Non}, $h$ ne peut pas être paire. La définition de la parité exige que pour \textbf{tout} $x$ du domaine, $-x$ soit aussi dans le domaine. Or $2\in[-1\,;\,2]$ et $-2\notin[-1\,;\,2]$ : l'expression $h(-2)$ n'a même pas de sens. Le domaine n'étant pas symétrique, $h$ est \textbf{ni paire ni impaire}, bien que sa formule $x^2$ soit « paire » sur $\mathbb{R}$.
\end{enumerate}
\end{corrige}

\begin{corrige}[Exercice 6 — Conjecture par table de valeurs, puis confirmation]
\begin{enumerate}
\item En calculant $f(x)=\dfrac{x^2-4}{x-2}$ pour chaque valeur :
\begin{center}
\begin{tabular}{|c|cccccc|}\hline
$x$ & $1,9$ & $1,99$ & $1,999$ & $2,001$ & $2,01$ & $2,1$ \\ \hline
$f(x)$ & $3,900$ & $3,990$ & $3,999$ & $4,001$ & $4,010$ & $4,100$ \\ \hline
\end{tabular}
\end{center}
Détail d'un calcul : $f(1,9)=\dfrac{1,9^2-4}{1,9-2}=\dfrac{3,61-4}{-0,1}=\dfrac{-0,39}{-0,1}=3,9$.
\item On observe que $f(x)$ se rapproche de $4$ quand $x$ se rapproche de $2$, aussi bien par valeurs inférieures que supérieures. On \textbf{conjecture} : $\displaystyle\lim_{x\to 2}f(x)=4$.
\item Pour $x\neq 2$ :
\[ f(x)=\frac{x^2-4}{x-2}=\frac{(x-2)(x+2)}{x-2}=x+2. \]
Or $\displaystyle\lim_{x\to 2}(x+2)=4$ (fonction polynomiale). Donc $\displaystyle\lim_{x\to 2}f(x)=4$ : la conjecture est \textbf{confirmée}.
\end{enumerate}
\end{corrige}

\begin{attention}[Point de vigilance]
La simplification par $x-2$ n'est licite que parce que $x\neq 2$ : c'est précisément le cas quand on cherche une limite en $2$ ($x$ s'approche de $2$ sans jamais l'atteindre). La fonction $f$ reste non définie en $2$.
\end{attention}

\begin{corrige}[Exercice 7 — Limites à gauche et à droite]
\begin{enumerate}
\item $x+2>0 \iff x>-2$ et $x+2<0 \iff x<-2$.
\item Le numérateur en $-2$ vaut $3\times(-2)-1=-7<0$.
\begin{itemize}
\item Quand $x\to -2^-$ : $x<-2$, donc $x+2<0$ et $x+2\to 0^-$. Le quotient $\dfrac{-7}{x+2}$ est le quotient de deux quantités négatives, donc positif, et devient arbitrairement grand :
\[ \lim_{x\to -2^-}f(x)=+\infty. \]
\item Quand $x\to -2^+$ : $x>-2$, donc $x+2>0$ et $x+2\to 0^+$. Le quotient $\dfrac{-7}{x+2}$ est négatif et devient arbitrairement grand en valeur absolue :
\[ \lim_{x\to -2^+}f(x)=-\infty. \]
\end{itemize}
\item \textbf{Non}, $f$ n'admet pas de limite en $-2$, car les limites à gauche ($+\infty$) et à droite ($-\infty$) sont différentes (et de plus infinies). La droite d'équation $x=-2$ est une asymptote verticale à la courbe de $f$.
\end{enumerate}
\end{corrige}

\begin{corrige}[Exercice 8 — Continuité d'une fonction définie par morceaux]
\begin{enumerate}
\item Comme $1\geq 1$, on utilise la deuxième expression : $f(1)=3-1=2$.
\item \begin{itemize}
\item À gauche de $1$ ($x<1$), $f(x)=x+1$ : $\displaystyle\lim_{x\to 1^-}f(x)=\lim_{x\to 1^-}(x+1)=2$.
\item À droite de $1$ ($x>1$), $f(x)=3-x$ : $\displaystyle\lim_{x\to 1^+}f(x)=\lim_{x\to 1^+}(3-x)=2$.
\end{itemize}
\item On a $\displaystyle\lim_{x\to 1^-}f(x)=\lim_{x\to 1^+}f(x)=2=f(1)$. Les trois conditions de la définition sont réunies : $f$ est définie en $1$, admet une limite finie en $1$, et cette limite vaut $f(1)$. Donc $f$ est \textbf{continue en $1$}.
\item La représentation graphique est formée de deux demi-droites qui se rejoignent au point $(1\,;\,2)$ : on trace la courbe sans lever le crayon, ce qui confirme la continuité.
\begin{popfigure}
\begin{tikzpicture}[scale=0.8]
\draw[->,thick] (-3.4,0) -- (4.4,0) node[below left]{$x$};
\draw[->,thick] (0,-2.4) -- (0,3.4) node[below left]{$y$};
\foreach \x in {-3,-2,-1,1,2,3,4} \draw (\x,0.07) -- (\x,-0.07) node[below]{\footnotesize $\x$};
\foreach \y in {-2,-1,1,2,3} \draw (0.07,\y) -- (-0.07,\y) node[left]{\footnotesize $\y$};
\draw[popblue,very thick,domain=-3:1,samples=2] plot (\x,{\x+1});
\draw[poporange,very thick,domain=1:4,samples=2] plot (\x,{3-\x});
\fill[popblue] (1,2) circle (2pt);
\end{tikzpicture}
\legende{Les deux morceaux se rejoignent au point $(1\,;\,2)$ : $f$ est continue en $1$.}
\end{popfigure}
\end{enumerate}
\end{corrige}

\begin{corrige}[Exercice 9 — Lecture graphique complète]
\emph{Barème indicatif (10 pts) : question 1 : 2 pts ; question 2 : 2 pts ; question 3 : 2 pts ; question 4 : 2 pts ; question 5 : 2 pts.}
\begin{enumerate}
\item La courbe n'est symétrique ni par rapport à l'axe des ordonnées (par exemple la branche de gauche s'envole vers $-\infty$ près de $-1$ alors que rien de tel n'apparaît près de $1$), ni par rapport à l'origine. Donc $f$ est \textbf{ni paire ni impaire}.
\item Lecture graphique au voisinage de $-1$ : quand $x$ s'approche de $-1$ par valeurs inférieures, la courbe (branche de gauche) descend vers le bas indéfiniment ; quand $x$ s'approche de $-1$ par valeurs supérieures, la courbe (branche centrale, une parabole) se rapproche d'une valeur bien précise, environ $-1$ :
\[ \lim_{x\to -1^-}f(x)=-\infty \qquad \text{et} \qquad \lim_{x\to -1^+}f(x)=-1. \]
La première limite est infinie, la seconde est finie : elles sont donc différentes, et $f$ \textbf{n'admet pas de limite} en $-1$ (la courbe présente un \og saut \fg{} en ce point).
\item De même au voisinage de $2$ : la branche centrale (la parabole) se rapproche d'une valeur finie, environ $2$, quand $x\to 2^-$ ; la branche de droite s'envole vers $+\infty$ quand $x\to 2^+$ :
\[ \lim_{x\to 2^-}f(x)=2 \qquad \text{et} \qquad \lim_{x\to 2^+}f(x)=+\infty. \]
Là encore, une limite finie et une limite infinie sont différentes : $f$ \textbf{n'admet pas de limite} en $2$.
\item On peut tracer la courbe sans lever le crayon sur chacun des intervalles $[-4\,;\,-1[$, $]-1\,;\,2[$ et $]2\,;\,4]$. La fonction $f$ semble donc continue sur :
\[ [-4\,;\,-1[\;\cup\;]-1\,;\,2[\;\cup\;]2\,;\,4]. \]
\item Sur $[-4\,;\,-1[$, $f$ est décroissante de $f(-4)=-\tfrac13$ vers $-\infty$. Sur $]-1\,;\,2[$, la branche centrale est la parabole d'équation $y=x^2-2$ (prolongée mentalement jusqu'aux bords) : elle vaut $-1$ tout près de $-1$, décroît jusqu'à son minimum $-2$ en $0$, puis croît jusqu'à une valeur proche de $2$ tout près de $2$. Sur $]2\,;\,4]$, $f$ décroît de $+\infty$ vers $f(4)=\tfrac12$.
\begin{center}
\begin{tabular}{|c|ccccccccc|}\hline
$x$ & $-4$ & & $-1$ & & $0$ & & $2$ & & $4$ \\ \hline
$f'(x)$ & & $-$ & \begin{tabular}{c}$\|$\\$\|$\end{tabular} & $-$ & & $+$ & \begin{tabular}{c}$\|$\\$\|$\end{tabular} & $-$ & \\ \hline
$f$ & $-\tfrac13$ & $\searrow$ & \begin{tabular}{c}$-\infty$\\$-1$\end{tabular} & $\searrow$ & $-2$ & $\nearrow$ & \begin{tabular}{c}$2$\\$+\infty$\end{tabular} & $\searrow$ & $\tfrac12$ \\ \hline
\end{tabular}
\end{center}
(Les doubles barres en $-1$ et $2$ matérialisent les valeurs interdites. En $-1$, la limite à gauche est infinie ($-\infty$) tandis que la limite à droite est finie ($-1$) : c'est un \og saut \fg. En $2$, c'est l'inverse : limite finie à gauche ($2$), limite infinie à droite ($+\infty$). La ligne $f'$ porte les signes de la dérivée, déduits des variations : $-$ pour décroissante, $+$ pour croissante.)
\end{enumerate}
\end{corrige}

\begin{attention}[Points de vigilance au Probatoire]
\begin{itemize}
\item Justifier chaque réponse par la phrase « par lecture graphique » et décrire ce que fait la courbe (elle « descend », « monte », « se rapproche de »).
\item Ne jamais écrire $f(-1)$ ou $f(2)$ : ces images n'existent pas, les valeurs $-1$ et $2$ sont interdites.
\item Dans le tableau de variations, les valeurs interdites imposent une double barre et des limites, pas des images. La ligne $f'$ doit être présente avec les signes de variation.
\end{itemize}
\end{attention}

\begin{corrige}[Exercice 10 — Le stand de beignets de Mme Etoundi]
\emph{Barème indicatif (10 pts) : question 1 : 1 pt ; question 2 : 2 pts ; question 3 : 2 pts ; question 4 : 2 pts ; question 5 : 2 pts ; question 6 : 1 pt.}
\begin{enumerate}
\item $30$ beignets correspondent à $x=3$ dizaines : comme $3<5$, $B(3)=2\times 3+4=10$ centaines de francs, soit $\SI{1000}{FCFA}$.
$50$ beignets correspondent à $x=5$ : comme $5\geq 5$, $B(5)=24-5=19$ centaines de francs, soit $\SI{1900}{FCFA}$.
\item La restriction de $B$ à $[-5\,;\,5]$ a pour domaine $[0\,;\,+\infty[\;\cap\;[-5\,;\,5]=[0\,;\,5]$. Ce domaine n'est \textbf{pas} symétrique par rapport à $0$ (par exemple $5\in[0\,;\,5]$ mais $-5\notin[0\,;\,5]$). On ne peut donc \textbf{pas} étudier la parité de cette restriction : elle n'est ni paire ni impaire.
\item \begin{itemize}
\item À gauche de $5$ ($x<5$), $B(x)=2x+4$ : $\displaystyle\lim_{x\to 5^-}B(x)=2\times 5+4=14$.
\item À droite de $5$ ($x>5$), $B(x)=24-x$ : $\displaystyle\lim_{x\to 5^+}B(x)=24-5=19$.
\end{itemize}
\item On a $\displaystyle\lim_{x\to 5^-}B(x)=14\neq 19=\lim_{x\to 5^+}B(x)$. La fonction $B$ n'admet donc pas de limite en $5$, et par conséquent $B$ est \textbf{discontinue en $5$} (discontinuité de type « saut »).
\textbf{Interprétation :} le modèle prévoit un « saut » de bénéfice autour de $50$ beignets : en vendant juste moins de $50$ beignets, le bénéfice approche $\SI{1400}{FCFA}$, tandis qu'à $50$ beignets il vaut $\SI{1900}{FCFA}$. Le modèle comporte une rupture (par exemple une prime ou un changement de tarif à partir de $50$ beignets) ; en réalité, un bénéfice évolue rarement par saut, ce qui invite Mme Etoundi à vérifier la validité du modèle.
\item Les fonctions $x\mapsto 2x+4$ et $x\mapsto 24-x$ sont des fonctions \textbf{polynomiales} (affines) ; or toute fonction polynomiale est continue sur $\mathbb{R}$, donc en particulier la première est continue sur $[0\,;\,5[$ et la seconde sur $]5\,;\,+\infty[$.
\item La représentation est formée de deux demi-droites : celle de $y=2x+4$ du point $(0\,;\,4)$ jusqu'au point $(5\,;\,14)$ \textbf{non inclus} (point creux), et celle de $y=24-x$ à partir du point $(5\,;\,19)$ \textbf{inclus} (point plein). Le saut entre les points $(5\,;\,14)$ et $(5\,;\,19)$ matérialise la discontinuité.
\begin{popfigure}
\begin{tikzpicture}[x=0.55cm,y=0.18cm]
\draw[->,thick] (-0.6,0) -- (10.6,0) node[below left]{$x$ (dizaines)};
\draw[->,thick] (0,-1.4) -- (0,26) node[below left]{$B(x)$ (centaines F)};
\foreach \x in {1,2,3,4,5,6,7,8,9,10} \draw (\x,0.5) -- (\x,-0.5) node[below]{\footnotesize $\x$};
\foreach \y in {4,8,12,16,20,24} \draw (0.1,\y) -- (-0.1,\y) node[left]{\footnotesize $\y$};
\draw[popblue,very thick,domain=0:5,samples=2] plot (\x,{2*\x+4});
\draw[poporange,very thick,domain=5:10,samples=2] plot (\x,{24-\x});
\draw[popblue,fill=white] (5,14) circle (2.5pt);
\fill[poporange] (5,19) circle (2.5pt);
\end{tikzpicture}
\legende{Point creux en $(5\,;\,14)$, point plein en $(5\,;\,19)$ : saut de discontinuité en $5$.}
\end{popfigure}
\end{enumerate}
\end{corrige}

\begin{attention}[Points de vigilance]
\begin{itemize}
\item Pour $x=5$, utiliser la bonne branche : la condition $x\geq 5$ impose $B(5)=24-5$.
\item Bien distinguer « $B$ n'admet pas de limite en $5$ » (limites unilatérales différentes) et « $B$ est discontinue en $5$ » (conséquence).
\item Sur le graphique, le point creux / point plein est exigé pour une lecture correcte.
\end{itemize}
\end{attention}

\begin{corrige}[Exercice 11 — Étude complète d'une fonction homographique]
\emph{Barème indicatif (12 pts) : question 1 : 1 pt ; question 2 : 2 pts ; question 3 : 3 pts ; question 4 : 2 pts ; question 5 : 2 pts ; question 6 : 2 pts.}
\begin{enumerate}
\item $f(x)$ existe si et seulement si $x+1\neq 0$, c'est-à-dire $x\neq -1$. Donc :
\[ D_f=\mathbb{R}\setminus\{-1\}=]-\infty\,;\,-1[\;\cup\;]-1\,;\,+\infty[. \]
\item $D_f$ \textbf{est symétrique} par rapport à $0$ : si $x\neq -1$, alors $-x\neq 1$... vérifions proprement : si $x\in D_f$ alors $x\neq -1$ ; son opposé $-x$ vérifie $-x\neq 1$, donc $-x\neq -1$ sauf si... Prenons un contre-exemple : $1\in D_f$ (car $1\neq -1$) mais $-1\notin D_f$. Donc $D_f$ n'est \textbf{pas} symétrique par rapport à $0$. Par conséquent $f$ ne peut être \textbf{ni paire ni impaire}.
\item \textbf{Étude des limites en $-1$.}
\begin{enumerate}
\item $x+1>0 \iff x>-1$ et $x+1<0 \iff x<-1$.
\item Le numérateur en $-1$ vaut $-1-3=-4<0$.
\begin{itemize}
\item Quand $x\to -1^-$ : $x<-1$, donc $x+1\to 0^-$. Le quotient $\dfrac{-4}{x+1}$ est positif (négatif divisé par négatif) et devient arbitrairement grand :
\[ \lim_{x\to -1^-}f(x)=+\infty. \]
\item Quand $x\to -1^+$ : $x>-1$, donc $x+1\to 0^+$. Le quotient $\dfrac{-4}{x+1}$ est négatif et devient arbitrairement grand en valeur absolue :
\[ \lim_{x\to -1^+}f(x)=-\infty. \]
\end{itemize}
\item Les limites à gauche ($+\infty$) et à droite ($-\infty$) sont différentes et infinies : $f$ \textbf{n'admet pas de limite} en $-1$. La droite d'équation $x=-1$ est asymptote verticale à $\mathcal{C}_f$.
\end{enumerate}
\item \textbf{Continuité.}
\begin{enumerate}
\item $f$ est une fonction \textbf{rationnelle} (quotient de deux polynômes) ; elle est donc continue sur son domaine de définition $D_f$. Comme $]-1\,;\,+\infty[\;\subset D_f$, $f$ est continue sur $]-1\,;\,+\infty[$.
\item \textbf{Non.} L'intervalle $[-2\,;\,0]$ contient $-1$, où $f$ n'est même pas définie (et admet des limites infinies). Donc $f$ n'est pas continue sur $[-2\,;\,0]$.
\end{enumerate}
\item \textbf{Approche numérique.} $f(-1,1)=\dfrac{-1,1-3}{-1,1+1}=\dfrac{-4,1}{-0,1}=41$ ; $f(-1,01)=\dfrac{-4,01}{-0,01}=401$ ; $f(-0,99)=\dfrac{-3,99}{0,01}=-399$ ; $f(-0,9)=\dfrac{-3,9}{0,1}=-39$.
\begin{center}
\begin{tabular}{|c|cccc|}\hline
$x$ & $-1,1$ & $-1,01$ & $-0,99$ & $-0,9$ \\ \hline
$f(x)$ & $41,00$ & $401,00$ & $-399,00$ & $-39,00$ \\ \hline
\end{tabular}
\end{center}
Cohérence : à gauche de $-1$ ($-1,1$ et $-1,01$), les valeurs sont positives et de plus en plus grandes ($41$, puis $401$) : cela confirme $\lim_{x\to -1^-}f(x)=+\infty$. À droite de $-1$ ($-0,99$ et $-0,9$), les valeurs sont négatives et de plus en plus grandes en valeur absolue ($-399$, puis $-39$ en s'éloignant) : cela confirme $\lim_{x\to -1^+}f(x)=-\infty$.
\item Au voisinage de la droite $x=-1$, la courbe « s'envole » vers $+\infty$ à gauche de l'asymptote et « plonge » vers $-\infty$ à droite :
\begin{popfigure}
\begin{tikzpicture}[scale=0.9]
\draw[->,thick] (-4.4,0) -- (3.4,0) node[below left]{$x$};
\draw[->,thick] (0,-4.4) -- (0,4.6) node[below left]{$y$};
\foreach \x in {-4,-3,-2,-1,1,2,3} \draw (\x,0.07) -- (\x,-0.07) node[below]{\footnotesize $\x$};
\foreach \y in {-4,-3,-2,-1,1,2,3,4} \draw (0.07,\y) -- (-0.07,\y) node[left]{\footnotesize $\y$};
\draw[dashed,popmuted] (-1,-4.3) -- (-1,4.5);
\draw[popblue,very thick,domain=-4:-1.25,samples=150] plot (\x,{(\x-3)/(\x+1)});
\draw[popblue,very thick,domain=-0.76:3,samples=150] plot (\x,{(\x-3)/(\x+1)});
\end{tikzpicture}
\legende{Allure de $\mathcal{C}_f$ : asymptote verticale d'équation $x=-1$.}
\end{popfigure}
\end{enumerate}
\end{corrige}

\begin{attention}[Points de vigilance au Probatoire]
\begin{itemize}
\item Question 2 : citer explicitement le contre-exemple ($1\in D_f$ mais $-1\notin D_f$) ; c'est la justification attendue.
\item Question 3 : rédiger le signe du dénominateur \textbf{avant} de conclure ; interdiction d'écrire $\frac{-4}{0}$.
\item Question 4b : la présence de la valeur interdite $-1$ dans l'intervalle suffit à conclure ; dire « $f$ n'est pas définie en $-1\in[-2\,;\,0]$ ».
\item Question 5 : les valeurs numériques doivent être cohérentes avec les limites trouvées ; un écart signale une erreur de calcul à corriger.
\end{itemize}
\end{attention}

\newpage

\coursec{Fiche bilan}

\begin{popfigure}
\begin{tikzpicture}[mindmap, grow cyclic, every node/.style={concept, font=\small},
  root concept/.append style={concept color=popblue, font=\bfseries},
  level 1/.append style={level distance=3.6cm, sibling angle=60},
  level 2/.append style={level distance=2.2cm, sibling angle=45, font=\scriptsize}]
\node[root concept]{Fonctions : parité, limites, continuité}
  child[concept color=poporange]{ node {Parité}
    child { node {Paire : $f(-x)=f(x)$} }
    child { node {Impaire : $f(-x)=-f(x)$} }
  }
  child[concept color=popgreen]{ node {Symétries}
    child { node {Paire : axe $(Oy)$} }
    child { node {Impaire : centre $O$} }
  }
  child[concept color=poppurple]{ node {Limite en $x_0$}
    child { node {Table de valeurs} }
    child { node {Lecture graphique} }
  }
  child[concept color=popgold]{ node {Limite à droite / à gauche}
    child { node {$\lim\limits_{x\to x_0^+} f$ ; $\lim\limits_{x\to x_0^-} f$} }
    child { node {Cas $\dfrac{ax+b}{cx+d}$ en $-\dfrac{d}{c}$} }
  }
  child[concept color=poppink]{ node {Continuité}
    child { node {En $x_0$ : $\lim f = f(x_0)$} }
    child { node {Sur $]a;b[$ : courbe sans saut} }
  }
  child[concept color=popblueL]{ node {Méthodes}
    child { node {Vérifier $D$ symétrique} }
    child { node {Rédaction rigoureuse} }
  };
\end{tikzpicture}
\legende{Carte mentale de la leçon : parité, limites et continuité.}
\end{popfigure}

\begin{aretenir}[L'essentiel de la leçon]
\textbf{1. Définitions clés}
\begin{itemize}
  \item \cle{Fonction paire} : pour tout $x\in D$, $-x\in D$ et $f(-x)=f(x)$. La courbe représentative est symétrique par rapport à l'axe des ordonnées.
  \item \cle{Fonction impaire} : pour tout $x\in D$, $-x\in D$ et $f(-x)=-f(x)$. La courbe représentative est symétrique par rapport à l'origine $O$ du repère.
  \item \cle{Condition préalable} : le domaine $D$ doit être \textbf{symétrique par rapport à $0$} (c'est-à-dire : si $x\in D$ alors $-x\in D$), sinon la fonction n'est ni paire ni impaire.
  \item \cle{Limite en $x_0$} : $f(x)$ se rapproche d'un réel $\ell$ quand $x$ se rapproche de $x_0$ ; on note $\lim\limits_{x\to x_0} f(x)=\ell$.
  \item \cle{Limites à droite et à gauche} : on note $\lim\limits_{x\to x_0^+} f(x)$ (limite par valeurs supérieures à $x_0$) et $\lim\limits_{x\to x_0^-} f(x)$ (limite par valeurs inférieures à $x_0$).
  \item \cle{Continuité en $x_0$} : $f$ est définie en $x_0$ et $\lim\limits_{x\to x_0} f(x)=f(x_0)$.
  \item \cle{Continuité sur $]a;b[$} : la courbe se trace \og sans lever le crayon \fg{} sur tout l'intervalle.
\end{itemize}

\textbf{2. Résultats à connaître par cœur}
\begin{center}
\begin{tabularx}{\linewidth}{|l|X|X|}
\hline
\rowcolor{popblueL} \textbf{Notion} & \textbf{Critère} & \textbf{Conséquence graphique} \\
\hline
Parité & $f(-x)=f(x)$ sur $D$ symétrique par rapport à $0$ & Symétrie par rapport à l'axe $(Oy)$ \\
\hline
Imparité & $f(-x)=-f(x)$ sur $D$ symétrique par rapport à $0$ & Symétrie par rapport à l'origine $O$ \\
\hline
Limite en $x_0$ & $\lim\limits_{x\to x_0^-}f=\lim\limits_{x\to x_0^+}f=\ell$ & Les deux branches de la courbe rejoignent la même valeur $\ell$ \\
\hline
Continuité en $x_0$ & $\lim\limits_{x\to x_0}f(x)=f(x_0)$ & Pas de saut ni de trou en $x_0$ \\
\hline
\end{tabularx}
\end{center}

\textbf{3. Méthodes express}
\begin{itemize}
  \item \cle{Parité} : (1) vérifier que $D$ est symétrique par rapport à $0$ ; (2) calculer $f(-x)$ ; (3) comparer à $f(x)$ et à $-f(x)$ ; (4) conclure.
  \item \cle{Limite par table de valeurs} : calculer $f(x)$ pour $x$ de plus en plus proche de $x_0$ (à gauche puis à droite) et observer la valeur approchée.
  \item \cle{Fonction homographique} $x\mapsto\dfrac{ax+b}{cx+d}$ : étudier les limites à gauche et à droite en $-\dfrac{d}{c}$ en exploitant la courbe de $x\mapsto\dfrac{a}{x}$ (asymptote verticale).
  \item \cle{Continuité} : justifier graphiquement (courbe sans interruption) ou identifier $f$ à une fonction continue du programme (polynôme, inverse sur son domaine, etc.).
  \item La notation $\dfrac{\ell}{0}$ ($\ell\in\mathbb{R}^*$) est \textbf{proscrite} : on écrit une limite infinie ou une limite à gauche/droite.
\end{itemize}
\end{aretenir}

\begin{attention}[Pièges fréquents au Probatoire]
\begin{itemize}
  \item Conclure sur la parité \textbf{sans avoir vérifié} que le domaine est symétrique par rapport à $0$ : c'est une erreur éliminatoire dans la justification.
  \item Confondre $f(-x)=-f(x)$ (fonction impaire) avec $-f(-x)=f(x)$ : les deux écritures sont en fait équivalentes, mais il faut savoir passer de l'une à l'autre proprement.
  \item Écrire $\lim\limits_{x\to x_0}f(x)$ alors que les limites à gauche et à droite sont \textbf{différentes} : dans ce cas, la limite en $x_0$ n'existe pas.
  \item Dire qu'une fonction est continue en $x_0$ alors qu'elle n'y est pas définie : la continuité exige d'abord que $f(x_0)$ existe.
  \item Utiliser la notation $\dfrac{\ell}{0}$ : elle est interdite. On rédige : \og la limite à droite est $+\infty$ \fg{} ou \og la limite à gauche est $-\infty$ \fg{} selon le cas.
\end{itemize}
\end{attention}

\begin{aretenir}[Checklist avant le Probatoire]
\begin{itemize}
  \item[$\square$] Je vérifie \textbf{toujours} d'abord que le domaine $D$ est symétrique par rapport à $0$ avant d'étudier la parité.
  \item[$\square$] Je sais calculer $f(-x)$ et le comparer à $f(x)$ et à $-f(x)$.
  \item[$\square$] Je connais les éléments de symétrie : axe $(Oy)$ (fonction paire), origine $O$ (fonction impaire).
  \item[$\square$] Je sais conjecturer une limite à l'aide d'une table de valeurs ou d'une lecture graphique.
  \item[$\square$] Je distingue limite à gauche ($x\to x_0^-$) et limite à droite ($x\to x_0^+$).
  \item[$\square$] Je sais exploiter la courbe de $x\mapsto\dfrac{a}{x}$ pour les limites de $x\mapsto\dfrac{ax+b}{cx+d}$ en $-\dfrac{d}{c}$.
  \item[$\square$] Je n'écris jamais $\dfrac{\ell}{0}$ : je parle de limite infinie ou de limites latérales.
  \item[$\square$] Je sais que $f$ est continue en $x_0$ si et seulement si $f$ est définie en $x_0$ et $\lim\limits_{x\to x_0}f(x)=f(x_0)$.
  \item[$\square$] Je justifie la continuité sur $]a;b[$ par une courbe \og sans lever le crayon \fg{} ou par identification à une fonction de référence.
  \item[$\square$] Je rédige mes justifications en trois temps : hypothèse, calcul, conclusion.
  \item[$\square$] Je sais repérer graphiquement une discontinuité (saut, trou, valeur interdite).
  \item[$\square$] Je relie limite et continuité : une fonction continue en $x_0$ admet en $x_0$ une limite égale à son image.
\end{itemize}
\end{aretenir}

\findecours

\end{document}
